% !TeX program = pdflatex
% papers/cristal_diversos.tex — GERADO por tools/cristal_tex.py; NÃO editar à mão.
% Fonte: cristal/cristal.jsonl (broca-so, última versão por conceito).
% Cada secção carrega o registo original em comentário %CRISTAL — a volta é
% tests/cristal_volta.js (resíduo 0 byte a byte contra a fonte).
\documentclass[11pt,a4paper]{article}
\usepackage[utf8]{inputenc}\usepackage[T1]{fontenc}\usepackage[brazil]{babel}
\usepackage{amsmath,amssymb}\usepackage{textcomp}
\usepackage[margin=2.4cm]{geometry}\usepackage{xcolor}
\definecolor{cinza}{gray}{0.35}
\newcommand{\prov}[1]{\par\smallskip\noindent{\small\color{cinza}#1}\par}
\setcounter{secnumdepth}{0}
% ─────────────────────────────────────────────────────────────────────────────
%  papers/gkcapa.tex — a capa do Reino, mínima, para os papers em `article`.
%
%  O estilo.tex completo é do `report` (títulos de capítulo, skak, …). Aqui fica
%  só o que a capa precisa, para o paper continuar a compilar sozinho com
%  pdflatex. O tradutor próprio (tests/tex) carrega o estilo.tex da raiz / o
%  symlink papers/estilo.tex e avalia o mesmo `\gkcapa`.
% ─────────────────────────────────────────────────────────────────────────────
\definecolor{ouro}{HTML}{B8912F}
\definecolor{ouroclaro}{HTML}{F3E9CE}
\definecolor{tinta}{HTML}{1E1B16}
\definecolor{regua}{HTML}{6E6858}
\providecommand{\gknota}{\fontsize{7.62}{11.05}\selectfont}
\providecommand{\gktexto}{\fontsize{10.50}{15.22}\selectfont}
\providecommand{\gksub}{\fontsize{12.33}{17.87}\selectfont}
\providecommand{\gksec}{\fontsize{14.47}{20.98}\selectfont}
\providecommand{\gkcap}{\fontsize{16.99}{24.63}\selectfont}
\providecommand{\gktit}{\fontsize{23.42}{33.95}\selectfont}
\providecommand{\Prj}{\mathbb{P}}
\providecommand{\Trf}{\mathcal{T}}
\providecommand{\gkcapa}[3]{%
\title{\vspace{-0.6cm}
{\color{ouro}\rule{\textwidth}{1.4pt}}\\[6mm]
\textbf{\gktit\color{tinta} Reino Dourado}\\[3mm]{\gksec\itshape\color{regua} Golden Kingdom}\\[8mm]
{\gkcap\scshape\color{ouro} #1}\\[5mm]
{\gksub\color{regua} #2}\\[2mm]
{\gktexto\color{regua}\itshape #3}\\[7mm]
{\color{ouro}\rule{\textwidth}{0.6pt}}\\[9mm]{\gknota\color{regua}\begin{minipage}{\textwidth}\setlength{\parindent}{0pt}%
\textbf{Aviso de Isenção Algorítmica:} O universo, as leis e as entidades contidas nestas páginas ---
incluindo felinos matriciais, aves de rapina algébricas e amuletos de controle de fluxo --- são
construções de pura ficção formal. Esta obra não descreve a realidade física, não serve como
engenharia reversa do cosmos e não constitui doutrina religiosa. Qualquer semelhança com bugs reais do seu
sistema operacional é mera coincidência (ou falta de reza). Prossiga por sua própria conta e
risco.\end{minipage}}}
\author{}\date{}}

\begin{document}
\gkcapa{O Cristal --- os diversos do cristal}
       {o que não coube nos outros grupos}
       {corpus recuperado do broca-so; 88 conceitos; projeção verificável da fonte}
\maketitle

%CRISTAL {"arestas": [{"alvo": "bai", "confianca": 0.117, "epistemico": "fato", "origem": "manual", "rel": "referencia"}, {"alvo": "bai_biblioteca_de_autorizacao_isomorfica", "confianca": 0.95, "epistemico": "fato", "origem": "manual", "rel": "parte_de"}], "confianca": 1.0, "contraexemplos": [], "descricao": "[ --- a proibição sempre vence] A escolha de política está na diferença de conjuntos: para toda linha r, [ r Ep(s,a) (grant g: r cg ) (proibição d: r cd ). ] Uma proibição domina grant sobre a mesma linha, independentemente de “especificidade” ou ordem. Assim, grant = 5 com proibição = A (e 5 A ) resulta: a proibição vence. O sistema é, por definição, --- declarado aqui para remover ambiguidade. (Casa com a semântica do CASL, onde sobrepõe.)", "epistemico": "fato", "exemplos": [], "id": "---_a_proibicao_sempre_vence_a_escolha_de_politica_esta_n", "memoria": "semantica", "meta": {}, "origem": "papers/casl-propagation.tex", "palavras_chave": [], "sinonimos": [], "tipo": "conceito", "titulo": "[ --- a proibição sempre vence] A escolha de política está n…"}
\section{[ --- a proibição sempre vence] A escolha de política está n…}
[ --- a proibição sempre vence] A escolha de política está na diferença de conjuntos: para toda linha r, [ r Ep(s,a) (grant g: r cg ) (proibição d: r cd ). ] Uma proibição domina grant sobre a mesma linha, independentemente de “especificidade” ou ordem. Assim, grant = 5 com proibição = A (e 5 A ) resulta: a proibição vence. O sistema é, por definição, --- declarado aqui para remover ambiguidade. (Casa com a semântica do CASL, onde sobrepõe.)
\prov{relações: referencia bai; parte\_de bai\_biblioteca\_de\_autorizacao\_isomorfica}
\prov{proveniência: papers/casl-propagation.tex \textperiodcentered{} fato \textperiodcentered{} confiança 1.0 \textperiodcentered{} domínio (sem) \textperiodcentered{} história 68 versões no jornal}

%CRISTAL {"arestas": [{"alvo": "corpo_glacial", "confianca": 1.0, "epistemico": "fato", "origem": "paper", "rel": "parte_de"}, {"alvo": "gato", "confianca": 1.0, "epistemico": "fato", "origem": "codigo", "rel": "confirma"}, {"alvo": "transformada_metalica", "confianca": 1.0, "epistemico": "fato", "origem": "wiki", "rel": "parte_de"}], "confianca": 1.0, "contraexemplos": ["Decomposição algébrica e_i M e_j ≠ semiótica triádica C.S. Peirce (homónimo)"], "descricao": "Resolvendo z z=z no produto estelar, os idempotentes não-triviais são eu=12+u, g(u,u)=-14 --- só no split, nulos para a norma, em um par complementar eu+e-u=1, eu e-u=0. Limite: a ortogonalidade bilateral eu ev=ev eu=0 força v=-u; logo o produto estelar de Cayley--Dickson dá exatamente um par de idempotentes em toda a torre (split-///sedeniões). Esse par são as duas células do caso mínimo --- não os canais. (split- M₂(): as duas células são E₁₁,E₂₂.)", "epistemico": "fato", "exemplos": [], "id": "a_decomposicao_de_peirce_celulas_e_canais", "memoria": "semantica", "meta": {"distincao": "algebra_vs_filosofia", "verificacao": "slug idêntico — enriquecer, não duplicar"}, "origem": "wiki", "palavras_chave": [], "sinonimos": [], "tipo": "conceito", "titulo": "A decomposição de Peirce: células e canais"}
\section{A decomposição de Peirce: células e canais}
Resolvendo z z=z no produto estelar, os idempotentes não-triviais são eu=12+u, g(u,u)=-14 --- só no split, nulos para a norma, em um par complementar eu+e-u=1, eu e-u=0. Limite: a ortogonalidade bilateral eu ev=ev eu=0 força v=-u; logo o produto estelar de Cayley--Dickson dá exatamente um par de idempotentes em toda a torre (split-///sedeniões). Esse par são as duas células do caso mínimo --- não os canais. (split- M\ensuremath{_{2}}(): as duas células são E\ensuremath{_{11}},E\ensuremath{_{22}}.)
\prov{contraexemplos: Decomposição algébrica e\_i M e\_j \ensuremath{\ne} semiótica triádica C.S. Peirce (homónimo)}
\prov{relações: parte\_de corpo\_glacial; confirma gato; parte\_de transformada\_metalica}
\prov{proveniência: wiki \textperiodcentered{} fato \textperiodcentered{} confiança 1.0 \textperiodcentered{} domínio (sem) \textperiodcentered{} história 49 versões no jornal}

%CRISTAL {"arestas": [{"alvo": "mecanica_quantica", "confianca": 1.0, "epistemico": "fato", "origem": "paper", "rel": "parte_de"}, {"alvo": "mecanica", "confianca": 1.0, "epistemico": "fato", "origem": "paper", "rel": "parte_de"}], "confianca": 1.0, "contraexemplos": [], "descricao": "A pressão algébrica =1-x² muda de sinal em |x|=1: para |x|<1, >0 --- regime dissipativo, a termodinâmica, a norma contrai; em |x|=1, =0 --- o ponto hermitiano, onde a norma é exatamente conservada. É ali (s=1, a=0) que a evolução é reversível e a mecânica quântica vive. O eixo dissipativohermitiano é a parábola: a termodinâmica habita o interior (|x|<1), a quântica a fronteira (|x|=1).", "epistemico": "fato", "exemplos": ["Holonomia Berry / transporte paralelo na esfera de Bloch"], "id": "a_dinamica_a_precessao_de_bloch", "memoria": "semantica", "meta": {"confirmado_web": "fase_geometrica_berry", "verificacao": "slug idêntico — enriquecer, não duplicar"}, "origem": "wiki", "palavras_chave": [], "sinonimos": [], "tipo": "conceito", "titulo": "A dinâmica: a precessão de Bloch"}
\section{A dinâmica: a precessão de Bloch}
A pressão algébrica =1-x\ensuremath{^{2}} muda de sinal em |x|=1: para |x|<1, >0 --- regime dissipativo, a termodinâmica, a norma contrai; em |x|=1, =0 --- o ponto hermitiano, onde a norma é exatamente conservada. É ali (s=1, a=0) que a evolução é reversível e a mecânica quântica vive. O eixo dissipativohermitiano é a parábola: a termodinâmica habita o interior (|x|<1), a quântica a fronteira (|x|=1).
\prov{exemplos: Holonomia Berry / transporte paralelo na esfera de Bloch}
\prov{relações: parte\_de mecanica\_quantica; parte\_de mecanica}
\prov{proveniência: wiki \textperiodcentered{} fato \textperiodcentered{} confiança 1.0 \textperiodcentered{} domínio (sem) \textperiodcentered{} história 21 versões no jornal}

%CRISTAL {"arestas": [{"alvo": "vetores", "confianca": 0.094, "epistemico": "fato", "origem": "manual", "rel": "referencia"}, {"alvo": "mecanica_quantica", "confianca": 1.0, "epistemico": "fato", "origem": "paper", "rel": "parte_de"}, {"alvo": "mecanica", "confianca": 1.0, "epistemico": "fato", "origem": "paper", "rel": "parte_de"}], "confianca": 1.0, "contraexemplos": [], "descricao": "Um qubit é a álgebra: C² S³. N funções de onda são o produto tensorial [ HN= C² C²N= C²N, ] de dimensão 2N. Aqui há uma escolha decisiva, e ela é honesta. A multiplicação de ouro sobre vetores dá a soma direta R⁴ R⁴ --- local-realista por construção: estados independentes, |S CHSH|=2. O emaranhamento genuíno exige o produto tensorial sobre operadores: só ali Bell atinge 22 e os estados coletivos [ |GHZN=12(|00+|11) C²N ] existem. Verificado: Bell =2", "epistemico": "fato", "exemplos": ["Berry phase γ = −½ solid angle (ciclo adiabático S²)", "Proc. R. Soc. A 1984"], "id": "a_funcao_de_onda_a_esfera_de_bloch", "memoria": "semantica", "meta": {"confirmado_web": "fase_geometrica_berry", "doi": "10.1098/rspa.1984.0023", "verificacao": "slug idêntico — enriquecer, não duplicar"}, "origem": "wiki", "palavras_chave": [], "sinonimos": [], "tipo": "conceito", "titulo": "A função de onda: a esfera de Bloch"}
\section{A função de onda: a esfera de Bloch}
Um qubit é a álgebra: C\ensuremath{^{2}} S\ensuremath{^{3}}. N funções de onda são o produto tensorial [ HN= C\ensuremath{^{2}} C\ensuremath{^{2}}N= C\ensuremath{^{2}}N, ] de dimensão 2N. Aqui há uma escolha decisiva, e ela é honesta. A multiplicação de ouro sobre vetores dá a soma direta R\ensuremath{^{4}} R\ensuremath{^{4}} --- local-realista por construção: estados independentes, |S CHSH|=2. O emaranhamento genuíno exige o produto tensorial sobre operadores: só ali Bell atinge 22 e os estados coletivos [ |GHZN=12(|00+|11) C\ensuremath{^{2}}N ] existem. Verificado: Bell =2
\prov{exemplos: Berry phase \ensuremath{\gamma} = \ensuremath{-}\ensuremath{\tfrac12} solid angle (ciclo adiabático S\ensuremath{^{2}}); Proc. R. Soc. A 1984}
\prov{relações: referencia vetores; parte\_de mecanica\_quantica; parte\_de mecanica}
\prov{proveniência: wiki \textperiodcentered{} fato \textperiodcentered{} confiança 1.0 \textperiodcentered{} domínio (sem) \textperiodcentered{} história 30 versões no jornal}

%CRISTAL {"arestas": [{"alvo": "delta", "confianca": 0.95, "epistemico": "fato", "origem": "manual", "rel": "define"}], "confianca": 1.0, "contraexemplos": [], "descricao": "[A intensidade não entra aqui] E_p depende de : a profundidade governa a (Seção~ ), não a avaliação. Quem detém uma capacidade a exerce plenamente, qualquer que seja o seu restante. Isso é o que torna o barato (Seção~ ).", "epistemico": "fato", "exemplos": [], "id": "a_intensidade_nao_entra_aqui_e_p_depende_de_a_profundida", "memoria": "semantica", "meta": {}, "origem": "papers/casl-propagation.tex", "palavras_chave": [], "sinonimos": [], "tipo": "conceito", "titulo": "[A intensidade não entra aqui] E_p depende de : a profundida…"}
\section{[A intensidade não entra aqui] E\_p depende de : a profundida…}
[A intensidade não entra aqui] E\_p depende de : a profundidade governa a (Seção\~{} ), não a avaliação. Quem detém uma capacidade a exerce plenamente, qualquer que seja o seu restante. Isso é o que torna o barato (Seção\~{} ).
\prov{relações: define delta}
\prov{proveniência: papers/casl-propagation.tex \textperiodcentered{} fato \textperiodcentered{} confiança 1.0 \textperiodcentered{} domínio (sem) \textperiodcentered{} história 61 versões no jornal}

%CRISTAL {"arestas": [{"alvo": "delta", "confianca": 0.95, "epistemico": "fato", "origem": "manual", "rel": "especializa"}, {"alvo": "borda_contemplacao", "confianca": 0.5, "epistemico": "fato", "origem": "wiki", "rel": "relaciona"}, {"alvo": "melhor_candidato", "confianca": 0.5, "epistemico": "fato", "origem": "wiki", "rel": "relaciona"}, {"alvo": "fase_2_tecido_de_conhecimento", "confianca": 0.5, "epistemico": "fato", "origem": "wiki", "rel": "relaciona"}, {"alvo": "inspecao_local_32", "confianca": 0.5, "epistemico": "fato", "origem": "wiki", "rel": "relaciona"}], "confianca": 1.0, "contraexemplos": [], "descricao": "A intensidade : profundidade de propagação", "epistemico": "fato", "exemplos": [], "id": "a_intensidade_profundidade_de_propagacao", "memoria": "semantica", "meta": {}, "origem": "papers/casl-propagation.tex", "palavras_chave": [], "sinonimos": [], "tipo": "conceito", "titulo": "A intensidade : profundidade de propagação"}
\section{A intensidade : profundidade de propagação}
A intensidade : profundidade de propagação
\prov{relações: especializa delta; relaciona borda\_contemplacao; relaciona melhor\_candidato; relaciona fase\_2\_tecido\_de\_conhecimento; relaciona inspecao\_local\_32}
\prov{proveniência: papers/casl-propagation.tex \textperiodcentered{} fato \textperiodcentered{} confiança 1.0 \textperiodcentered{} domínio (sem) \textperiodcentered{} história 65 versões no jornal}

%CRISTAL {"arestas": [{"alvo": "geometria", "confianca": 0.039, "epistemico": "fato", "origem": "manual", "rel": "referencia"}, {"alvo": "mecanica_quantica", "confianca": 1.0, "epistemico": "fato", "origem": "paper", "rel": "parte_de"}, {"alvo": "mecanica", "confianca": 1.0, "epistemico": "fato", "origem": "paper", "rel": "parte_de"}], "confianca": 1.0, "contraexemplos": [], "descricao": "O estado é um versor de S³ (Fundamentos da Geometria Riemanniana); a evolução unitária gira o versor sem mudar a norma. A fibração de Hopf S³ S² dá a esfera de Bloch: o qubit é a base S²¹ (estados puros), e a fase global é a fibra U(1), não observável. A logo da Engenharia de Ouro --- o yin-yang --- é esse estado: duas amplitudes, uma fase, norma um.", "epistemico": "fato", "exemplos": ["Hopf S³→S²", "fibra U(1)", "Wikipedia Bloch sphere"], "id": "a_medida_o_calor_de_volta", "memoria": "semantica", "meta": {"confirmado_web": "fibracao_hopf_bloch", "verificacao": "slug idêntico — enriquecer, não duplicar"}, "origem": "wiki", "palavras_chave": [], "sinonimos": [], "tipo": "conceito", "titulo": "A medida: o calor de volta"}
\section{A medida: o calor de volta}
O estado é um versor de S\ensuremath{^{3}} (Fundamentos da Geometria Riemanniana); a evolução unitária gira o versor sem mudar a norma. A fibração de Hopf S\ensuremath{^{3}} S\ensuremath{^{2}} dá a esfera de Bloch: o qubit é a base S\ensuremath{^{21}} (estados puros), e a fase global é a fibra U(1), não observável. A logo da Engenharia de Ouro --- o yin-yang --- é esse estado: duas amplitudes, uma fase, norma um.
\prov{exemplos: Hopf S\ensuremath{^{3}}\ensuremath{\to}S\ensuremath{^{2}}; fibra U(1); Wikipedia Bloch sphere}
\prov{relações: referencia geometria; parte\_de mecanica\_quantica; parte\_de mecanica}
\prov{proveniência: wiki \textperiodcentered{} fato \textperiodcentered{} confiança 1.0 \textperiodcentered{} domínio (sem) \textperiodcentered{} história 23 versões no jornal}

%CRISTAL {"arestas": [{"alvo": "reticulado", "confianca": -0.023, "epistemico": "fato", "origem": "manual", "rel": "referencia"}, {"alvo": "casl-propagation", "confianca": 0.95, "epistemico": "fato", "origem": "manual", "rel": "parte_de"}, {"alvo": "reticulados", "confianca": 0.95, "epistemico": "fato", "origem": "manual", "rel": "define"}], "confianca": 1.0, "contraexemplos": [], "descricao": "Apresentamos a: biblioteca de sobre reticulados booleanos. Primitivas: =(c, o), (propagação discreta), () (estrutura vs. alcance), operadores T id (atenuação), P id (propagação) e = ( _ ) sobre candidatos vigentes _ =: >0. Teoremas: não-escalada, propagação limitada, fail-closed --- independentes de domínio. Toda instância é determinada pelo quinteto (T, ); CASL, e são interpretações. instanciam; analogias físicas ficam no Apêndice.", "epistemico": "fato", "exemplos": [], "id": "apresentamos_a_biblioteca_de_sobre_reticulados_booleanos", "memoria": "semantica", "meta": {}, "origem": "papers/casl-propagation.tex", "palavras_chave": [], "sinonimos": [], "tipo": "conceito", "titulo": "Apresentamos a : biblioteca de sobre reticulados booleanos. …"}
\section{Apresentamos a : biblioteca de sobre reticulados booleanos. …}
Apresentamos a: biblioteca de sobre reticulados booleanos. Primitivas: =(c, o), (propagação discreta), () (estrutura vs. alcance), operadores T id (atenuação), P id (propagação) e = ( \_ ) sobre candidatos vigentes \_ =: >0. Teoremas: não-escalada, propagação limitada, fail-closed --- independentes de domínio. Toda instância é determinada pelo quinteto (T, ); CASL, e são interpretações. instanciam; analogias físicas ficam no Apêndice.
\prov{relações: referencia reticulado; parte\_de casl-propagation; define reticulados}
\prov{proveniência: papers/casl-propagation.tex \textperiodcentered{} fato \textperiodcentered{} confiança 1.0 \textperiodcentered{} domínio (sem) \textperiodcentered{} história 123 versões no jornal}

%CRISTAL {"arestas": [{"alvo": "bai", "confianca": 0.016, "epistemico": "fato", "origem": "manual", "rel": "referencia"}, {"alvo": "monotono", "confianca": 0.95, "epistemico": "fato", "origem": "manual", "rel": "especializa"}, {"alvo": "bai_biblioteca_de_autorizacao_isomorfica", "confianca": 0.95, "epistemico": "fato", "origem": "manual", "rel": "parte_de"}], "confianca": 1.0, "contraexemplos": [], "descricao": "[Associativo estrutural] Em instâncias com endereçamento associativo (overlap, extent SQL), opera sobre o associativo; captura o que o endereçamento não conserva. Os dois eixos são independentes por construção.", "epistemico": "fato", "exemplos": [], "id": "associativo_estrutural_em_instancias_com_enderecamento_ass", "memoria": "semantica", "meta": {}, "origem": "papers/casl-propagation.tex", "palavras_chave": [], "sinonimos": [], "tipo": "conceito", "titulo": "[Associativo estrutural] Em instâncias com endereçamento ass…"}
\section{[Associativo estrutural] Em instâncias com endereçamento ass…}
[Associativo estrutural] Em instâncias com endereçamento associativo (overlap, extent SQL), opera sobre o associativo; captura o que o endereçamento não conserva. Os dois eixos são independentes por construção.
\prov{relações: referencia bai; especializa monotono; parte\_de bai\_biblioteca\_de\_autorizacao\_isomorfica}
\prov{proveniência: papers/casl-propagation.tex \textperiodcentered{} fato \textperiodcentered{} confiança 1.0 \textperiodcentered{} domínio (sem) \textperiodcentered{} história 122 versões no jornal}

%CRISTAL {"arestas": [{"alvo": "kappa", "confianca": 0.95, "epistemico": "fato", "origem": "manual", "rel": "explica"}, {"alvo": "delta", "confianca": 0.95, "epistemico": "fato", "origem": "manual", "rel": "explica"}], "confianca": 1.0, "contraexemplos": [], "descricao": "[Atenuação e propagação] Sobre extents: T:2X 2X é ( T id, x y T(x) T(y) ). Sobre detentores: P:2^ P 2^ P é ( P id ). Re-delegar por saltos é P^; = corresponde ao menor ponto fixo P= ₙ Pⁿ(H₀).", "epistemico": "fato", "exemplos": [], "id": "atenuacao_e_propagacao_sobre_extents_t2x_2x_e_t_id", "memoria": "semantica", "meta": {}, "origem": "papers/casl-propagation.tex", "palavras_chave": [], "sinonimos": [], "tipo": "conceito", "titulo": "[Atenuação e propagação] Sobre extents: T:2X 2X é ( T id,…"}
\section{[Atenuação e propagação] Sobre extents: T:2X 2X é ( T id,…}
[Atenuação e propagação] Sobre extents: T:2X 2X é ( T id, x y T(x) T(y) ). Sobre detentores: P:2\^{} P 2\^{} P é ( P id ). Re-delegar por saltos é P\^{}; = corresponde ao menor ponto fixo P= \ensuremath{_{n}} P\ensuremath{^{n}}(H\ensuremath{_{0}}).
\prov{relações: explica kappa; explica delta}
\prov{proveniência: papers/casl-propagation.tex \textperiodcentered{} fato \textperiodcentered{} confiança 1.0 \textperiodcentered{} domínio (sem) \textperiodcentered{} história 121 versões no jornal}

%CRISTAL {"arestas": [{"alvo": "casl-propagation", "confianca": 1.0, "epistemico": "fato", "origem": "codigo", "rel": "especializa"}, {"alvo": "tiffany", "confianca": 1.0, "epistemico": "fato", "origem": "paper", "rel": "referencia"}, {"alvo": "broca_os", "confianca": 1.0, "epistemico": "fato", "origem": "paper", "rel": "referencia"}, {"alvo": "colapso", "confianca": 1.0, "epistemico": "fato", "origem": "paper", "rel": "explica"}, {"alvo": "bai", "confianca": 1.0, "epistemico": "fato", "origem": "paper", "rel": "eh_um"}, {"alvo": "reticulados", "confianca": 1.0, "epistemico": "fato", "origem": "codigo", "rel": "especializa"}, {"alvo": "monotono", "confianca": 1.0, "epistemico": "fato", "origem": "codigo", "rel": "depende"}, {"alvo": "kappa", "confianca": 1.0, "epistemico": "fato", "origem": "codigo", "rel": "define"}, {"alvo": "delta", "confianca": 1.0, "epistemico": "fato", "origem": "codigo", "rel": "define"}], "confianca": 1.0, "contraexemplos": ["Peirce (decomposição algébrica) ≠ BAI — BAI usa reticulados", "BAI ≠ RBAC/ACL clássico — delegação capacitária monótona", "Geometria riemanniana ≠ reticulado BAI — não isomorfos", "Ψ BAI ≠ colapso postulado da função de onda (MQ)"], "descricao": "Biblioteca de Autorização Isomórfica: delegação capacitária monótona sobre reticulados booleanos. Tese-mãe: instância determinada pelo quinteto (κ,δ,Φ,T,Ψ).", "epistemico": "fato", "exemplos": ["papers/casl-propagation.pdf — tese-mãe e universalidade", "theorem_tese-mae — toda instância pelo quinteto", "fail-closed: união vazia se nenhum grant detido", "conversa/colapso.py — Ψ = Collapse(𝒞_δ) sobre candidatos vigentes"], "id": "bai_biblioteca_de_autorizacao_isomorfica", "memoria": "semantica", "meta": {"confirmado": "bai_quinteto_hub", "fechado_fase7": true}, "origem": "codigo", "palavras_chave": ["κ", "δ", "Φ", "Ψ", "autorização", "isomorfismo"], "sinonimos": ["BAI", "biblioteca de autorização isomórfica", "autorização isomórfica"], "tipo": "conceito", "titulo": "bai_biblioteca_de_autorizacao_isomorfica"}
\section{bai\_biblioteca\_de\_autorizacao\_isomorfica}
Biblioteca de Autorização Isomórfica: delegação capacitária monótona sobre reticulados booleanos. Tese-mãe: instância determinada pelo quinteto (\ensuremath{\kappa},\ensuremath{\delta},\ensuremath{\Phi},T,\ensuremath{\Psi}).
\prov{exemplos: papers/casl-propagation.pdf — tese-mãe e universalidade; theorem\_tese-mae — toda instância pelo quinteto; fail-closed: união vazia se nenhum grant detido; conversa/colapso.py — \ensuremath{\Psi} = Collapse(\ensuremath{\mathcal{C}}\_\ensuremath{\delta}) sobre candidatos vigentes}
\prov{contraexemplos: Peirce (decomposição algébrica) \ensuremath{\ne} BAI — BAI usa reticulados; BAI \ensuremath{\ne} RBAC/ACL clássico — delegação capacitária monótona; Geometria riemanniana \ensuremath{\ne} reticulado BAI — não isomorfos; \ensuremath{\Psi} BAI \ensuremath{\ne} colapso postulado da função de onda (MQ)}
\prov{sinónimos: BAI; biblioteca de autorização isomórfica; autorização isomórfica}
\prov{relações: especializa casl-propagation; referencia tiffany; referencia broca\_os; explica colapso; eh\_um bai; especializa reticulados; depende monotono; define kappa; define delta}
\prov{proveniência: codigo \textperiodcentered{} fato \textperiodcentered{} confiança 1.0 \textperiodcentered{} domínio (sem) \textperiodcentered{} história 484 versões no jornal}

%CRISTAL {"arestas": [{"alvo": "broca_so", "confianca": 1.0, "epistemico": "fato", "origem": "manual", "rel": "parte_de"}, {"alvo": "processos", "confianca": 1.0, "epistemico": "fato", "origem": "codigo", "rel": "usa"}, {"alvo": "syscall", "confianca": 1.0, "epistemico": "fato", "origem": "codigo", "rel": "referencia"}, {"alvo": "kernel", "confianca": 1.0, "epistemico": "fato", "origem": "codigo", "rel": "parte_de"}, {"alvo": "gato", "confianca": 1.0, "epistemico": "fato", "origem": "codigo", "rel": "usa"}, {"alvo": "colapso", "confianca": 1.0, "epistemico": "fato", "origem": "codigo", "rel": "explica"}, {"alvo": "hopfield", "confianca": 1.0, "epistemico": "fato", "origem": "codigo", "rel": "explica"}], "confianca": 1.0, "contraexemplos": ["SO clássico (Linux, UI, logs) = camada de contemplação, não de trabalho", "escalonador por fatia de tempo — a célula soma sem ordem, sem buffer", "hypervisor/VM genérico — não há isomorfismo álgebra↔SO", "broca_borda (visualização na borda) — projeção, não o núcleo de trabalho"], "descricao": "**Decisão:** SO clássico = contemplação. Broca-SO = trabalho (memória + álgebra).", "epistemico": "fato", "exemplos": ["processo = idempotente de Peirce (célula e_i em broca_so.h)", "IPC = canal de Peirce e_i M e_j", "medição = operador gato A sobre bytes", "rede = memória associativa de Hopfield"], "id": "broca-so_camada_de_trabalho_isomorfica", "memoria": "semantica", "meta": {"confirmado": "broca_os_hub", "fechado_fase7": true, "nivel": 6, "serve": "Camada de trabalho do SO fractal — memória + álgebra; separa trabalho de contemplação clássica"}, "origem": "codigo", "palavras_chave": [], "sinonimos": [], "tipo": "conceito", "titulo": "Broca-SO — camada de trabalho isomórfica"}
\section{Broca-SO — camada de trabalho isomórfica}
**Decisão:** SO clássico = contemplação. Broca-SO = trabalho (memória + álgebra).
\prov{exemplos: processo = idempotente de Peirce (célula e\_i em broca\_so.h); IPC = canal de Peirce e\_i M e\_j; medição = operador gato A sobre bytes; rede = memória associativa de Hopfield}
\prov{contraexemplos: SO clássico (Linux, UI, logs) = camada de contemplação, não de trabalho; escalonador por fatia de tempo — a célula soma sem ordem, sem buffer; hypervisor/VM genérico — não há isomorfismo álgebra\ensuremath{\leftrightarrow}SO; broca\_borda (visualização na borda) — projeção, não o núcleo de trabalho}
\prov{relações: parte\_de broca\_so; usa processos; referencia syscall; parte\_de kernel; usa gato; explica colapso; explica hopfield}
\prov{proveniência: codigo \textperiodcentered{} fato \textperiodcentered{} confiança 1.0 \textperiodcentered{} domínio (sem) \textperiodcentered{} história 182 versões no jornal}

%CRISTAL {"arestas": [{"alvo": "broca_os", "confianca": 1.0, "epistemico": "fato", "origem": "manual", "rel": "parte_de"}, {"alvo": "destruir_a_parabola", "confianca": 0.5, "epistemico": "fato", "origem": "wiki", "rel": "relaciona"}, {"alvo": "word", "confianca": 0.5, "epistemico": "fato", "origem": "wiki", "rel": "relaciona"}, {"alvo": "vontade_de_poder", "confianca": 0.5, "epistemico": "fato", "origem": "wiki", "rel": "relaciona"}, {"alvo": "virtualizacao", "confianca": 0.5, "epistemico": "fato", "origem": "wiki", "rel": "relaciona"}], "confianca": 1.0, "contraexemplos": [], "descricao": "Tudo deriva do gato **A = [[1,1],[1,0]]**: φ, U=e⁻iA, o espectro (Cooley–Tukey), a função de onda.", "epistemico": "fato", "exemplos": [], "id": "broca-so_o_gato_no_metal", "memoria": "semantica", "meta": {}, "origem": "README.md", "palavras_chave": [], "sinonimos": [], "tipo": "conceito", "titulo": "broca-so — o gato no metal"}
\section{broca-so — o gato no metal}
Tudo deriva do gato **A = [[1,1],[1,0]]**: \ensuremath{\varphi}, U=e\ensuremath{^{-}}iA, o espectro (Cooley–Tukey), a função de onda.
\prov{relações: parte\_de broca\_os; relaciona destruir\_a\_parabola; relaciona word; relaciona vontade\_de\_poder; relaciona virtualizacao}
\prov{proveniência: README.md \textperiodcentered{} fato \textperiodcentered{} confiança 1.0 \textperiodcentered{} domínio (sem) \textperiodcentered{} história 12 versões no jornal}

%CRISTAL {"arestas": [{"alvo": "computacao", "confianca": 1.0, "epistemico": "fato", "origem": "curriculo", "rel": "parte_de"}, {"alvo": "goldchain", "confianca": 1.0, "epistemico": "fato", "origem": "curriculo", "rel": "usa"}, {"alvo": "kernel", "confianca": 1.0, "epistemico": "fato", "origem": "paper", "rel": "parte_de"}, {"alvo": "gato", "confianca": 1.0, "epistemico": "fato", "origem": "codigo", "rel": "depende"}, {"alvo": "processos", "confianca": 1.0, "epistemico": "fato", "origem": "codigo", "rel": "depende"}, {"alvo": "ipc", "confianca": 1.0, "epistemico": "fato", "origem": "codigo", "rel": "usa"}, {"alvo": "hopfield", "confianca": 1.0, "epistemico": "fato", "origem": "codigo", "rel": "usa"}, {"alvo": "colapso", "confianca": 1.0, "epistemico": "fato", "origem": "codigo", "rel": "usa"}, {"alvo": "tiffany", "confianca": 1.0, "epistemico": "fato", "origem": "codigo", "rel": "explica"}, {"alvo": "broca-so_camada_de_trabalho_isomorfica", "confianca": 1.0, "epistemico": "fato", "origem": "codigo", "rel": "especializa"}, {"alvo": "sistema_operacional", "confianca": 1.0, "epistemico": "fato", "origem": "wiki", "rel": "especializa"}, {"alvo": "filosofia_unix", "confianca": 1.0, "epistemico": "fato", "origem": "wiki", "rel": "usa"}], "confianca": 1.0, "contraexemplos": ["Broca OS ≠ Linux kernel — contemplação corre em borda clássica", "Broca OS ≠ Von Neumann puro — unifica Peirce (estado+canal), não nega Turing", "broca_os (SO fractal) ≠ broca_so (header C) — mesmo design, níveis distintos", "Tiffany ≠ substituto do SO — cabeça conversacional parte de Broca OS", "Broca OS ≠ Linux kernel", "broca_os ≠ broca_so (header C)"], "descricao": "SO fractal: processo=idempotente, IPC=Peirce, rede=Hopfield, gato A mede bytes.", "epistemico": "fato", "exemplos": ["papers/broca_os.tex", "neuronio/neuronio.c", "neuronio/radio.py", "papers/broca_os.tex — primitivas SO fractal", "neuronio/neuronio.c, neuronio/radio.py — implementação", "GATO_STREAM → TORO → REGISTRO_N → COLAPSO_6 → BANDA → NONCE", "ula/BROCA_SO.md — manual da camada C"], "id": "broca_os", "memoria": "semantica", "meta": {"confirmado": "broca_os_hub", "fechado_fase7": true}, "origem": "codigo", "palavras_chave": ["Peirce", "gato", "Tiffany", "Hopfield"], "sinonimos": ["Broca OS", "broca-so", "BrocaOS"], "tipo": "conceito", "titulo": "broca_os"}
\section{broca\_os}
SO fractal: processo=idempotente, IPC=Peirce, rede=Hopfield, gato A mede bytes.
\prov{exemplos: papers/broca\_os.tex; neuronio/neuronio.c; neuronio/radio.py; papers/broca\_os.tex — primitivas SO fractal; neuronio/neuronio.c, neuronio/radio.py — implementação; GATO\_STREAM \ensuremath{\to} TORO \ensuremath{\to} REGISTRO\_N \ensuremath{\to} COLAPSO\_6 \ensuremath{\to} BANDA \ensuremath{\to} NONCE; ula/BROCA\_SO.md — manual da camada C}
\prov{contraexemplos: Broca OS \ensuremath{\ne} Linux kernel — contemplação corre em borda clássica; Broca OS \ensuremath{\ne} Von Neumann puro — unifica Peirce (estado+canal), não nega Turing; broca\_os (SO fractal) \ensuremath{\ne} broca\_so (header C) — mesmo design, níveis distintos; Tiffany \ensuremath{\ne} substituto do SO — cabeça conversacional parte de Broca OS; Broca OS \ensuremath{\ne} Linux kernel; broca\_os \ensuremath{\ne} broca\_so (header C)}
\prov{sinónimos: Broca OS; broca-so; BrocaOS}
\prov{relações: parte\_de computacao; usa goldchain; parte\_de kernel; depende gato; depende processos; usa ipc; usa hopfield; usa colapso; explica tiffany; especializa broca-so\_camada\_de\_trabalho\_isomorfica; especializa sistema\_operacional; usa filosofia\_unix}
\prov{proveniência: codigo \textperiodcentered{} fato \textperiodcentered{} confiança 1.0 \textperiodcentered{} domínio (sem) \textperiodcentered{} história 619 versões no jornal}

%CRISTAL {"arestas": [{"alvo": "broca_os_hub_broca_os_hub_do_sistema_fractal", "confianca": 1.0, "epistemico": "fato", "origem": "broca_os_hub.md", "rel": "parte_de"}, {"alvo": "tiffany", "confianca": 0.95, "epistemico": "fato", "origem": "wiki", "rel": "explica"}, {"alvo": "colapso", "confianca": 0.95, "epistemico": "fato", "origem": "wiki", "rel": "usa"}, {"alvo": "hopfield", "confianca": 0.95, "epistemico": "fato", "origem": "wiki", "rel": "usa"}, {"alvo": "ipc", "confianca": 0.95, "epistemico": "fato", "origem": "wiki", "rel": "usa"}, {"alvo": "a_decomposicao_de_peirce_celulas_e_canais", "confianca": 0.95, "epistemico": "fato", "origem": "wiki", "rel": "usa"}, {"alvo": "broca-so_camada_de_trabalho_isomorfica", "confianca": 0.95, "epistemico": "fato", "origem": "wiki", "rel": "especializa"}, {"alvo": "goldchain", "confianca": 0.95, "epistemico": "fato", "origem": "wiki", "rel": "usa"}, {"alvo": "computacao", "confianca": 0.95, "epistemico": "fato", "origem": "wiki", "rel": "parte_de"}], "confianca": 1.0, "contraexemplos": [], "descricao": "| Broca OS | Nó | relação | |----------|-----|---------| | núcleo | broca_os | confirma | | álgebra gato | gato | depende | | IPC Peirce | ipc | usa | | decomposição | a_decomposicao_de_peirce_celulas_e_canais | usa | | processos SO | processos | depende | | computação | computacao | parte_de | | cabeça | tiffany | explica | | colapso | colapso | usa | | Hopfield | hopfield | usa | | camada C | broca-so_camada_de_trabalho_isomorfica | especializa |", "epistemico": "fato", "exemplos": [], "id": "broca_os_hub_ponte_broca-so", "memoria": "semantica", "meta": {"indexar": false, "serve": "Mapeamento Broca OS ↔ nós do grafo de conhecimento"}, "origem": "codigo", "palavras_chave": [], "sinonimos": [], "tipo": "conceito", "titulo": "Ponte Broca-SO"}
\section{Ponte Broca-SO}
| Broca OS | Nó | relação | |----------|-----|---------| | núcleo | broca\_os | confirma | | álgebra gato | gato | depende | | IPC Peirce | ipc | usa | | decomposição | a\_decomposicao\_de\_peirce\_celulas\_e\_canais | usa | | processos SO | processos | depende | | computação | computacao | parte\_de | | cabeça | tiffany | explica | | colapso | colapso | usa | | Hopfield | hopfield | usa | | camada C | broca-so\_camada\_de\_trabalho\_isomorfica | especializa |
\prov{relações: parte\_de broca\_os\_hub\_broca\_os\_hub\_do\_sistema\_fractal; explica tiffany; usa colapso; usa hopfield; usa ipc; usa a\_decomposicao\_de\_peirce\_celulas\_e\_canais; especializa broca-so\_camada\_de\_trabalho\_isomorfica; usa goldchain; parte\_de computacao}
\prov{proveniência: codigo \textperiodcentered{} fato \textperiodcentered{} confiança 1.0 \textperiodcentered{} domínio (sem) \textperiodcentered{} história 19 versões no jornal}

%CRISTAL {"arestas": [{"alvo": "possui_gerenciamento_de_processos", "confianca": 0.55, "epistemico": "inferencia", "origem": "semilla", "rel": "relaciona"}, {"alvo": "broca_os", "confianca": 0.95, "epistemico": "fato", "origem": "paper", "rel": "parte_de"}], "confianca": 0.55, "contraexemplos": [], "descricao": "Fato (inferencia): Broca OS relaciona possui gerenciamento de processos. Origem: semilla.", "epistemico": "inferencia", "exemplos": [], "id": "broca_os_relaciona_possui_gerenciamento_de_processos", "memoria": "semantica", "meta": {"predicado": "possui gerenciamento de processos", "sujeito": "Broca OS"}, "origem": "semilla", "palavras_chave": ["Broca OS", "possui gerenciamento de processos"], "sinonimos": [], "tipo": "fato", "titulo": "Broca OS relaciona possui gerenciamento de processos"}
\section{Broca OS relaciona possui gerenciamento de processos}
Fato (inferencia): Broca OS relaciona possui gerenciamento de processos. Origem: semilla.
\prov{relações: relaciona possui\_gerenciamento\_de\_processos; parte\_de broca\_os}
\prov{proveniência: semilla \textperiodcentered{} inferencia \textperiodcentered{} confiança 0.55 \textperiodcentered{} domínio (sem) \textperiodcentered{} história 3 versões no jornal}

%CRISTAL {"arestas": [{"alvo": "grava_fato_resposta_normal_em_seguida", "confianca": 1.0, "epistemico": "fato", "origem": "CONHECIMENTO.md", "rel": "parte_de"}, {"alvo": "semantica", "confianca": 1.0, "epistemico": "fato", "origem": "paper", "rel": "usa"}, {"alvo": "colapso", "confianca": 0.95, "epistemico": "fato", "origem": "manual", "rel": "especializa"}, {"alvo": "word", "confianca": 0.5, "epistemico": "fato", "origem": "wiki", "rel": "relaciona"}, {"alvo": "virtualizacao", "confianca": 0.5, "epistemico": "fato", "origem": "wiki", "rel": "relaciona"}], "confianca": 1.0, "contraexemplos": [], "descricao": "Ordem: **dialogos (8)** → **conhecimento (6)** → **corpus (4)**.", "epistemico": "fato", "exemplos": [], "id": "camada_no_colapso", "memoria": "semantica", "meta": {}, "origem": "CONHECIMENTO.md", "palavras_chave": [], "sinonimos": [], "tipo": "conceito", "titulo": "Camada no colapso"}
\section{Camada no colapso}
Ordem: **dialogos (8)** \ensuremath{\to} **conhecimento (6)** \ensuremath{\to} **corpus (4)**.
\prov{relações: parte\_de grava\_fato\_resposta\_normal\_em\_seguida; usa semantica; especializa colapso; relaciona word; relaciona virtualizacao}
\prov{proveniência: CONHECIMENTO.md \textperiodcentered{} fato \textperiodcentered{} confiança 1.0 \textperiodcentered{} domínio (sem) \textperiodcentered{} história 116 versões no jornal}

%CRISTAL {"arestas": [{"alvo": "bai", "confianca": 0.008, "epistemico": "fato", "origem": "manual", "rel": "referencia"}, {"alvo": "bai_biblioteca_de_autorizacao_isomorfica", "confianca": 0.95, "epistemico": "fato", "origem": "manual", "rel": "define"}], "confianca": 1.0, "contraexemplos": [], "descricao": "[Capacidade CASL] Specialização da Definição~ com X= (linhas), =(a,s) (ação e subject), ( = 1 é proibição, ) e como orçamento (Definição~ ). Formalmente [ =(a, s, c, ), ] autoriza (ou nega, se = 1 ) a sobre s nas linhas de c, e admite re-delegação conforme.", "epistemico": "fato", "exemplos": [], "id": "capacidade_casl_specializacao_da_definicao_com_x_linhas", "memoria": "semantica", "meta": {}, "origem": "papers/casl-propagation.tex", "palavras_chave": [], "sinonimos": [], "tipo": "conceito", "titulo": "[Capacidade CASL] Specialização da Definição~ com X= (linhas…"}
\section{[Capacidade CASL] Specialização da Definição\~{} com X= (linhas…}
[Capacidade CASL] Specialização da Definição\~{} com X= (linhas), =(a,s) (ação e subject), ( = 1 é proibição, ) e como orçamento (Definição\~{} ). Formalmente [ =(a, s, c, ), ] autoriza (ou nega, se = 1 ) a sobre s nas linhas de c, e admite re-delegação conforme.
\prov{relações: referencia bai; define bai\_biblioteca\_de\_autorizacao\_isomorfica}
\prov{proveniência: papers/casl-propagation.tex \textperiodcentered{} fato \textperiodcentered{} confiança 1.0 \textperiodcentered{} domínio (sem) \textperiodcentered{} história 64 versões no jornal}

%CRISTAL {"arestas": [{"alvo": "kappa", "confianca": 0.95, "epistemico": "fato", "origem": "manual", "rel": "define"}, {"alvo": "bai_kappa", "confianca": 1.0, "epistemico": "fato", "origem": "wiki", "rel": "confirma"}], "confianca": 1.0, "contraexemplos": [], "descricao": "[Capacidade] Uma é uma tupla [ =(c, o), ] onde é o, c a condição, o orçamento (Seção~ ), o peso (grant/deny quando aplicável) e o a origem. Nenhum domínio entra na definição.", "epistemico": "fato", "exemplos": [], "id": "capacidade_uma_e_uma_tupla_c_o", "memoria": "semantica", "meta": {}, "origem": "papers/casl-propagation.tex", "palavras_chave": [], "sinonimos": [], "tipo": "conceito", "titulo": "[Capacidade] Uma é uma tupla [ =(c, o), …"}
\section{[Capacidade] Uma é uma tupla [ =(c, o), …}
[Capacidade] Uma é uma tupla [ =(c, o), ] onde é o, c a condição, o orçamento (Seção\~{} ), o peso (grant/deny quando aplicável) e o a origem. Nenhum domínio entra na definição.
\prov{relações: define kappa; confirma bai\_kappa}
\prov{proveniência: papers/casl-propagation.tex \textperiodcentered{} fato \textperiodcentered{} confiança 1.0 \textperiodcentered{} domínio (sem) \textperiodcentered{} história 63 versões no jornal}

%CRISTAL {"arestas": [{"alvo": "reticulado", "confianca": 0.016, "epistemico": "fato", "origem": "manual", "rel": "referencia"}, {"alvo": "bai", "confianca": 0.016, "epistemico": "fato", "origem": "manual", "rel": "referencia"}, {"alvo": "reticulados", "confianca": 0.95, "epistemico": "fato", "origem": "manual", "rel": "especializa"}, {"alvo": "bai_biblioteca_de_autorizacao_isomorfica", "confianca": 0.95, "epistemico": "fato", "origem": "manual", "rel": "parte_de"}, {"alvo": "bai_t_atenuacao", "confianca": 1.0, "epistemico": "fato", "origem": "wiki", "rel": "confirma"}], "confianca": 1.0, "contraexemplos": [], "descricao": "[Capacidades detidas] Hold (p) é o menor conjunto que contém as capacidades intrínsecas de p (se p é ou dono) e toda tal que existe uma aresta válida u p .", "epistemico": "fato", "exemplos": [], "id": "capacidades_detidas_hold_p_e_o_menor_conjunto_que_contem", "memoria": "semantica", "meta": {}, "origem": "papers/casl-propagation.tex", "palavras_chave": [], "sinonimos": [], "tipo": "conceito", "titulo": "[Capacidades detidas] Hold (p) é o menor conjunto que contém…"}
\section{[Capacidades detidas] Hold (p) é o menor conjunto que contém…}
[Capacidades detidas] Hold (p) é o menor conjunto que contém as capacidades intrínsecas de p (se p é ou dono) e toda tal que existe uma aresta válida u p .
\prov{relações: referencia reticulado; referencia bai; especializa reticulados; parte\_de bai\_biblioteca\_de\_autorizacao\_isomorfica; confirma bai\_t\_atenuacao}
\prov{proveniência: papers/casl-propagation.tex \textperiodcentered{} fato \textperiodcentered{} confiança 1.0 \textperiodcentered{} domínio (sem) \textperiodcentered{} história 126 versões no jornal}

%CRISTAL {"arestas": [{"alvo": "quatro_ciencias", "confianca": 1.0, "epistemico": "fato", "origem": "manual", "rel": "generaliza"}, {"alvo": "casl-propagation", "confianca": 1.0, "epistemico": "fato", "origem": "manual", "rel": "relaciona"}, {"alvo": "fisica", "confianca": 1.0, "epistemico": "fato", "origem": "manual", "rel": "parte_de"}, {"alvo": "matematica", "confianca": 1.0, "epistemico": "fato", "origem": "manual", "rel": "parte_de"}, {"alvo": "biologia", "confianca": 1.0, "epistemico": "fato", "origem": "manual", "rel": "parte_de"}], "confianca": 1.0, "contraexemplos": [], "descricao": "Rede fractal de domínios — matemática, física, vida e computação entrelaçam-se por analogia e isomorfismo.", "epistemico": "fato", "exemplos": ["quatro quadrantes do gato", "BAI ↔ reticulados ↔ Broca OS"], "id": "ciencias", "memoria": "semantica", "meta": {"verificacao": "slug idêntico — enriquecer, não duplicar"}, "origem": "manual", "palavras_chave": [], "sinonimos": [], "tipo": "conceito", "titulo": "ciencias"}
\section{ciencias}
Rede fractal de domínios — matemática, física, vida e computação entrelaçam-se por analogia e isomorfismo.
\prov{exemplos: quatro quadrantes do gato; BAI \ensuremath{\leftrightarrow} reticulados \ensuremath{\leftrightarrow} Broca OS}
\prov{relações: generaliza quatro\_ciencias; relaciona casl-propagation; parte\_de fisica; parte\_de matematica; parte\_de biologia}
\prov{proveniência: manual \textperiodcentered{} fato \textperiodcentered{} confiança 1.0 \textperiodcentered{} domínio (sem) \textperiodcentered{} história 338 versões no jornal}

%CRISTAL {"arestas": [{"alvo": "kappa", "confianca": 1.0, "epistemico": "fato", "origem": "codigo", "rel": "depende"}, {"alvo": "gato", "confianca": 1.0, "epistemico": "fato", "origem": "codigo", "rel": "usa"}, {"alvo": "semantica", "confianca": 1.0, "epistemico": "fato", "origem": "codigo", "rel": "usa"}, {"alvo": "olho_de_koch", "confianca": 1.0, "epistemico": "fato", "origem": "paper", "rel": "usa"}, {"alvo": "painel_estelar", "confianca": 1.0, "epistemico": "fato", "origem": "manual", "rel": "referencia"}, {"alvo": "codigo_conversa_colapso", "confianca": 1.0, "epistemico": "fato", "origem": "codigo", "rel": "define"}, {"alvo": "metais", "confianca": 1.0, "epistemico": "fato", "origem": "paper", "rel": "usa"}, {"alvo": "delta", "confianca": 1.0, "epistemico": "fato", "origem": "codigo", "rel": "depende"}, {"alvo": "tiffany", "confianca": 1.0, "epistemico": "fato", "origem": "codigo", "rel": "explica"}, {"alvo": "camada_no_colapso", "confianca": 1.0, "epistemico": "fato", "origem": "codigo", "rel": "especializa"}, {"alvo": "bai_biblioteca_de_autorizacao_isomorfica", "confianca": 1.0, "epistemico": "fato", "origem": "codigo", "rel": "especializa"}], "confianca": 1.0, "contraexemplos": ["Ψ Broca ≠ colapso postulado ψ MQ", "Dougherty vórtice — analogia, não implementação", "Ψ Broca ≠ colapso postulado da função de onda (MQ)", "Analogia Dougherty vórtice/toro ≠ implementação Ψ", "machine tecido peso 0 sem --full"], "descricao": "Operador Ψ = Collapse(𝒞_δ): selecciona face vencedora no tecido Koch (overlap prosa + camadas curadas). Motor conversacional da Tiffany — não colapso MQ.", "epistemico": "fato", "exemplos": ["colapso.py seleciona face vencedora", "Ψ com orçamento δ", "pow_gato", "painel_estelar", "hipótese Dougherty (vórtice/toro)", "colapso.py — camadas dialogos(8) > conhecimento(6) > corpus(4)", "BAI quinteto κ,δ,Φ,T,Ψ", "lab.py porque <fala>", "conversa/colapso.py — dialogos(8) > conhecimento(6) > corpus(4)", "lab.py porque <fala> — explica face vencedora", "contexto.bits + ultimo.face — estado pós-Ψ", "BAI quinteto: κ, δ, Φ, T, Ψ"], "id": "colapso", "memoria": "semantica", "meta": {"serve": "Seleção Koch; núcleo da Tiffany conversa", "verificacao": "slug idêntico — enriquecer, não duplicar"}, "origem": "manual", "palavras_chave": [], "sinonimos": [], "tipo": "conceito", "titulo": "colapso"}
\section{colapso}
Operador \ensuremath{\Psi} = Collapse(\ensuremath{\mathcal{C}}\_\ensuremath{\delta}): selecciona face vencedora no tecido Koch (overlap prosa + camadas curadas). Motor conversacional da Tiffany — não colapso MQ.
\prov{exemplos: colapso.py seleciona face vencedora; \ensuremath{\Psi} com orçamento \ensuremath{\delta}; pow\_gato; painel\_estelar; hipótese Dougherty (vórtice/toro); colapso.py — camadas dialogos(8) > conhecimento(6) > corpus(4); BAI quinteto \ensuremath{\kappa},\ensuremath{\delta},\ensuremath{\Phi},T,\ensuremath{\Psi}; lab.py porque <fala>; conversa/colapso.py — dialogos(8) > conhecimento(6) > corpus(4); lab.py porque <fala> — explica face vencedora; contexto.bits + ultimo.face — estado pós-\ensuremath{\Psi}; BAI quinteto: \ensuremath{\kappa}, \ensuremath{\delta}, \ensuremath{\Phi}, T, \ensuremath{\Psi}}
\prov{contraexemplos: \ensuremath{\Psi} Broca \ensuremath{\ne} colapso postulado \ensuremath{\psi} MQ; Dougherty vórtice — analogia, não implementação; \ensuremath{\Psi} Broca \ensuremath{\ne} colapso postulado da função de onda (MQ); Analogia Dougherty vórtice/toro \ensuremath{\ne} implementação \ensuremath{\Psi}; machine tecido peso 0 sem --full}
\prov{relações: depende kappa; usa gato; usa semantica; usa olho\_de\_koch; referencia painel\_estelar; define codigo\_conversa\_colapso; usa metais; depende delta; explica tiffany; especializa camada\_no\_colapso; especializa bai\_biblioteca\_de\_autorizacao\_isomorfica}
\prov{proveniência: manual \textperiodcentered{} fato \textperiodcentered{} confiança 1.0 \textperiodcentered{} domínio (sem) \textperiodcentered{} história 393 versões no jornal}

%CRISTAL {"arestas": [{"alvo": "bai", "confianca": -0.016, "epistemico": "fato", "origem": "manual", "rel": "referencia"}, {"alvo": "bai_biblioteca_de_autorizacao_isomorfica", "confianca": 0.95, "epistemico": "fato", "origem": "manual", "rel": "parte_de"}], "confianca": 1.0, "contraexemplos": [], "descricao": "[Compartilhamento controlado] Conceder uma capacidade cross- com =0 garante que o destinatário a exerce mas ninguém abaixo dele a herda --- exatamente ``dar acesso sem propagar''. Com =k , a difusão é confinada à vizinhança de re-delegação de raio k .", "epistemico": "fato", "exemplos": [], "id": "compartilhamento_controlado_conceder_uma_capacidade_cross-", "memoria": "semantica", "meta": {}, "origem": "papers/casl-propagation.tex", "palavras_chave": [], "sinonimos": [], "tipo": "conceito", "titulo": "[Compartilhamento controlado] Conceder uma capacidade cross-…"}
\section{[Compartilhamento controlado] Conceder uma capacidade cross-…}
[Compartilhamento controlado] Conceder uma capacidade cross- com =0 garante que o destinatário a exerce mas ninguém abaixo dele a herda --- exatamente ``dar acesso sem propagar''. Com =k , a difusão é confinada à vizinhança de re-delegação de raio k .
\prov{relações: referencia bai; parte\_de bai\_biblioteca\_de\_autorizacao\_isomorfica}
\prov{proveniência: papers/casl-propagation.tex \textperiodcentered{} fato \textperiodcentered{} confiança 1.0 \textperiodcentered{} domínio (sem) \textperiodcentered{} história 63 versões no jornal}

%CRISTAL {"arestas": [{"alvo": "delta", "confianca": 0.95, "epistemico": "fato", "origem": "manual", "rel": "explica"}], "confianca": 1.0, "contraexemplos": [], "descricao": "[Concessão válida / atenuação] A aresta u w é se existe ₀ Hold (u) tal que [ c_ c_ ₀, (s,a) ₀ (s,a), _ = _ ₀, _ ₀ 1, _ _ ₀ 1, _ _ ₀, ] com a convenção 1=. Isto é, u só concede o que detém, a condição, a profundidade e a validade (expiração, Seção~ ). Em particular _ ₀ =0 proíbe qualquer concessão.", "epistemico": "fato", "exemplos": [], "id": "concessao_valida_atenuacao_a_aresta_u_w_e_se_existe_0_h", "memoria": "semantica", "meta": {}, "origem": "papers/casl-propagation.tex", "palavras_chave": [], "sinonimos": [], "tipo": "conceito", "titulo": "[Concessão válida / atenuação] A aresta u w é se existe _0 H…"}
\section{[Concessão válida / atenuação] A aresta u w é se existe \_0 H…}
[Concessão válida / atenuação] A aresta u w é se existe \ensuremath{_{0}} Hold (u) tal que [ c\_ c\_ \ensuremath{_{0}}, (s,a) \ensuremath{_{0}} (s,a), \_ = \_ \ensuremath{_{0}}, \_ \ensuremath{_{0}} 1, \_ \_ \ensuremath{_{0}} 1, \_ \_ \ensuremath{_{0}}, ] com a convenção 1=. Isto é, u só concede o que detém, a condição, a profundidade e a validade (expiração, Seção\~{} ). Em particular \_ \ensuremath{_{0}} =0 proíbe qualquer concessão.
\prov{relações: explica delta}
\prov{proveniência: papers/casl-propagation.tex \textperiodcentered{} fato \textperiodcentered{} confiança 1.0 \textperiodcentered{} domínio (sem) \textperiodcentered{} história 61 versões no jornal}

%CRISTAL {"arestas": [{"alvo": "ponte_quantica_conjunto_dougherty_e_orbitais_de_hidrogenio", "confianca": 1.0, "epistemico": "fato", "origem": "ponte_quantica_dougherty.md", "rel": "parte_de"}, {"alvo": "geometria", "confianca": 0.52, "epistemico": "hipotese", "origem": "ingest", "rel": "analogia"}, {"alvo": "colapso", "confianca": 0.52, "epistemico": "hipotese", "origem": "ingest", "rel": "analogia"}, {"alvo": "a_medida_o_calor_de_volta", "confianca": 1.0, "epistemico": "fato", "origem": "ingest", "rel": "analogia"}], "confianca": 0.95, "contraexemplos": ["Toro φ-Dougherty ≠ fibrado Hopf S³→S² (estruturas topológicas distintas)", "Fase geométrica Berry (MQ standard) ≠ correção log_φ(r) Dougherty"], "descricao": "Geometria dinâmica projectiva 4-D: toros de dilatação interligados, floresta de vórtices, planos radicais “sphere-kissing”, auto-similaridade em escala φ.", "epistemico": "fato", "exemplos": ["Filamentos aurora ~100 m (NASA) — analogia vórtice, não prova φ", "Paralelo Ginzburg 3D-SST — toro/hélice independente, analogia não prova φ"], "id": "conjunto_dougherty", "memoria": "semantica", "meta": {"verificacao": "slug idêntico — enriquecer, não duplicar"}, "origem": "wiki", "palavras_chave": [], "sinonimos": [], "tipo": "conceito", "titulo": "Conjunto Dougherty"}
\section{Conjunto Dougherty}
Geometria dinâmica projectiva 4-D: toros de dilatação interligados, floresta de vórtices, planos radicais “sphere-kissing”, auto-similaridade em escala \ensuremath{\varphi}.
\prov{exemplos: Filamentos aurora \~{}100 m (NASA) — analogia vórtice, não prova \ensuremath{\varphi}; Paralelo Ginzburg 3D-SST — toro/hélice independente, analogia não prova \ensuremath{\varphi}}
\prov{contraexemplos: Toro \ensuremath{\varphi}-Dougherty \ensuremath{\ne} fibrado Hopf S\ensuremath{^{3}}\ensuremath{\to}S\ensuremath{^{2}} (estruturas topológicas distintas); Fase geométrica Berry (MQ standard) \ensuremath{\ne} correção log\_\ensuremath{\varphi}(r) Dougherty}
\prov{relações: parte\_de ponte\_quantica\_conjunto\_dougherty\_e\_orbitais\_de\_hidrogenio; analogia geometria; analogia colapso; analogia a\_medida\_o\_calor\_de\_volta}
\prov{proveniência: wiki \textperiodcentered{} fato \textperiodcentered{} confiança 0.95 \textperiodcentered{} domínio (sem) \textperiodcentered{} história 51 versões no jornal}

%CRISTAL {"arestas": [{"alvo": "colapso", "confianca": 0.95, "epistemico": "fato", "origem": "manual", "rel": "explica"}, {"alvo": "kappa", "confianca": 0.95, "epistemico": "fato", "origem": "manual", "rel": "usa"}], "confianca": 1.0, "contraexemplos": [], "descricao": "[Conjunto vigente e colapso ] Seja = ᵢ capacidades detidas, filtradas por vigência e por cᵢ q para consulta q. Defina [ _ =: _ >0 ou avaliação local, (q)= ( _,q) = _ _ score (q) ] (quando o máximo existe; senão (q)= --- ). é de restante na leitura: só entra na formação de via delegação prévia.", "epistemico": "fato", "exemplos": [], "id": "conjunto_vigente_e_colapso_seja_i_capacidades_detida", "memoria": "semantica", "meta": {}, "origem": "papers/casl-propagation.tex", "palavras_chave": [], "sinonimos": [], "tipo": "conceito", "titulo": "[Conjunto vigente e colapso ] Seja = ᵢ capacidades detida…"}
\section{[Conjunto vigente e colapso ] Seja = \ensuremath{_{i}} capacidades detida…}
[Conjunto vigente e colapso ] Seja = \ensuremath{_{i}} capacidades detidas, filtradas por vigência e por c\ensuremath{_{i}} q para consulta q. Defina [ \_ =: \_ >0 ou avaliação local, (q)= ( \_,q) = \_ \_ score (q) ] (quando o máximo existe; senão (q)= --- ). é de restante na leitura: só entra na formação de via delegação prévia.
\prov{relações: explica colapso; usa kappa}
\prov{proveniência: papers/casl-propagation.tex \textperiodcentered{} fato \textperiodcentered{} confiança 1.0 \textperiodcentered{} domínio (sem) \textperiodcentered{} história 121 versões no jornal}

%CRISTAL {"arestas": [{"alvo": "microscopia_de_fotoionizacao_orbitais_de_hidrogenio_stodolna_2013", "confianca": 1.0, "epistemico": "fato", "origem": "microscopia_fotoionizacao_hidrogenio.md", "rel": "parte_de"}, {"alvo": "fibracao_de_hopf_e_esfera_de_bloch", "confianca": 1.0, "epistemico": "fato", "origem": "fibracao_hopf_bloch.md", "rel": "parte_de"}, {"alvo": "correntes_de_birkeland_filamentos_de_plasma", "confianca": 1.0, "epistemico": "fato", "origem": "birkeland_correntes_plasma.md", "rel": "parte_de"}, {"alvo": "semiotica_triadica_de_charles_sanders_peirce", "confianca": 1.0, "epistemico": "fato", "origem": "semiotica_peirce_triade.md", "rel": "parte_de"}, {"alvo": "ginzburg_toro_helice_e_3d-sst", "confianca": 1.0, "epistemico": "fato", "origem": "ginzburg_toro_helice.md", "rel": "parte_de"}, {"alvo": "fase_geometrica_de_berry_pancharatnamberry", "confianca": 1.0, "epistemico": "fato", "origem": "fase_geometrica_berry.md", "rel": "parte_de"}, {"alvo": "colapso_koch_operador_ψ_da_tiffany", "confianca": 1.0, "epistemico": "fato", "origem": "colapso_koch_tiffany.md", "rel": "parte_de"}, {"alvo": "tiffany_cabeca_conversacional_broca-so", "confianca": 1.0, "epistemico": "fato", "origem": "tiffany_cabeca_broca.md", "rel": "parte_de"}, {"alvo": "geometria_hub_broca-so", "confianca": 1.0, "epistemico": "fato", "origem": "geometria_hub_broca.md", "rel": "parte_de"}], "confianca": 1.0, "contraexemplos": [], "descricao": "Broca OS **≠** Linux kernel — contemplação corre em borda clássica.", "epistemico": "fato", "exemplos": [], "id": "contraexemplos", "memoria": "semantica", "meta": {"indexar": false, "verificacao": "slug idêntico — enriquecer, não duplicar"}, "origem": "manual", "palavras_chave": [], "sinonimos": [], "tipo": "conceito", "titulo": "Contraexemplos"}
\section{Contraexemplos}
Broca OS **\ensuremath{\ne}** Linux kernel — contemplação corre em borda clássica.
\prov{relações: parte\_de microscopia\_de\_fotoionizacao\_orbitais\_de\_hidrogenio\_stodolna\_2013; parte\_de fibracao\_de\_hopf\_e\_esfera\_de\_bloch; parte\_de correntes\_de\_birkeland\_filamentos\_de\_plasma; parte\_de semiotica\_triadica\_de\_charles\_sanders\_peirce; parte\_de ginzburg\_toro\_helice\_e\_3d-sst; parte\_de fase\_geometrica\_de\_berry\_pancharatnamberry; parte\_de colapso\_koch\_operador\_\ensuremath{\psi}\_da\_tiffany; parte\_de tiffany\_cabeca\_conversacional\_broca-so; parte\_de geometria\_hub\_broca-so}
\prov{proveniência: manual \textperiodcentered{} fato \textperiodcentered{} confiança 1.0 \textperiodcentered{} domínio (sem) \textperiodcentered{} história 51 versões no jornal}

%CRISTAL {"arestas": [{"alvo": "word", "confianca": 0.5, "epistemico": "fato", "origem": "wiki", "rel": "relaciona"}, {"alvo": "virtualizacao", "confianca": 0.5, "epistemico": "fato", "origem": "wiki", "rel": "relaciona"}], "confianca": 1.0, "contraexemplos": [], "descricao": "Orçamento discreto de propagação; limita profundidade do colapso Koch.", "epistemico": "fato", "exemplos": ["δ orçamento limita profundidade Ψ / colapso Koch", "Orçamento δ limita profundidade do colapso Koch / Ψ"], "id": "delta", "memoria": "semantica", "meta": {}, "origem": "codigo", "palavras_chave": [], "sinonimos": [], "tipo": "conceito", "titulo": "delta"}
\section{delta}
Orçamento discreto de propagação; limita profundidade do colapso Koch.
\prov{exemplos: \ensuremath{\delta} orçamento limita profundidade \ensuremath{\Psi} / colapso Koch; Orçamento \ensuremath{\delta} limita profundidade do colapso Koch / \ensuremath{\Psi}}
\prov{relações: relaciona word; relaciona virtualizacao}
\prov{proveniência: codigo \textperiodcentered{} fato \textperiodcentered{} confiança 1.0 \textperiodcentered{} domínio (sem) \textperiodcentered{} história 80 versões no jornal}

%CRISTAL {"arestas": [{"alvo": "bai", "confianca": -0.078, "epistemico": "fato", "origem": "manual", "rel": "referencia"}, {"alvo": "termodinamica", "confianca": 0.95, "epistemico": "fato", "origem": "manual", "rel": "analogia"}, {"alvo": "bai_biblioteca_de_autorizacao_isomorfica", "confianca": 0.95, "epistemico": "fato", "origem": "manual", "rel": "referencia"}], "confianca": 1.0, "contraexemplos": [], "descricao": "[ é uma energia bem-fundada ciclos são inofensivos] Como toda aresta válida decrementa a profundidade ( _ _ _0 -1 , Definição~ ), é um sobre as cadeias de re-delegação: qualquer caminho de concessões com finito tem comprimento e termina. --- um ciclo u v u apenas faz cair a cada volta até 0 , quando a propagação para. A terminação (e a não-escalada) vem da atenuação, não da estrutura do grafo. Só = ignora a energia; reserve-o às autoridades intrínsecas.", "epistemico": "fato", "exemplos": [], "id": "e_uma_energia_bem-fundada_ciclos_sao_inofensivos_como_tod", "memoria": "semantica", "meta": {}, "origem": "papers/casl-propagation.tex", "palavras_chave": [], "sinonimos": [], "tipo": "conceito", "titulo": "[ é uma energia bem-fundada ciclos são inofensivos] Como tod…"}
\section{[ é uma energia bem-fundada ciclos são inofensivos] Como tod…}
[ é uma energia bem-fundada ciclos são inofensivos] Como toda aresta válida decrementa a profundidade ( \_ \_ \_0 -1 , Definição\~{} ), é um sobre as cadeias de re-delegação: qualquer caminho de concessões com finito tem comprimento e termina. --- um ciclo u v u apenas faz cair a cada volta até 0 , quando a propagação para. A terminação (e a não-escalada) vem da atenuação, não da estrutura do grafo. Só = ignora a energia; reserve-o às autoridades intrínsecas.
\prov{relações: referencia bai; analogia termodinamica; referencia bai\_biblioteca\_de\_autorizacao\_isomorfica}
\prov{proveniência: papers/casl-propagation.tex \textperiodcentered{} fato \textperiodcentered{} confiança 1.0 \textperiodcentered{} domínio (sem) \textperiodcentered{} história 122 versões no jornal}

%CRISTAL {"arestas": [{"alvo": "fase_2_tecido_de_conhecimento", "confianca": 1.0, "epistemico": "fato", "origem": "CONHECIMENTO.md", "rel": "parte_de"}, {"alvo": "virtualizacao", "confianca": 0.5, "epistemico": "fato", "origem": "wiki", "rel": "relaciona"}, {"alvo": "word", "confianca": 0.5, "epistemico": "fato", "origem": "wiki", "rel": "relaciona"}], "confianca": 1.0, "contraexemplos": [], "descricao": "```bash export TIFFANY_APRENDER=1 python3 code/tiffany.py \"Broca OS é um sistema operacional.\"", "epistemico": "fato", "exemplos": [], "id": "ensino_automatico_dialogo", "memoria": "semantica", "meta": {}, "origem": "CONHECIMENTO.md", "palavras_chave": [], "sinonimos": [], "tipo": "conceito", "titulo": "Ensino automático (diálogo)"}
\section{Ensino automático (diálogo)}
```bash export TIFFANY\_APRENDER=1 python3 code/tiffany.py "Broca OS é um sistema operacional."
\prov{relações: parte\_de fase\_2\_tecido\_de\_conhecimento; relaciona virtualizacao; relaciona word}
\prov{proveniência: CONHECIMENTO.md \textperiodcentered{} fato \textperiodcentered{} confiança 1.0 \textperiodcentered{} domínio (sem) \textperiodcentered{} história 4 versões no jornal}

%CRISTAL {"arestas": [{"alvo": "bai", "confianca": -0.008, "epistemico": "fato", "origem": "manual", "rel": "referencia"}, {"alvo": "bai_biblioteca_de_autorizacao_isomorfica", "confianca": 0.95, "epistemico": "fato", "origem": "manual", "rel": "parte_de"}], "confianca": 1.0, "contraexemplos": [], "descricao": "[Esboço] Itens (i)--(iv) são exatamente as hipóteses das Definições~ , e~ . Mapear estados para , extents para c , delegações para arestas com decrementado, e avaliação para . Não-escalada vem de T id ; propagação limitada de como variante bem-fundada; fail-closed de ser o neutro de .", "epistemico": "fato", "exemplos": [], "id": "esboco_itens_i--iv_sao_exatamente_as_hipoteses_das_def", "memoria": "semantica", "meta": {}, "origem": "papers/casl-propagation.tex", "palavras_chave": [], "sinonimos": [], "tipo": "conceito", "titulo": "[Esboço] Itens (i)--(iv) são exatamente as hipóteses das Def…"}
\section{[Esboço] Itens (i)--(iv) são exatamente as hipóteses das Def…}
[Esboço] Itens (i)--(iv) são exatamente as hipóteses das Definições\~{} , e\~{} . Mapear estados para , extents para c , delegações para arestas com decrementado, e avaliação para . Não-escalada vem de T id ; propagação limitada de como variante bem-fundada; fail-closed de ser o neutro de .
\prov{relações: referencia bai; parte\_de bai\_biblioteca\_de\_autorizacao\_isomorfica}
\prov{proveniência: papers/casl-propagation.tex \textperiodcentered{} fato \textperiodcentered{} confiança 1.0 \textperiodcentered{} domínio (sem) \textperiodcentered{} história 63 versões no jornal}

%CRISTAL {"arestas": [{"alvo": "casl-propagation", "confianca": 0.95, "epistemico": "fato", "origem": "manual", "rel": "parte_de"}], "confianca": 1.0, "contraexemplos": [], "descricao": "Este documento apresenta uma --- não uma aplicação. Autoridade flui, atenua ( T id ), propaga com orçamento , organiza-se ( ) e colapsa ( ) ou falha fechada ( ).", "epistemico": "fato", "exemplos": [], "id": "este_documento_apresenta_uma_---_nao_uma_aplicacao_autorida", "memoria": "semantica", "meta": {}, "origem": "papers/casl-propagation.tex", "palavras_chave": [], "sinonimos": [], "tipo": "conceito", "titulo": "Este documento apresenta uma --- não uma aplicação. Autorida…"}
\section{Este documento apresenta uma --- não uma aplicação. Autorida…}
Este documento apresenta uma --- não uma aplicação. Autoridade flui, atenua ( T id ), propaga com orçamento , organiza-se ( ) e colapsa ( ) ou falha fechada ( ).
\prov{relações: parte\_de casl-propagation}
\prov{proveniência: papers/casl-propagation.tex \textperiodcentered{} fato \textperiodcentered{} confiança 1.0 \textperiodcentered{} domínio (sem) \textperiodcentered{} história 60 versões no jornal}

%CRISTAL {"arestas": [{"alvo": "bai_biblioteca_de_autorizacao_isomorfica", "confianca": 0.031, "epistemico": "fato", "origem": "manual", "rel": "referencia"}, {"alvo": "bai", "confianca": -0.023, "epistemico": "fato", "origem": "manual", "rel": "referencia"}, {"alvo": "colapso", "confianca": 0.85, "epistemico": "fato", "origem": "manual", "rel": "referencia"}], "confianca": 1.0, "contraexemplos": [], "descricao": "| Parte | Conteúdo | |-------|----------| | **I** | Biblioteca matemática — κ, δ, (Φ,δ), T, P, Ψ | | **II** | CASL — RLS, MongoQuery, grafo de delegação | | **III** | Broca OS — tecidos Koch, `colapso.py` | | **IV** | Tiffany — conversa, `mede_flocos.py` | | **Apêndice** | Isomorfismos, interpretação física, plano |", "epistemico": "fato", "exemplos": [], "id": "estrutura_do_documento", "memoria": "semantica", "meta": {}, "origem": "CASL.md", "palavras_chave": [], "sinonimos": [], "tipo": "conceito", "titulo": "Estrutura do documento"}
\section{Estrutura do documento}
| Parte | Conteúdo | |-------|----------| | **I** | Biblioteca matemática — \ensuremath{\kappa}, \ensuremath{\delta}, (\ensuremath{\Phi},\ensuremath{\delta}), T, P, \ensuremath{\Psi} | | **II** | CASL — RLS, MongoQuery, grafo de delegação | | **III** | Broca OS — tecidos Koch, `colapso.py` | | **IV** | Tiffany — conversa, `mede\_flocos.py` | | **Apêndice** | Isomorfismos, interpretação física, plano |
\prov{relações: referencia bai\_biblioteca\_de\_autorizacao\_isomorfica; referencia bai; referencia colapso}
\prov{proveniência: CASL.md \textperiodcentered{} fato \textperiodcentered{} confiança 1.0 \textperiodcentered{} domínio (sem) \textperiodcentered{} história 11 versões no jornal}

%CRISTAL {"arestas": [{"alvo": "bai", "confianca": 0.85, "epistemico": "fato", "origem": "manual", "rel": "referencia"}, {"alvo": "bai_biblioteca_de_autorizacao_isomorfica", "confianca": 0.95, "epistemico": "fato", "origem": "manual", "rel": "confirma"}], "confianca": 1.0, "contraexemplos": [], "descricao": "[Fail-closed = união vazia] Se nenhum grant detido casa (s,a) , então E_p(s,a)= : acesso negado. É o da álgebra, uniforme --- sem piso de tenant especial.", "epistemico": "fato", "exemplos": [], "id": "fail-closed_uniao_vazia_se_nenhum_grant_detido_casa_sa", "memoria": "semantica", "meta": {}, "origem": "papers/casl-propagation.tex", "palavras_chave": [], "sinonimos": [], "tipo": "conceito", "titulo": "[Fail-closed = união vazia] Se nenhum grant detido casa (s,a…"}
\section{[Fail-closed = união vazia] Se nenhum grant detido casa (s,a…}
[Fail-closed = união vazia] Se nenhum grant detido casa (s,a) , então E\_p(s,a)= : acesso negado. É o da álgebra, uniforme --- sem piso de tenant especial.
\prov{relações: referencia bai; confirma bai\_biblioteca\_de\_autorizacao\_isomorfica}
\prov{proveniência: papers/casl-propagation.tex \textperiodcentered{} fato \textperiodcentered{} confiança 1.0 \textperiodcentered{} domínio (sem) \textperiodcentered{} história 63 versões no jornal}

%CRISTAL {"arestas": [{"alvo": "ponte_quantica_conjunto_dougherty_e_orbitais_de_hidrogenio", "confianca": 1.0, "epistemico": "fato", "origem": "ponte_quantica_dougherty.md", "rel": "parte_de"}, {"alvo": "contradicao", "confianca": 0.52, "epistemico": "hipotese", "origem": "ingest", "rel": "explica"}, {"alvo": "ginzburg_toro_helice_e_3d-sst", "confianca": 1.0, "epistemico": "fato", "origem": "ginzburg_toro_helice.md", "rel": "parte_de"}], "confianca": 0.95, "contraexemplos": [], "descricao": "Critérios de refutação da ponte Dougherty. Teste 1: topologia nodal PRL 2013. Teste 2: filamentos plasma (Birkeland). Teste 3: espectros MQ.", "epistemico": "fato", "exemplos": ["Teste 2: subestruturas φ em filamentos Birkeland — ainda aberto", "Ginzburg 3D-SST: previsões superluminais não confirmadas mainstream"], "id": "falsificabilidade", "memoria": "semantica", "meta": {"verificacao": "slug idêntico — enriquecer, não duplicar"}, "origem": "wiki", "palavras_chave": [], "sinonimos": [], "tipo": "conceito", "titulo": "Falsificabilidade"}
\section{Falsificabilidade}
Critérios de refutação da ponte Dougherty. Teste 1: topologia nodal PRL 2013. Teste 2: filamentos plasma (Birkeland). Teste 3: espectros MQ.
\prov{exemplos: Teste 2: subestruturas \ensuremath{\varphi} em filamentos Birkeland — ainda aberto; Ginzburg 3D-SST: previsões superluminais não confirmadas mainstream}
\prov{relações: parte\_de ponte\_quantica\_conjunto\_dougherty\_e\_orbitais\_de\_hidrogenio; explica contradicao; parte\_de ginzburg\_toro\_helice\_e\_3d-sst}
\prov{proveniência: wiki \textperiodcentered{} fato \textperiodcentered{} confiança 0.95 \textperiodcentered{} domínio (sem) \textperiodcentered{} história 24 versões no jornal}

%CRISTAL {"arestas": [{"alvo": "utilitarismo", "confianca": 0.5, "epistemico": "fato", "origem": "wiki", "rel": "relaciona"}, {"alvo": "rust", "confianca": 0.5, "epistemico": "fato", "origem": "wiki", "rel": "relaciona"}, {"alvo": "o_que_e_no_projecto", "confianca": 0.5, "epistemico": "fato", "origem": "wiki", "rel": "relaciona"}, {"alvo": "grupos", "confianca": 0.5, "epistemico": "fato", "origem": "wiki", "rel": "relaciona"}, {"alvo": "recursao", "confianca": 0.5, "epistemico": "fato", "origem": "wiki", "rel": "relaciona"}, {"alvo": "funil_cumulativo", "confianca": 0.5, "epistemico": "fato", "origem": "wiki", "rel": "relaciona"}, {"alvo": "melhor_candidato", "confianca": 0.5, "epistemico": "fato", "origem": "wiki", "rel": "relaciona"}, {"alvo": "funcao", "confianca": 0.5, "epistemico": "fato", "origem": "wiki", "rel": "relaciona"}, {"alvo": "pow_bit_taxa", "confianca": 0.5, "epistemico": "fato", "origem": "wiki", "rel": "relaciona"}, {"alvo": "mecanica_quantica_hub", "confianca": 0.5, "epistemico": "fato", "origem": "wiki", "rel": "relaciona"}, {"alvo": "geometria_hub_broca", "confianca": 0.5, "epistemico": "fato", "origem": "wiki", "rel": "relaciona"}, {"alvo": "pell", "confianca": 0.5, "epistemico": "fato", "origem": "wiki", "rel": "relaciona"}, {"alvo": "teste_ordem", "confianca": 0.5, "epistemico": "fato", "origem": "wiki", "rel": "relaciona"}, {"alvo": "motivos", "confianca": 0.5, "epistemico": "fato", "origem": "wiki", "rel": "relaciona"}, {"alvo": "regua_associativa_e_o_risco_de_vazamento", "confianca": 0.5, "epistemico": "fato", "origem": "wiki", "rel": "relaciona"}, {"alvo": "topologia", "confianca": 0.5, "epistemico": "fato", "origem": "wiki", "rel": "relaciona"}, {"alvo": "ipc", "confianca": 0.5, "epistemico": "fato", "origem": "wiki", "rel": "relaciona"}, {"alvo": "nucleo_broca_vivo", "confianca": 0.5, "epistemico": "fato", "origem": "wiki", "rel": "relaciona"}], "confianca": 1.0, "contraexemplos": [], "descricao": "Base viva incremental: fatos, inferências e conceitos no mesmo tecido Koch.", "epistemico": "fato", "exemplos": [], "id": "fase_2_tecido_de_conhecimento", "memoria": "semantica", "meta": {}, "origem": "CONHECIMENTO.md", "palavras_chave": [], "sinonimos": [], "tipo": "conceito", "titulo": "Fase 2 — Tecido de conhecimento"}
\section{Fase 2 — Tecido de conhecimento}
Base viva incremental: fatos, inferências e conceitos no mesmo tecido Koch.
\prov{relações: relaciona utilitarismo; relaciona rust; relaciona o\_que\_e\_no\_projecto; relaciona grupos; relaciona recursao; relaciona funil\_cumulativo; relaciona melhor\_candidato; relaciona funcao; relaciona pow\_bit\_taxa; relaciona mecanica\_quantica\_hub; relaciona geometria\_hub\_broca; relaciona pell; relaciona teste\_ordem; relaciona motivos; relaciona regua\_associativa\_e\_o\_risco\_de\_vazamento; relaciona topologia; relaciona ipc; relaciona nucleo\_broca\_vivo}
\prov{proveniência: CONHECIMENTO.md \textperiodcentered{} fato \textperiodcentered{} confiança 1.0 \textperiodcentered{} domínio (sem) \textperiodcentered{} história 19 versões no jornal}

%CRISTAL {"arestas": [{"alvo": "fase_2_tecido_de_conhecimento", "confianca": 1.0, "epistemico": "fato", "origem": "CONHECIMENTO.md", "rel": "parte_de"}], "confianca": 1.0, "contraexemplos": [], "descricao": "| Epistemico | Confiança típica | Origem | |------------|------------------|--------| | **fato** | ≥ 0.9 | usuário disse, documento, semilla | | **inferencia** | 0.5–0.7 | \"provavelmente\", modelo | | **hipotese** | < 0.5 | exploratório |", "epistemico": "fato", "exemplos": [], "id": "fato_vs_inferencia", "memoria": "semantica", "meta": {}, "origem": "CONHECIMENTO.md", "palavras_chave": [], "sinonimos": [], "tipo": "conceito", "titulo": "Fato vs inferência"}
\section{Fato vs inferência}
| Epistemico | Confiança típica | Origem | |------------|------------------|--------| | **fato** | \ensuremath{\ge} 0.9 | usuário disse, documento, semilla | | **inferencia** | 0.5–0.7 | "provavelmente", modelo | | **hipotese** | < 0.5 | exploratório |
\prov{relações: parte\_de fase\_2\_tecido\_de\_conhecimento}
\prov{proveniência: CONHECIMENTO.md \textperiodcentered{} fato \textperiodcentered{} confiança 1.0 \textperiodcentered{} domínio (sem) \textperiodcentered{} história 2 versões no jornal}

%CRISTAL {"arestas": [{"alvo": "kappa", "confianca": 0.95, "epistemico": "fato", "origem": "manual", "rel": "parte_de"}, {"alvo": "utilitarismo", "confianca": 0.5, "epistemico": "fato", "origem": "wiki", "rel": "relaciona"}, {"alvo": "parabola_isa", "confianca": 0.5, "epistemico": "fato", "origem": "wiki", "rel": "relaciona"}, {"alvo": "resultado_completo", "confianca": 0.5, "epistemico": "fato", "origem": "wiki", "rel": "relaciona"}, {"alvo": "memoria_virtual", "confianca": 0.5, "epistemico": "fato", "origem": "wiki", "rel": "relaciona"}, {"alvo": "word", "confianca": 0.5, "epistemico": "fato", "origem": "wiki", "rel": "relaciona"}], "confianca": 1.0, "contraexemplos": [], "descricao": "Fixamos (), =, ( = ) e o curinga.", "epistemico": "fato", "exemplos": [], "id": "fixamos_e_o_curinga", "memoria": "semantica", "meta": {}, "origem": "papers/casl-propagation.tex", "palavras_chave": [], "sinonimos": [], "tipo": "conceito", "titulo": "Fixamos (), =, ( = ) e o curinga."}
\section{Fixamos (), =, ( = ) e o curinga.}
Fixamos (), =, ( = ) e o curinga.
\prov{relações: parte\_de kappa; relaciona utilitarismo; relaciona parabola\_isa; relaciona resultado\_completo; relaciona memoria\_virtual; relaciona word}
\prov{proveniência: papers/casl-propagation.tex \textperiodcentered{} fato \textperiodcentered{} confiança 1.0 \textperiodcentered{} domínio (sem) \textperiodcentered{} história 67 versões no jornal}

%CRISTAL {"arestas": [{"alvo": "ponte_quantica_conjunto_dougherty_e_orbitais_de_hidrogenio", "confianca": 1.0, "epistemico": "fato", "origem": "ponte_quantica_dougherty.md", "rel": "parte_de"}, {"alvo": "fibracao_de_hopf_e_esfera_de_bloch", "confianca": 1.0, "epistemico": "fato", "origem": "fibracao_hopf_bloch.md", "rel": "parte_de"}, {"alvo": "correntes_de_birkeland_filamentos_de_plasma", "confianca": 1.0, "epistemico": "fato", "origem": "birkeland_correntes_plasma.md", "rel": "parte_de"}, {"alvo": "semiotica_triadica_de_charles_sanders_peirce", "confianca": 1.0, "epistemico": "fato", "origem": "semiotica_peirce_triade.md", "rel": "parte_de"}, {"alvo": "ginzburg_toro_helice_e_3d-sst", "confianca": 1.0, "epistemico": "fato", "origem": "ginzburg_toro_helice.md", "rel": "parte_de"}, {"alvo": "fase_geometrica_de_berry_pancharatnamberry", "confianca": 1.0, "epistemico": "fato", "origem": "fase_geometrica_berry.md", "rel": "parte_de"}, {"alvo": "colapso_koch_operador_ψ_da_tiffany", "confianca": 1.0, "epistemico": "fato", "origem": "colapso_koch_tiffany.md", "rel": "parte_de"}, {"alvo": "tiffany_cabeca_conversacional_broca-so", "confianca": 1.0, "epistemico": "fato", "origem": "tiffany_cabeca_broca.md", "rel": "parte_de"}, {"alvo": "geometria_hub_broca-so", "confianca": 1.0, "epistemico": "fato", "origem": "geometria_hub_broca.md", "rel": "parte_de"}], "confianca": 1.0, "contraexemplos": [], "descricao": "- `ula/BROCA_SO.md` - `papers/broca_os.tex` - `neuronio/primitivas_broca.py`", "epistemico": "fato", "exemplos": [], "id": "fontes", "memoria": "semantica", "meta": {"indexar": false, "verificacao": "slug idêntico — enriquecer, não duplicar"}, "origem": "manual", "palavras_chave": [], "sinonimos": [], "tipo": "conceito", "titulo": "Fontes"}
\section{Fontes}
- `ula/BROCA\_SO.md` - `papers/broca\_os.tex` - `neuronio/primitivas\_broca.py`
\prov{relações: parte\_de ponte\_quantica\_conjunto\_dougherty\_e\_orbitais\_de\_hidrogenio; parte\_de fibracao\_de\_hopf\_e\_esfera\_de\_bloch; parte\_de correntes\_de\_birkeland\_filamentos\_de\_plasma; parte\_de semiotica\_triadica\_de\_charles\_sanders\_peirce; parte\_de ginzburg\_toro\_helice\_e\_3d-sst; parte\_de fase\_geometrica\_de\_berry\_pancharatnamberry; parte\_de colapso\_koch\_operador\_\ensuremath{\psi}\_da\_tiffany; parte\_de tiffany\_cabeca\_conversacional\_broca-so; parte\_de geometria\_hub\_broca-so}
\prov{proveniência: manual \textperiodcentered{} fato \textperiodcentered{} confiança 1.0 \textperiodcentered{} domínio (sem) \textperiodcentered{} história 50 versões no jornal}

%CRISTAL {"arestas": [{"alvo": "mamifero", "confianca": 1.0, "epistemico": "fato", "origem": "semilla", "rel": "eh_um"}, {"alvo": "ponte_simbolica_gentil", "confianca": 1.0, "epistemico": "fato", "origem": "codigo", "rel": "confirma"}, {"alvo": "corpo_tropical", "confianca": 1.0, "epistemico": "fato", "origem": "codigo", "rel": "confirma"}, {"alvo": "a_decomposicao_de_peirce_celulas_e_canais", "confianca": 1.0, "epistemico": "fato", "origem": "codigo", "rel": "depende"}, {"alvo": "colapso", "confianca": 1.0, "epistemico": "fato", "origem": "codigo", "rel": "usa"}, {"alvo": "painel_estelar", "confianca": 1.0, "epistemico": "fato", "origem": "codigo", "rel": "explica"}, {"alvo": "matrizes", "confianca": 1.0, "epistemico": "fato", "origem": "codigo", "rel": "especializa"}], "confianca": 1.0, "contraexemplos": ["gato A ≠ gato felino (semilla mamífero)", "Operador A ≠ matriz identidade", "Medição gato ≠ triade semiótica Peirce completa", "Analogia Bloch/Dougherty ≠ implementação do operador A", "gato A ≠ gato felino (semilla)"], "descricao": "Operador algébrico 2×2 A=[[1,1],[1,0]] — medição PoW sobre bytes. Núcleo Broca-SO; homónimo mamífero (semilla) é contexto distinto.", "epistemico": "fato", "exemplos": ["operador A=[[1,1],[1,0]] no broca-so", "streaming byte a byte em neuronio.c", "A=[[1,1],[1,0]] — papers/computacao_fundamentos_broca.tex", "neuronio/neuronio.c — GATO_STREAM byte a byte", "ula/PAINEL_ESTELAR.md — três braços do operador", "ponte_simbolica_gentil — gato 2×2 → símbolos Gentil"], "id": "gato", "memoria": "semantica", "meta": {"confirmado": "gato_operador_hub", "fechado_fase7": true, "serve": "Operador algébrico PoW; não confundir com mamífero"}, "origem": "codigo", "palavras_chave": [], "sinonimos": [], "tipo": "conceito", "titulo": "gato"}
\section{gato}
Operador algébrico 2$\times$2 A=[[1,1],[1,0]] — medição PoW sobre bytes. Núcleo Broca-SO; homónimo mamífero (semilla) é contexto distinto.
\prov{exemplos: operador A=[[1,1],[1,0]] no broca-so; streaming byte a byte em neuronio.c; A=[[1,1],[1,0]] — papers/computacao\_fundamentos\_broca.tex; neuronio/neuronio.c — GATO\_STREAM byte a byte; ula/PAINEL\_ESTELAR.md — três braços do operador; ponte\_simbolica\_gentil — gato 2$\times$2 \ensuremath{\to} símbolos Gentil}
\prov{contraexemplos: gato A \ensuremath{\ne} gato felino (semilla mamífero); Operador A \ensuremath{\ne} matriz identidade; Medição gato \ensuremath{\ne} triade semiótica Peirce completa; Analogia Bloch/Dougherty \ensuremath{\ne} implementação do operador A; gato A \ensuremath{\ne} gato felino (semilla)}
\prov{relações: eh\_um mamifero; confirma ponte\_simbolica\_gentil; confirma corpo\_tropical; depende a\_decomposicao\_de\_peirce\_celulas\_e\_canais; usa colapso; explica painel\_estelar; especializa matrizes}
\prov{proveniência: codigo \textperiodcentered{} fato \textperiodcentered{} confiança 1.0 \textperiodcentered{} domínio (sem) \textperiodcentered{} história 304 versões no jornal}

%CRISTAL {"arestas": [{"alvo": "mamifero", "confianca": 0.92, "epistemico": "fato", "origem": "semilla", "rel": "eh_um"}, {"alvo": "o_que_e_no_projecto", "confianca": 0.5, "epistemico": "fato", "origem": "wiki", "rel": "relaciona"}, {"alvo": "vida-morte", "confianca": 0.5, "epistemico": "fato", "origem": "wiki", "rel": "relaciona"}, {"alvo": "word", "confianca": 0.5, "epistemico": "fato", "origem": "wiki", "rel": "relaciona"}, {"alvo": "vontade_de_poder", "confianca": 0.5, "epistemico": "fato", "origem": "wiki", "rel": "relaciona"}], "confianca": 0.92, "contraexemplos": [], "descricao": "Fato (fato): gato eh um mamífero. Origem: semilla.", "epistemico": "fato", "exemplos": [], "id": "gato_eh_um_mamifero", "memoria": "semantica", "meta": {"predicado": "mamífero", "sujeito": "gato"}, "origem": "semilla", "palavras_chave": ["gato", "mamífero"], "sinonimos": [], "tipo": "fato", "titulo": "gato eh um mamífero"}
\section{gato eh um mamífero}
Fato (fato): gato eh um mamífero. Origem: semilla.
\prov{relações: eh\_um mamifero; relaciona o\_que\_e\_no\_projecto; relaciona vida-morte; relaciona word; relaciona vontade\_de\_poder}
\prov{proveniência: semilla \textperiodcentered{} fato \textperiodcentered{} confiança 0.92 \textperiodcentered{} domínio (sem) \textperiodcentered{} história 5 versões no jornal}

%CRISTAL {"arestas": [{"alvo": "analise", "confianca": 1.0, "epistemico": "fato", "origem": "curriculo", "rel": "especializa"}, {"alvo": "espaco_vetorial", "confianca": 1.0, "epistemico": "fato", "origem": "paper", "rel": "depende"}, {"alvo": "topologia", "confianca": 1.0, "epistemico": "fato", "origem": "paper", "rel": "generaliza"}], "confianca": 1.0, "contraexemplos": ["Toro φ-Dougherty ≠ fibrado Hopf S³→S² (topologias distintas)", "PoW espectral ≠ acelerador SHA (Broca = coordenadas)", "Geometria riemanniana ≠ reticulado BAI (não isomorfos)", "Hipótese Ginzburg/Birkeland ≠ prova escala φ sem teste espectral"], "descricao": "Forma, distância, curvatura. Hub Broca: S³/Bloch/Hopf (fato), PoW espectral, hipóteses Dougherty/Ginzburg (analogia).", "epistemico": "fato", "exemplos": ["métrica de Sitter", "espaço euclidiano", "fibração de Hopf", "métrica de Sitter — a_metrica_de_sitter_global", "fibração Hopf S³→S² — fibracao_hopf_bloch", "PoW espectral calor/fase² — POW_GEOMETRIA.md", "fase geométrica Berry — ciclo adiabático Bloch"], "id": "geometria", "memoria": "semantica", "meta": {"serve": "Forma, distância, curvatura; ponte entre física e álgebra", "verificacao": "slug idêntico — enriquecer, não duplicar"}, "origem": "manual", "palavras_chave": [], "sinonimos": [], "tipo": "conceito", "titulo": "geometria"}
\section{geometria}
Forma, distância, curvatura. Hub Broca: S\ensuremath{^{3}}/Bloch/Hopf (fato), PoW espectral, hipóteses Dougherty/Ginzburg (analogia).
\prov{exemplos: métrica de Sitter; espaço euclidiano; fibração de Hopf; métrica de Sitter — a\_metrica\_de\_sitter\_global; fibração Hopf S\ensuremath{^{3}}\ensuremath{\to}S\ensuremath{^{2}} — fibracao\_hopf\_bloch; PoW espectral calor/fase\ensuremath{^{2}} — POW\_GEOMETRIA.md; fase geométrica Berry — ciclo adiabático Bloch}
\prov{contraexemplos: Toro \ensuremath{\varphi}-Dougherty \ensuremath{\ne} fibrado Hopf S\ensuremath{^{3}}\ensuremath{\to}S\ensuremath{^{2}} (topologias distintas); PoW espectral \ensuremath{\ne} acelerador SHA (Broca = coordenadas); Geometria riemanniana \ensuremath{\ne} reticulado BAI (não isomorfos); Hipótese Ginzburg/Birkeland \ensuremath{\ne} prova escala \ensuremath{\varphi} sem teste espectral}
\prov{relações: especializa analise; depende espaco\_vetorial; generaliza topologia}
\prov{proveniência: manual \textperiodcentered{} fato \textperiodcentered{} confiança 1.0 \textperiodcentered{} domínio (sem) \textperiodcentered{} história 172 versões no jornal}

%CRISTAL {"arestas": [{"alvo": "cifra_de_ouro_branco", "confianca": 1.0, "epistemico": "fato", "origem": "curriculo", "rel": "relaciona"}, {"alvo": "cifra", "confianca": 1.0, "epistemico": "fato", "origem": "paper", "rel": "usa"}, {"alvo": "broca_os", "confianca": 1.0, "epistemico": "fato", "origem": "curriculo", "rel": "usa"}, {"alvo": "blockchain", "confianca": 1.0, "epistemico": "fato", "origem": "codigo", "rel": "especializa"}, {"alvo": "valor", "confianca": 1.0, "epistemico": "fato", "origem": "codigo", "rel": "especializa"}, {"alvo": "dinheiro", "confianca": 1.0, "epistemico": "fato", "origem": "codigo", "rel": "especializa"}, {"alvo": "sha256", "confianca": 1.0, "epistemico": "fato", "origem": "codigo", "rel": "usa"}, {"alvo": "contabilidade", "confianca": 1.0, "epistemico": "fato", "origem": "codigo", "rel": "explica"}, {"alvo": "gato", "confianca": 1.0, "epistemico": "fato", "origem": "codigo", "rel": "usa"}], "confianca": 1.0, "contraexemplos": ["GoldChain ≠ DeFi genérico com gas ETH", "GoldChain ≠ blockchain pública qualquer — álgebra Broca", "Pell (cifra tesouraria) ≠ BTC on-chain directamente", "Mercado Broca ≠ especulação sem PoW/work precedente"], "descricao": "Álgebra dos contratos: valor, certificados sha256, liquidação (broca-so/goldchain/).", "epistemico": "fato", "exemplos": ["kernel/algebra_dos_contratos.tex", "goldchain/fluxo.py", "goldchain/fluxo.py — pipeline contrato→liquidação", "goldchain/producao.py — tecido tesouraria central", "kernel/algebra_dos_contratos.tex — formalização", "goldchain/README.md — PoW→Pell→BTC lastro"], "id": "goldchain", "memoria": "semantica", "meta": {"confirmado": "goldchain_mercado_hub", "fechado_fase7": true, "serve": "Núcleo mercado financeiro — contratos, certificação, liquidação"}, "origem": "codigo", "palavras_chave": ["GoldCoin", "Pell", "tesouraria"], "sinonimos": ["GoldChain", "gold chain"], "tipo": "conceito", "titulo": "goldchain"}
\section{goldchain}
Álgebra dos contratos: valor, certificados sha256, liquidação (broca-so/goldchain/).
\prov{exemplos: kernel/algebra\_dos\_contratos.tex; goldchain/fluxo.py; goldchain/fluxo.py — pipeline contrato\ensuremath{\to}liquidação; goldchain/producao.py — tecido tesouraria central; kernel/algebra\_dos\_contratos.tex — formalização; goldchain/README.md — PoW\ensuremath{\to}Pell\ensuremath{\to}BTC lastro}
\prov{contraexemplos: GoldChain \ensuremath{\ne} DeFi genérico com gas ETH; GoldChain \ensuremath{\ne} blockchain pública qualquer — álgebra Broca; Pell (cifra tesouraria) \ensuremath{\ne} BTC on-chain directamente; Mercado Broca \ensuremath{\ne} especulação sem PoW/work precedente}
\prov{sinónimos: GoldChain; gold chain}
\prov{relações: relaciona cifra\_de\_ouro\_branco; usa cifra; usa broca\_os; especializa blockchain; especializa valor; especializa dinheiro; usa sha256; explica contabilidade; usa gato}
\prov{proveniência: codigo \textperiodcentered{} fato \textperiodcentered{} confiança 1.0 \textperiodcentered{} domínio (sem) \textperiodcentered{} história 265 versões no jornal}

%CRISTAL {"arestas": [{"alvo": "reticulado", "confianca": 0.039, "epistemico": "fato", "origem": "manual", "rel": "referencia"}, {"alvo": "bai", "confianca": 0.07, "epistemico": "fato", "origem": "manual", "rel": "referencia"}, {"alvo": "reticulados", "confianca": 0.95, "epistemico": "fato", "origem": "manual", "rel": "usa"}, {"alvo": "bai_biblioteca_de_autorizacao_isomorfica", "confianca": 0.95, "epistemico": "fato", "origem": "manual", "rel": "parte_de"}], "confianca": 1.0, "contraexemplos": [], "descricao": "[Grafo de delegação] Um é um (não necessariamente acíclico) cujos vértices são principals e cujas arestas u w representam `` u concede a capacidade a w ''. w é principal --- sem exigência de árvore, de vizinhança ou de aciclicidade.", "epistemico": "fato", "exemplos": [], "id": "grafo_de_delegacao_um_e_um_nao_necessariamente_aciclico", "memoria": "semantica", "meta": {}, "origem": "papers/casl-propagation.tex", "palavras_chave": [], "sinonimos": [], "tipo": "conceito", "titulo": "[Grafo de delegação] Um é um (não necessariamente acíclico) …"}
\section{[Grafo de delegação] Um é um (não necessariamente acíclico) …}
[Grafo de delegação] Um é um (não necessariamente acíclico) cujos vértices são principals e cujas arestas u w representam `` u concede a capacidade a w ''. w é principal --- sem exigência de árvore, de vizinhança ou de aciclicidade.
\prov{relações: referencia reticulado; referencia bai; usa reticulados; parte\_de bai\_biblioteca\_de\_autorizacao\_isomorfica}
\prov{proveniência: papers/casl-propagation.tex \textperiodcentered{} fato \textperiodcentered{} confiança 1.0 \textperiodcentered{} domínio (sem) \textperiodcentered{} história 125 versões no jornal}

%CRISTAL {"arestas": [{"alvo": "kernel", "confianca": 0.734, "epistemico": "fato", "origem": "manual", "rel": "referencia"}, {"alvo": "koch_atlas", "confianca": 0.5, "epistemico": "fato", "origem": "wiki", "rel": "relaciona"}, {"alvo": "pow_gato", "confianca": 0.5, "epistemico": "fato", "origem": "wiki", "rel": "relaciona"}, {"alvo": "sessao_interativa", "confianca": 0.5, "epistemico": "fato", "origem": "wiki", "rel": "relaciona"}, {"alvo": "thread", "confianca": 0.5, "epistemico": "fato", "origem": "wiki", "rel": "relaciona"}, {"alvo": "prova_isa", "confianca": 0.5, "epistemico": "fato", "origem": "wiki", "rel": "relaciona"}, {"alvo": "vida-morte", "confianca": 0.5, "epistemico": "fato", "origem": "wiki", "rel": "relaciona"}, {"alvo": "resposta_a_pergunta_dificil", "confianca": 0.5, "epistemico": "fato", "origem": "wiki", "rel": "relaciona"}, {"alvo": "linker", "confianca": 0.5, "epistemico": "fato", "origem": "wiki", "rel": "relaciona"}, {"alvo": "hopfield_rede_hub", "confianca": 0.5, "epistemico": "fato", "origem": "wiki", "rel": "relaciona"}], "confianca": 1.0, "contraexemplos": [], "descricao": "```", "epistemico": "fato", "exemplos": [], "id": "grava_fato_resposta_normal_em_seguida", "memoria": "semantica", "meta": {}, "origem": "CONHECIMENTO.md", "palavras_chave": [], "sinonimos": [], "tipo": "conceito", "titulo": "grava fato; resposta normal em seguida"}
\section{grava fato; resposta normal em seguida}
```
\prov{relações: referencia kernel; relaciona koch\_atlas; relaciona pow\_gato; relaciona sessao\_interativa; relaciona thread; relaciona prova\_isa; relaciona vida-morte; relaciona resposta\_a\_pergunta\_dificil; relaciona linker; relaciona hopfield\_rede\_hub}
\prov{proveniência: CONHECIMENTO.md \textperiodcentered{} fato \textperiodcentered{} confiança 1.0 \textperiodcentered{} domínio (sem) \textperiodcentered{} história 13 versões no jornal}

%CRISTAL {"arestas": [{"alvo": "tiffany", "confianca": 1.0, "epistemico": "fato", "origem": "curriculo", "rel": "especializa"}, {"alvo": "broca_os", "confianca": 1.0, "epistemico": "fato", "origem": "paper", "rel": "explica"}, {"alvo": "mcculloch_pitts", "confianca": 1.0, "epistemico": "fato", "origem": "codigo", "rel": "especializa"}, {"alvo": "gato", "confianca": 1.0, "epistemico": "fato", "origem": "codigo", "rel": "especializa"}, {"alvo": "inteligencia_artificial", "confianca": 1.0, "epistemico": "fato", "origem": "codigo", "rel": "especializa"}, {"alvo": "koch_como_recipiente_memoria_nao-associativa", "confianca": 1.0, "epistemico": "fato", "origem": "wiki", "rel": "relaciona"}, {"alvo": "floco_phi", "confianca": 1.0, "epistemico": "fato", "origem": "wiki", "rel": "relaciona"}], "confianca": 1.0, "contraexemplos": ["Hopfield ≠ backprop / deep learning", "Hopfield ≠ Tiffany — Tiffany acrescenta colapso Ψ + BAI", "Hopfield ≠ LLM autoregressivo", "Rede histórica ≠ implementação completa da cabeça Tiffany"], "descricao": "Rede associativa Hopfield (1982): memória content-addressable — recall descendo a energia ao atrator, sinapse Hebbian (ξξᵀ). No Broca: McCulloch-Pitts → Hopfield → Tiffany; confirma o operador gato A. É a metade ASSOCIATIVA da memória do sistema — a reta 𝓛(w) que RECONHECE por ressonância (overlap). Sua contraparte ORTOGONAL é o recipiente não-associativo: o floco Φ, que guarda justamente o que o recall associativo dissipa (ver koch_dissipa_nao_codifica, floco_phi). Aterrado no metal em reconhece.c.", "epistemico": "fato", "exemplos": ["currículo IA — mcculloch_pitts → hopfield → tiffany", "papers/tiffany_cabeca_broca.md — cadeia associativa", "neuronio/experimentos_tiffany.py — overlap sobre tecido", "broca_os — rede=Hopfield associativa"], "id": "hopfield", "memoria": "semantica", "meta": {"confirmado": "hopfield_rede_hub", "fechado_fase7": true, "serve": "Memória associativa — predecessor da Tiffany"}, "origem": "codigo", "palavras_chave": ["atrator", "Hebb", "pattern completion"], "sinonimos": ["Hopfield", "rede de Hopfield"], "tipo": "conceito", "titulo": "hopfield"}
\section{hopfield}
Rede associativa Hopfield (1982): memória content-addressable — recall descendo a energia ao atrator, sinapse Hebbian (\ensuremath{\xi}\ensuremath{\xi}\ensuremath{^{T}}). No Broca: McCulloch-Pitts \ensuremath{\to} Hopfield \ensuremath{\to} Tiffany; confirma o operador gato A. É a metade ASSOCIATIVA da memória do sistema — a reta \ensuremath{\mathcal{L}}(w) que RECONHECE por ressonância (overlap). Sua contraparte ORTOGONAL é o recipiente não-associativo: o floco \ensuremath{\Phi}, que guarda justamente o que o recall associativo dissipa (ver koch\_dissipa\_nao\_codifica, floco\_phi). Aterrado no metal em reconhece.c.
\prov{exemplos: currículo IA — mcculloch\_pitts \ensuremath{\to} hopfield \ensuremath{\to} tiffany; papers/tiffany\_cabeca\_broca.md — cadeia associativa; neuronio/experimentos\_tiffany.py — overlap sobre tecido; broca\_os — rede=Hopfield associativa}
\prov{contraexemplos: Hopfield \ensuremath{\ne} backprop / deep learning; Hopfield \ensuremath{\ne} Tiffany — Tiffany acrescenta colapso \ensuremath{\Psi} + BAI; Hopfield \ensuremath{\ne} LLM autoregressivo; Rede histórica \ensuremath{\ne} implementação completa da cabeça Tiffany}
\prov{sinónimos: Hopfield; rede de Hopfield}
\prov{relações: especializa tiffany; explica broca\_os; especializa mcculloch\_pitts; especializa gato; especializa inteligencia\_artificial; relaciona koch\_como\_recipiente\_memoria\_nao-associativa; relaciona floco\_phi}
\prov{proveniência: codigo \textperiodcentered{} fato \textperiodcentered{} confiança 1.0 \textperiodcentered{} domínio (sem) \textperiodcentered{} história 177 versões no jornal}

%CRISTAL {"arestas": [{"alvo": "bai", "confianca": -0.109, "epistemico": "fato", "origem": "manual", "rel": "referencia"}, {"alvo": "semantica", "confianca": 0.95, "epistemico": "fato", "origem": "manual", "rel": "especializa"}, {"alvo": "bai_biblioteca_de_autorizacao_isomorfica", "confianca": 0.95, "epistemico": "fato", "origem": "manual", "rel": "parte_de"}], "confianca": 1.0, "contraexemplos": [], "descricao": "[Identificação semântica --- nunca sintática] Toda é identificada ao seu: os construtores sintáticos do CASL mapeiam-se nas operações da álgebra, [: [c₁,c₂] = c₁ c₂,: [c₁,c₂] = c₁ c₂,: c = c. ] Raciocinamos sempre sobre, jamais sobre a forma da query --- duas queries sintaticamente distintas com o mesmo extent são, para nós, a condição. (Isto não trivializa a implementação: o teste de inclusão no é conservador justamente porque decidir sintaticamente é incompleto; ver", "epistemico": "fato", "exemplos": [], "id": "identificacao_semantica_---_nunca_sintatica_toda_e_identif", "memoria": "semantica", "meta": {}, "origem": "papers/casl-propagation.tex", "palavras_chave": [], "sinonimos": [], "tipo": "conceito", "titulo": "[Identificação semântica --- nunca sintática] Toda é identif…"}
\section{[Identificação semântica --- nunca sintática] Toda é identif…}
[Identificação semântica --- nunca sintática] Toda é identificada ao seu: os construtores sintáticos do CASL mapeiam-se nas operações da álgebra, [: [c\ensuremath{_{1}},c\ensuremath{_{2}}] = c\ensuremath{_{1}} c\ensuremath{_{2}},: [c\ensuremath{_{1}},c\ensuremath{_{2}}] = c\ensuremath{_{1}} c\ensuremath{_{2}},: c = c. ] Raciocinamos sempre sobre, jamais sobre a forma da query --- duas queries sintaticamente distintas com o mesmo extent são, para nós, a condição. (Isto não trivializa a implementação: o teste de inclusão no é conservador justamente porque decidir sintaticamente é incompleto; ver
\prov{relações: referencia bai; especializa semantica; parte\_de bai\_biblioteca\_de\_autorizacao\_isomorfica}
\prov{proveniência: papers/casl-propagation.tex \textperiodcentered{} fato \textperiodcentered{} confiança 1.0 \textperiodcentered{} domínio (sem) \textperiodcentered{} história 120 versões no jornal}

%CRISTAL {"arestas": [{"alvo": "casl-propagation", "confianca": 0.95, "epistemico": "fato", "origem": "manual", "rel": "parte_de"}, {"alvo": "broca-so_camada_de_trabalho_isomorfica", "confianca": 0.95, "epistemico": "fato", "origem": "manual", "rel": "referencia"}], "confianca": 1.0, "contraexemplos": [], "descricao": "[Instanciações] Parte~II: =(a,s) , . Parte~III: é identificador de nó no tecido. Parte~IV: é um turno dialogado .", "epistemico": "fato", "exemplos": [], "id": "instanciacoes_parteii_as_parteiii_e_identifica", "memoria": "semantica", "meta": {}, "origem": "papers/casl-propagation.tex", "palavras_chave": [], "sinonimos": [], "tipo": "conceito", "titulo": "[Instanciações] Parte~II: =(a,s) , . Parte~III: é identifica…"}
\section{[Instanciações] Parte\~{}II: =(a,s) , . Parte\~{}III: é identifica…}
[Instanciações] Parte\~{}II: =(a,s) , . Parte\~{}III: é identificador de nó no tecido. Parte\~{}IV: é um turno dialogado .
\prov{relações: parte\_de casl-propagation; referencia broca-so\_camada\_de\_trabalho\_isomorfica}
\prov{proveniência: papers/casl-propagation.tex \textperiodcentered{} fato \textperiodcentered{} confiança 1.0 \textperiodcentered{} domínio (sem) \textperiodcentered{} história 116 versões no jornal}

%CRISTAL {"arestas": [{"alvo": "por_sessao_conversadadostelemetriasessaojsonl", "confianca": 1.0, "epistemico": "fato", "origem": "CASL.md", "rel": "parte_de"}, {"alvo": "word", "confianca": 0.5, "epistemico": "fato", "origem": "wiki", "rel": "relaciona"}, {"alvo": "virtualizacao", "confianca": 0.5, "epistemico": "fato", "origem": "wiki", "rel": "relaciona"}, {"alvo": "syscall", "confianca": 0.5, "epistemico": "fato", "origem": "wiki", "rel": "relaciona"}], "confianca": 1.0, "contraexemplos": [], "descricao": "Ver [`README.md`](README.md) para instrumentos (`observar.py`, `cadeias.py`, `mede_flocos.py`).", "epistemico": "fato", "exemplos": [], "id": "lab_conversa", "memoria": "semantica", "meta": {}, "origem": "CASL.md", "palavras_chave": [], "sinonimos": [], "tipo": "conceito", "titulo": "Lab conversa/"}
\section{Lab conversa/}
Ver [`README.md`](README.md) para instrumentos (`observar.py`, `cadeias.py`, `mede\_flocos.py`).
\prov{relações: parte\_de por\_sessao\_conversadadostelemetriasessaojsonl; relaciona word; relaciona virtualizacao; relaciona syscall}
\prov{proveniência: CASL.md \textperiodcentered{} fato \textperiodcentered{} confiança 1.0 \textperiodcentered{} domínio (sem) \textperiodcentered{} história 5 versões no jornal}

%CRISTAL {"arestas": [{"alvo": "fase_2_tecido_de_conhecimento", "confianca": 1.0, "epistemico": "fato", "origem": "CONHECIMENTO.md", "rel": "parte_de"}, {"alvo": "reticulado", "confianca": 0.85, "epistemico": "fato", "origem": "manual", "rel": "referencia"}], "confianca": 1.0, "contraexemplos": [], "descricao": "```bash python3 conversa/conhecimento/lab.py semilla python3 conversa/conhecimento/lab.py conceito reticulado python3 conversa/conhecimento/lab.py grafo reticulado python3 conversa/conhecimento/lab.py aprender \"O gato é um mamífero.\" python3 conversa/conhecimento/lab.py sugerir reticulado python3 conversa/conhecimento/lab.py conectar gato mamifero --rel eh_um --auto python3 conversa/conhecimento/lab.py porque \"o que é reticulado\"", "epistemico": "fato", "exemplos": [], "id": "laboratorio", "memoria": "semantica", "meta": {}, "origem": "CONHECIMENTO.md", "palavras_chave": [], "sinonimos": [], "tipo": "conceito", "titulo": "Laboratório"}
\section{Laboratório}
```bash python3 conversa/conhecimento/lab.py semilla python3 conversa/conhecimento/lab.py conceito reticulado python3 conversa/conhecimento/lab.py grafo reticulado python3 conversa/conhecimento/lab.py aprender "O gato é um mamífero." python3 conversa/conhecimento/lab.py sugerir reticulado python3 conversa/conhecimento/lab.py conectar gato mamifero --rel eh\_um --auto python3 conversa/conhecimento/lab.py porque "o que é reticulado"
\prov{relações: parte\_de fase\_2\_tecido\_de\_conhecimento; referencia reticulado}
\prov{proveniência: CONHECIMENTO.md \textperiodcentered{} fato \textperiodcentered{} confiança 1.0 \textperiodcentered{} domínio (sem) \textperiodcentered{} história 5 versões no jornal}

%CRISTAL {"arestas": [{"alvo": "bai", "confianca": 0.047, "epistemico": "fato", "origem": "manual", "rel": "referencia"}, {"alvo": "bai_biblioteca_de_autorizacao_isomorfica", "confianca": 0.95, "epistemico": "fato", "origem": "manual", "rel": "explica"}], "confianca": 1.0, "contraexemplos": [], "descricao": "[Leituras por domínio --- sem alterar a definição] @ ll@ CASL & re-delegação de permissões conversa / memória & profundidade de cadeia ou persistência replicação & alcance em grafo distribuído", "epistemico": "fato", "exemplos": [], "id": "leituras_por_dominio_---_sem_alterar_a_definicao_ll_cas", "memoria": "semantica", "meta": {}, "origem": "papers/casl-propagation.tex", "palavras_chave": [], "sinonimos": [], "tipo": "conceito", "titulo": "[Leituras por domínio --- sem alterar a definição] @ ll@ CAS…"}
\section{[Leituras por domínio --- sem alterar a definição] @ ll@ CAS…}
[Leituras por domínio --- sem alterar a definição] @ ll@ CASL \& re-delegação de permissões conversa / memória \& profundidade de cadeia ou persistência replicação \& alcance em grafo distribuído
\prov{relações: referencia bai; explica bai\_biblioteca\_de\_autorizacao\_isomorfica}
\prov{proveniência: papers/casl-propagation.tex \textperiodcentered{} fato \textperiodcentered{} confiança 1.0 \textperiodcentered{} domínio (sem) \textperiodcentered{} história 64 versões no jornal}

%CRISTAL {"arestas": [{"alvo": "kappa", "confianca": 0.95, "epistemico": "fato", "origem": "manual", "rel": "parte_de"}], "confianca": 1.0, "contraexemplos": [], "descricao": "[Linhas e condições] Seja o universo de todas as linhas. Cada r tem atributos ( r., r., ). Uma c é um predicado c: identificado ao seu c = r:c(r). As condições formam a álgebra booleana (2^, ) --- as do CASL.", "epistemico": "fato", "exemplos": [], "id": "linhas_e_condicoes_seja_o_universo_de_todas_as_linhas_cad", "memoria": "semantica", "meta": {}, "origem": "papers/casl-propagation.tex", "palavras_chave": [], "sinonimos": [], "tipo": "conceito", "titulo": "[Linhas e condições] Seja o universo de todas as linhas. Cad…"}
\section{[Linhas e condições] Seja o universo de todas as linhas. Cad…}
[Linhas e condições] Seja o universo de todas as linhas. Cada r tem atributos ( r., r., ). Uma c é um predicado c: identificado ao seu c = r:c(r). As condições formam a álgebra booleana (2\^{}, ) --- as do CASL.
\prov{relações: parte\_de kappa}
\prov{proveniência: papers/casl-propagation.tex \textperiodcentered{} fato \textperiodcentered{} confiança 1.0 \textperiodcentered{} domínio (sem) \textperiodcentered{} história 62 versões no jornal}

%CRISTAL {"arestas": [{"alvo": "kernel", "confianca": 0.672, "epistemico": "fato", "origem": "manual", "rel": "referencia"}, {"alvo": "transcricao", "confianca": 0.5, "epistemico": "fato", "origem": "wiki", "rel": "relaciona"}, {"alvo": "teste_ordem", "confianca": 0.5, "epistemico": "fato", "origem": "wiki", "rel": "relaciona"}, {"alvo": "parabola_pa", "confianca": 0.5, "epistemico": "fato", "origem": "wiki", "rel": "relaciona"}, {"alvo": "testes", "confianca": 0.5, "epistemico": "fato", "origem": "wiki", "rel": "relaciona"}, {"alvo": "por_sessao_conversadadostelemetriasessaojsonl", "confianca": 0.5, "epistemico": "fato", "origem": "wiki", "rel": "relaciona"}, {"alvo": "resposta_a_pergunta_dificil", "confianca": 0.5, "epistemico": "fato", "origem": "wiki", "rel": "relaciona"}, {"alvo": "topologia", "confianca": 0.5, "epistemico": "fato", "origem": "wiki", "rel": "relaciona"}, {"alvo": "vida-morte", "confianca": 0.5, "epistemico": "fato", "origem": "wiki", "rel": "relaciona"}, {"alvo": "tiffany_cabeca_broca", "confianca": 0.5, "epistemico": "fato", "origem": "wiki", "rel": "relaciona"}, {"alvo": "metais", "confianca": 0.5, "epistemico": "fato", "origem": "wiki", "rel": "relaciona"}, {"alvo": "motivos", "confianca": 0.5, "epistemico": "fato", "origem": "wiki", "rel": "relaciona"}, {"alvo": "pow_metal_funil", "confianca": 0.5, "epistemico": "fato", "origem": "wiki", "rel": "relaciona"}, {"alvo": "projeto_cristaljson", "confianca": 0.5, "epistemico": "fato", "origem": "wiki", "rel": "relaciona"}, {"alvo": "pow_metal_pre_fecho", "confianca": 0.5, "epistemico": "fato", "origem": "wiki", "rel": "relaciona"}, {"alvo": "pow_metal_setor", "confianca": 0.5, "epistemico": "fato", "origem": "wiki", "rel": "relaciona"}, {"alvo": "pow_geometria", "confianca": 0.5, "epistemico": "fato", "origem": "wiki", "rel": "relaciona"}, {"alvo": "pow_metal_ressonancia", "confianca": 0.5, "epistemico": "fato", "origem": "wiki", "rel": "relaciona"}, {"alvo": "pilha", "confianca": 0.5, "epistemico": "fato", "origem": "wiki", "rel": "relaciona"}, {"alvo": "numa", "confianca": 0.5, "epistemico": "fato", "origem": "wiki", "rel": "relaciona"}, {"alvo": "ponteiro", "confianca": 0.5, "epistemico": "fato", "origem": "wiki", "rel": "relaciona"}, {"alvo": "parabola_destruir", "confianca": 0.5, "epistemico": "fato", "origem": "wiki", "rel": "relaciona"}, {"alvo": "virtualizacao", "confianca": 0.5, "epistemico": "fato", "origem": "wiki", "rel": "relaciona"}, {"alvo": "syscall", "confianca": 0.5, "epistemico": "fato", "origem": "wiki", "rel": "relaciona"}], "confianca": 1.0, "contraexemplos": [], "descricao": "Secção de CASL.md.", "epistemico": "fato", "exemplos": [], "id": "log_legado_conversadadostelemetriajsonl", "memoria": "semantica", "meta": {"indexar": false, "verificacao": "slug idêntico — enriquecer, não duplicar"}, "origem": "manual", "palavras_chave": [], "sinonimos": [], "tipo": "conceito", "titulo": "log_legado_conversadadostelemetriajsonl"}
\section{log\_legado\_conversadadostelemetriajsonl}
Secção de CASL.md.
\prov{relações: referencia kernel; relaciona transcricao; relaciona teste\_ordem; relaciona parabola\_pa; relaciona testes; relaciona por\_sessao\_conversadadostelemetriasessaojsonl; relaciona resposta\_a\_pergunta\_dificil; relaciona topologia; relaciona vida-morte; relaciona tiffany\_cabeca\_broca; relaciona metais; relaciona motivos; relaciona pow\_metal\_funil; relaciona projeto\_cristaljson; relaciona pow\_metal\_pre\_fecho; relaciona pow\_metal\_setor; relaciona pow\_geometria; relaciona pow\_metal\_ressonancia; relaciona pilha; relaciona numa; relaciona ponteiro; relaciona parabola\_destruir; relaciona virtualizacao; relaciona syscall}
\prov{proveniência: manual \textperiodcentered{} fato \textperiodcentered{} confiança 1.0 \textperiodcentered{} domínio (sem) \textperiodcentered{} história 83 versões no jornal}

%CRISTAL {"arestas": [{"alvo": "python", "confianca": 0.5, "epistemico": "fato", "origem": "wiki", "rel": "relaciona"}, {"alvo": "prova_isa", "confianca": 0.5, "epistemico": "fato", "origem": "wiki", "rel": "relaciona"}, {"alvo": "memoria_virtual", "confianca": 0.5, "epistemico": "fato", "origem": "wiki", "rel": "relaciona"}, {"alvo": "pow_gato", "confianca": 0.5, "epistemico": "fato", "origem": "wiki", "rel": "relaciona"}, {"alvo": "pow_espectral", "confianca": 0.5, "epistemico": "fato", "origem": "wiki", "rel": "relaciona"}, {"alvo": "word", "confianca": 0.5, "epistemico": "fato", "origem": "wiki", "rel": "relaciona"}, {"alvo": "virtualizacao", "confianca": 0.5, "epistemico": "fato", "origem": "wiki", "rel": "relaciona"}], "confianca": 1.0, "contraexemplos": [], "descricao": "Classe de vertebrados homeotermos com glândulas mamárias.", "epistemico": "fato", "exemplos": [], "id": "mamifero", "memoria": "semantica", "meta": {}, "origem": "semilla", "palavras_chave": [], "sinonimos": [], "tipo": "conceito", "titulo": "mamifero"}
\section{mamifero}
Classe de vertebrados homeotermos com glândulas mamárias.
\prov{relações: relaciona python; relaciona prova\_isa; relaciona memoria\_virtual; relaciona pow\_gato; relaciona pow\_espectral; relaciona word; relaciona virtualizacao}
\prov{proveniência: semilla \textperiodcentered{} fato \textperiodcentered{} confiança 1.0 \textperiodcentered{} domínio (sem) \textperiodcentered{} história 8 versões no jornal}

%CRISTAL {"arestas": [{"alvo": "vetores", "confianca": 1.0, "epistemico": "fato", "origem": "paper", "rel": "usa"}, {"alvo": "analise", "confianca": 1.0, "epistemico": "fato", "origem": "paper", "rel": "depende"}, {"alvo": "utilitarismo", "confianca": 0.5, "epistemico": "fato", "origem": "wiki", "rel": "relaciona"}, {"alvo": "pedagogia", "confianca": 0.5, "epistemico": "fato", "origem": "wiki", "rel": "relaciona"}, {"alvo": "tipos_de_interpretante", "confianca": 0.5, "epistemico": "fato", "origem": "wiki", "rel": "relaciona"}, {"alvo": "virtualizacao", "confianca": 0.5, "epistemico": "fato", "origem": "wiki", "rel": "relaciona"}, {"alvo": "scheduler", "confianca": 0.5, "epistemico": "fato", "origem": "wiki", "rel": "relaciona"}, {"alvo": "vida-morte", "confianca": 0.5, "epistemico": "fato", "origem": "wiki", "rel": "relaciona"}, {"alvo": "ruido", "confianca": 0.5, "epistemico": "fato", "origem": "wiki", "rel": "relaciona"}, {"alvo": "transcricao", "confianca": 0.5, "epistemico": "fato", "origem": "wiki", "rel": "relaciona"}, {"alvo": "resultado_anterior_identidade_nao_cobertura", "confianca": 0.5, "epistemico": "fato", "origem": "wiki", "rel": "relaciona"}, {"alvo": "prova_da_maquina_espectral", "confianca": 0.5, "epistemico": "fato", "origem": "wiki", "rel": "relaciona"}, {"alvo": "veredito_experimental", "confianca": 0.5, "epistemico": "fato", "origem": "wiki", "rel": "relaciona"}, {"alvo": "pitagorismo", "confianca": 0.5, "epistemico": "fato", "origem": "wiki", "rel": "relaciona"}, {"alvo": "paper_2_reta_ouro", "confianca": 0.5, "epistemico": "fato", "origem": "wiki", "rel": "relaciona"}, {"alvo": "resultado_completo", "confianca": 0.5, "epistemico": "fato", "origem": "wiki", "rel": "relaciona"}, {"alvo": "modo_espectral_spect", "confianca": 0.5, "epistemico": "fato", "origem": "wiki", "rel": "relaciona"}, {"alvo": "pow_metal_ressonancia_val", "confianca": 0.5, "epistemico": "fato", "origem": "wiki", "rel": "relaciona"}, {"alvo": "teoria_do_valor", "confianca": 0.5, "epistemico": "fato", "origem": "wiki", "rel": "relaciona"}], "confianca": 1.0, "contraexemplos": [], "descricao": "Paper: mecanica.tex", "epistemico": "fato", "exemplos": ["Lei de Newton", "sistema massa-mola"], "id": "mecanica", "memoria": "semantica", "meta": {"serve": "Movimento e forças; usa vetores e análise", "verificacao": "slug idêntico — enriquecer, não duplicar"}, "origem": "manual", "palavras_chave": [], "sinonimos": [], "tipo": "conceito", "titulo": "mecanica"}
\section{mecanica}
Paper: mecanica.tex
\prov{exemplos: Lei de Newton; sistema massa-mola}
\prov{relações: usa vetores; depende analise; relaciona utilitarismo; relaciona pedagogia; relaciona tipos\_de\_interpretante; relaciona virtualizacao; relaciona scheduler; relaciona vida-morte; relaciona ruido; relaciona transcricao; relaciona resultado\_anterior\_identidade\_nao\_cobertura; relaciona prova\_da\_maquina\_espectral; relaciona veredito\_experimental; relaciona pitagorismo; relaciona paper\_2\_reta\_ouro; relaciona resultado\_completo; relaciona modo\_espectral\_spect; relaciona pow\_metal\_ressonancia\_val; relaciona teoria\_do\_valor}
\prov{proveniência: manual \textperiodcentered{} fato \textperiodcentered{} confiança 1.0 \textperiodcentered{} domínio (sem) \textperiodcentered{} história 172 versões no jornal}

%CRISTAL {"arestas": [{"alvo": "corpo_dual", "confianca": 1.0, "epistemico": "fato", "origem": "paper", "rel": "parte_de"}, {"alvo": "termodinamica", "confianca": 1.0, "epistemico": "fato", "origem": "curriculo", "rel": "continua"}, {"alvo": "geometria", "confianca": 1.0, "epistemico": "fato", "origem": "curriculo", "rel": "referencia"}, {"alvo": "matematica", "confianca": 1.0, "epistemico": "fato", "origem": "codigo", "rel": "depende"}, {"alvo": "mecanica", "confianca": 1.0, "epistemico": "fato", "origem": "codigo", "rel": "especializa"}, {"alvo": "fisica", "confianca": 1.0, "epistemico": "fato", "origem": "codigo", "rel": "parte_de"}, {"alvo": "fibracao_hopf_bloch", "confianca": 1.0, "epistemico": "fato", "origem": "codigo", "rel": "explica"}, {"alvo": "microscopia_fotoionizacao_hidrogenio", "confianca": 1.0, "epistemico": "fato", "origem": "codigo", "rel": "explica"}, {"alvo": "fase_geometrica_berry", "confianca": 1.0, "epistemico": "fato", "origem": "codigo", "rel": "explica"}, {"alvo": "gato", "confianca": 1.0, "epistemico": "fato", "origem": "codigo", "rel": "analogia"}], "confianca": 1.0, "contraexemplos": ["Hipótese Dougherty ≠ MQ standard — toro φ é conjectura", "Analogia gato A ≠ qubit de laboratório", "Hopf/Bloch confirmados ≠ prova escala φ-Dougherty", "Colapso Ψ Broca ≠ colapso postulado da função de onda"], "descricao": "Estados, operadores, |ψ|² — física standard. No Broca: Hopf/Bloch, Stodolna H, Berry confirmados (fato); Dougherty φ = hipótese analogia; gato A = analogia espectral.", "epistemico": "fato", "exemplos": ["gato A", "APS v6/58 orbitais H", "ψ_nlm standard", "papers/fibracao_hopf_bloch.md — S³→S², Bloch", "papers/microscopia_fotoionizacao_hidrogenio.md — PRL 2013 nodos H", "papers/fase_geometrica_berry.md — fase geométrica adiabática", "ψ_nlm standard + gato A (analogia PoW)"], "id": "mecanica_quantica", "memoria": "semantica", "meta": {"confirmado": "mecanica_quantica_hub", "fechado_fase7": true, "serve": "Domínio físico-quântico; pontes verificadas vs hipóteses φ"}, "origem": "codigo", "palavras_chave": ["operador", "função de onda", "Peirce"], "sinonimos": ["mecânica quântica", "MQ"], "tipo": "conceito", "titulo": "mecanica_quantica"}
\section{mecanica\_quantica}
Estados, operadores, |\ensuremath{\psi}|\ensuremath{^{2}} — física standard. No Broca: Hopf/Bloch, Stodolna H, Berry confirmados (fato); Dougherty \ensuremath{\varphi} = hipótese analogia; gato A = analogia espectral.
\prov{exemplos: gato A; APS v6/58 orbitais H; \ensuremath{\psi}\_nlm standard; papers/fibracao\_hopf\_bloch.md — S\ensuremath{^{3}}\ensuremath{\to}S\ensuremath{^{2}}, Bloch; papers/microscopia\_fotoionizacao\_hidrogenio.md — PRL 2013 nodos H; papers/fase\_geometrica\_berry.md — fase geométrica adiabática; \ensuremath{\psi}\_nlm standard + gato A (analogia PoW)}
\prov{contraexemplos: Hipótese Dougherty \ensuremath{\ne} MQ standard — toro \ensuremath{\varphi} é conjectura; Analogia gato A \ensuremath{\ne} qubit de laboratório; Hopf/Bloch confirmados \ensuremath{\ne} prova escala \ensuremath{\varphi}-Dougherty; Colapso \ensuremath{\Psi} Broca \ensuremath{\ne} colapso postulado da função de onda}
\prov{sinónimos: mecânica quântica; MQ}
\prov{relações: parte\_de corpo\_dual; continua termodinamica; referencia geometria; depende matematica; especializa mecanica; parte\_de fisica; explica fibracao\_hopf\_bloch; explica microscopia\_fotoionizacao\_hidrogenio; explica fase\_geometrica\_berry; analogia gato}
\prov{proveniência: codigo \textperiodcentered{} fato \textperiodcentered{} confiança 1.0 \textperiodcentered{} domínio (sem) \textperiodcentered{} história 204 versões no jornal}

%CRISTAL {"arestas": [{"alvo": "colapso", "confianca": 0.016, "epistemico": "fato", "origem": "manual", "rel": "referencia"}, {"alvo": "kappa", "confianca": 0.95, "epistemico": "fato", "origem": "manual", "rel": "especializa"}], "confianca": 1.0, "contraexemplos": [], "descricao": "[Medida ] Fixe um espaço de estados S (configurações, impressões, snapshots). Uma é um funcional [: S R _ 0, ] monótono sob agregação observável: estados mais estruturados (menos dispersos, mais estáveis sob transformação) carregam mais informativo. não entra no colapso --- observa o da propagação/atenuação, complementar a.", "epistemico": "fato", "exemplos": [], "id": "medida_fixe_um_espaco_de_estados_s_configuracoes_impres", "memoria": "semantica", "meta": {}, "origem": "papers/casl-propagation.tex", "palavras_chave": [], "sinonimos": [], "tipo": "conceito", "titulo": "[Medida ] Fixe um espaço de estados S (configurações, impres…"}
\section{[Medida ] Fixe um espaço de estados S (configurações, impres…}
[Medida ] Fixe um espaço de estados S (configurações, impressões, snapshots). Uma é um funcional [: S R \_ 0, ] monótono sob agregação observável: estados mais estruturados (menos dispersos, mais estáveis sob transformação) carregam mais informativo. não entra no colapso --- observa o da propagação/atenuação, complementar a.
\prov{relações: referencia colapso; especializa kappa}
\prov{proveniência: papers/casl-propagation.tex \textperiodcentered{} fato \textperiodcentered{} confiança 1.0 \textperiodcentered{} domínio (sem) \textperiodcentered{} história 64 versões no jornal}

%CRISTAL {"arestas": [{"alvo": "kappa", "confianca": 0.95, "epistemico": "fato", "origem": "manual", "rel": "referencia"}], "confianca": 1.0, "contraexemplos": [], "descricao": "[Métricas por domínio] Cada instância escolhe sua métrica concreta de: @ ll@ CASL & entropia de políticas, cardinalidade de grants ativos & dissipação não-associativa acumulada (recipiente) & frames, atratores únicos, bits de estabilidade por turno", "epistemico": "fato", "exemplos": [], "id": "metricas_por_dominio_cada_instancia_escolhe_sua_metrica_co", "memoria": "semantica", "meta": {}, "origem": "papers/casl-propagation.tex", "palavras_chave": [], "sinonimos": [], "tipo": "conceito", "titulo": "[Métricas por domínio] Cada instância escolhe sua métrica co…"}
\section{[Métricas por domínio] Cada instância escolhe sua métrica co…}
[Métricas por domínio] Cada instância escolhe sua métrica concreta de: @ ll@ CASL \& entropia de políticas, cardinalidade de grants ativos \& dissipação não-associativa acumulada (recipiente) \& frames, atratores únicos, bits de estabilidade por turno
\prov{relações: referencia kappa}
\prov{proveniência: papers/casl-propagation.tex \textperiodcentered{} fato \textperiodcentered{} confiança 1.0 \textperiodcentered{} domínio (sem) \textperiodcentered{} história 61 versões no jornal}

%CRISTAL {"arestas": [{"alvo": "utilitarismo", "confianca": 0.5, "epistemico": "fato", "origem": "wiki", "rel": "relaciona"}, {"alvo": "recursao", "confianca": 0.5, "epistemico": "fato", "origem": "wiki", "rel": "relaciona"}, {"alvo": "resultado_anterior_identidade_nao_cobertura", "confianca": 0.5, "epistemico": "fato", "origem": "wiki", "rel": "relaciona"}, {"alvo": "topologia", "confianca": 0.5, "epistemico": "fato", "origem": "wiki", "rel": "relaciona"}], "confianca": 1.0, "contraexemplos": [], "descricao": "Operador que preserva ordem.", "epistemico": "fato", "exemplos": ["join preserva autorização", "revogação só diminui capacidade"], "id": "monotono", "memoria": "semantica", "meta": {"serve": "Propriedade central da autorização BAI", "verificacao": "slug idêntico — enriquecer, não duplicar"}, "origem": "manual", "palavras_chave": [], "sinonimos": [], "tipo": "conceito", "titulo": "monotono"}
\section{monotono}
Operador que preserva ordem.
\prov{exemplos: join preserva autorização; revogação só diminui capacidade}
\prov{relações: relaciona utilitarismo; relaciona recursao; relaciona resultado\_anterior\_identidade\_nao\_cobertura; relaciona topologia}
\prov{proveniência: manual \textperiodcentered{} fato \textperiodcentered{} confiança 1.0 \textperiodcentered{} domínio (sem) \textperiodcentered{} história 61 versões no jornal}

%CRISTAL {"arestas": [{"alvo": "delta", "confianca": 0.95, "epistemico": "fato", "origem": "manual", "rel": "prova"}], "confianca": 1.0, "contraexemplos": [], "descricao": "[Não-escalada por atenuação] Para todo p e toda Hold (p), existe uma autoridade intrínseca ₀ e uma cadeia de arestas válidas de sua origem até p tais que c_ c_ ₀. Logo [ Ep(s,a) _ ₀ (raiz que alcança p) c_ ₀. ] Nenhum principal acessa uma linha que nenhuma autoridade intrínseca em sua cadeia podia conceder. Indução no comprimento da cadeia de delegação. Caso base: capacidade intrínseca. Passo: a Definição~ exige c_ c_ ₀ a cada aresta; a inclusão é transitiva. A cota de Ep segue p", "epistemico": "fato", "exemplos": [], "id": "nao-escalada_por_atenuacao_para_todo_p_e_toda_hold_p_e", "memoria": "semantica", "meta": {}, "origem": "papers/casl-propagation.tex", "palavras_chave": [], "sinonimos": [], "tipo": "conceito", "titulo": "[Não-escalada por atenuação] Para todo p e toda Hold (p) , e…"}
\section{[Não-escalada por atenuação] Para todo p e toda Hold (p) , e…}
[Não-escalada por atenuação] Para todo p e toda Hold (p), existe uma autoridade intrínseca \ensuremath{_{0}} e uma cadeia de arestas válidas de sua origem até p tais que c\_ c\_ \ensuremath{_{0}}. Logo [ Ep(s,a) \_ \ensuremath{_{0}} (raiz que alcança p) c\_ \ensuremath{_{0}}. ] Nenhum principal acessa uma linha que nenhuma autoridade intrínseca em sua cadeia podia conceder. Indução no comprimento da cadeia de delegação. Caso base: capacidade intrínseca. Passo: a Definição\~{} exige c\_ c\_ \ensuremath{_{0}} a cada aresta; a inclusão é transitiva. A cota de Ep segue p
\prov{relações: prova delta}
\prov{proveniência: papers/casl-propagation.tex \textperiodcentered{} fato \textperiodcentered{} confiança 1.0 \textperiodcentered{} domínio (sem) \textperiodcentered{} história 61 versões no jornal}

%CRISTAL {"arestas": [{"alvo": "ponte_quantica_conjunto_dougherty_e_orbitais_de_hidrogenio", "confianca": 1.0, "epistemico": "fato", "origem": "ponte_quantica_dougherty.md", "rel": "parte_de"}, {"alvo": "semantica", "confianca": 1.0, "epistemico": "fato", "origem": "ingest", "rel": "analogia"}], "confianca": 0.95, "contraexemplos": [], "descricao": "Interpretante Peirce (SEP): efeito/entendimento que o signo determina — significação triádica signo↔objecto↔interpretante. Fato filosófico.", "epistemico": "fato", "exemplos": ["Stanford SEP peirce-semiotics", "regressão infinita de interpretantes"], "id": "observador_e_observado", "memoria": "semantica", "meta": {"confirmado_web": "semiotica_peirce_triade", "verificacao": "slug idêntico — enriquecer, não duplicar"}, "origem": "wiki", "palavras_chave": [], "sinonimos": [], "tipo": "conceito", "titulo": "Observador e observado"}
\section{Observador e observado}
Interpretante Peirce (SEP): efeito/entendimento que o signo determina — significação triádica signo\ensuremath{\leftrightarrow}objecto\ensuremath{\leftrightarrow}interpretante. Fato filosófico.
\prov{exemplos: Stanford SEP peirce-semiotics; regressão infinita de interpretantes}
\prov{relações: parte\_de ponte\_quantica\_conjunto\_dougherty\_e\_orbitais\_de\_hidrogenio; analogia semantica}
\prov{proveniência: wiki \textperiodcentered{} fato \textperiodcentered{} confiança 0.95 \textperiodcentered{} domínio (sem) \textperiodcentered{} história 41 versões no jornal}

%CRISTAL {"arestas": [{"alvo": "ponte_quantica_conjunto_dougherty_e_orbitais_de_hidrogenio", "confianca": 1.0, "epistemico": "fato", "origem": "ponte_quantica_dougherty.md", "rel": "parte_de"}, {"alvo": "mecanica_quantica", "confianca": 0.52, "epistemico": "hipotese", "origem": "ingest", "rel": "referencia"}], "confianca": 0.95, "contraexemplos": [], "descricao": "Imagens APS/PRL 2013 (Stodolna): microscopia de fotoionização com lente electrostática ~20 000×; nodos 0–3 observados em ξ=r+z. Fato experimental.", "epistemico": "fato", "exemplos": ["PRL 110, 213001", "Physics APS v6/58"], "id": "orbitais_de_hidrogenio_evidencia_empirica", "memoria": "semantica", "meta": {"doi": "10.1103/PhysRevLett.110.213001", "fonte": "PRL 2013", "verificacao": "slug idêntico — enriquecer, não duplicar"}, "origem": "wiki", "palavras_chave": [], "sinonimos": [], "tipo": "conceito", "titulo": "Orbitais de hidrogénio (evidência empírica)"}
\section{Orbitais de hidrogénio (evidência empírica)}
Imagens APS/PRL 2013 (Stodolna): microscopia de fotoionização com lente electrostática \~{}20 000$\times$; nodos 0–3 observados em \ensuremath{\xi}=r+z. Fato experimental.
\prov{exemplos: PRL 110, 213001; Physics APS v6/58}
\prov{relações: parte\_de ponte\_quantica\_conjunto\_dougherty\_e\_orbitais\_de\_hidrogenio; referencia mecanica\_quantica}
\prov{proveniência: wiki \textperiodcentered{} PRL 2013 \textperiodcentered{} fato \textperiodcentered{} confiança 0.95 \textperiodcentered{} domínio (sem) \textperiodcentered{} história 13 versões no jornal}

%CRISTAL {"arestas": [{"alvo": "colapso", "confianca": 0.023, "epistemico": "fato", "origem": "manual", "rel": "referencia"}, {"alvo": "delta", "confianca": 0.95, "epistemico": "fato", "origem": "manual", "rel": "define"}], "confianca": 1.0, "contraexemplos": [], "descricao": "[Orçamento de propagação] representa um : quantas re-delegações (saltos, turnos, réplicas) a capacidade ainda admite. =0 --- usa mas não repassa; =k --- mais k saltos, decrementando; = --- fecho transitivo ( P ). A verificação de ocorre no , nunca na avaliação/colapso.", "epistemico": "fato", "exemplos": [], "id": "orcamento_de_propagacao_representa_um_quantas_re-delegac", "memoria": "semantica", "meta": {}, "origem": "papers/casl-propagation.tex", "palavras_chave": [], "sinonimos": [], "tipo": "conceito", "titulo": "[Orçamento de propagação] representa um : quantas re-delegaç…"}
\section{[Orçamento de propagação] representa um : quantas re-delegaç…}
[Orçamento de propagação] representa um : quantas re-delegações (saltos, turnos, réplicas) a capacidade ainda admite. =0 --- usa mas não repassa; =k --- mais k saltos, decrementando; = --- fecho transitivo ( P ). A verificação de ocorre no , nunca na avaliação/colapso.
\prov{relações: referencia colapso; define delta}
\prov{proveniência: papers/casl-propagation.tex \textperiodcentered{} fato \textperiodcentered{} confiança 1.0 \textperiodcentered{} domínio (sem) \textperiodcentered{} história 63 versões no jornal}

%CRISTAL {"arestas": [{"alvo": "delta", "confianca": 0.95, "epistemico": "fato", "origem": "manual", "rel": "parte_de"}], "confianca": 1.0, "contraexemplos": [], "descricao": "[Par ( , ) ] (alcance, saltos restantes). (organização estrutural do estado). Diálogos ou delegações longas consomem e, em geral, aumentam --- mais estrutura acumulada --- sem violar atenuação.", "epistemico": "fato", "exemplos": [], "id": "par_alcance_saltos_restantes_organizacao_estru", "memoria": "semantica", "meta": {}, "origem": "papers/casl-propagation.tex", "palavras_chave": [], "sinonimos": [], "tipo": "conceito", "titulo": "[Par ( , ) ] (alcance, saltos restantes). (organização estru…"}
\section{[Par ( , ) ] (alcance, saltos restantes). (organização estru…}
[Par ( , ) ] (alcance, saltos restantes). (organização estrutural do estado). Diálogos ou delegações longas consomem e, em geral, aumentam --- mais estrutura acumulada --- sem violar atenuação.
\prov{relações: parte\_de delta}
\prov{proveniência: papers/casl-propagation.tex \textperiodcentered{} fato \textperiodcentered{} confiança 1.0 \textperiodcentered{} domínio (sem) \textperiodcentered{} história 60 versões no jornal}

%CRISTAL {"arestas": [{"alvo": "microscopia_de_fotoionizacao_orbitais_de_hidrogenio_stodolna_2013", "confianca": 1.0, "epistemico": "fato", "origem": "microscopia_fotoionizacao_hidrogenio.md", "rel": "parte_de"}, {"alvo": "correntes_de_birkeland_filamentos_de_plasma", "confianca": 1.0, "epistemico": "fato", "origem": "birkeland_correntes_plasma.md", "rel": "parte_de"}, {"alvo": "semiotica_triadica_de_charles_sanders_peirce", "confianca": 1.0, "epistemico": "fato", "origem": "semiotica_peirce_triade.md", "rel": "parte_de"}, {"alvo": "ginzburg_toro_helice_e_3d-sst", "confianca": 1.0, "epistemico": "fato", "origem": "ginzburg_toro_helice.md", "rel": "parte_de"}, {"alvo": "a_escala_logaritmica_do_ouro", "confianca": 0.52, "epistemico": "hipotese", "origem": "ingest", "rel": "analogia"}, {"alvo": "birkeland_correntes_plasma", "confianca": 0.52, "epistemico": "hipotese", "origem": "ingest", "rel": "analogia"}, {"alvo": "fase_geometrica_de_berry_pancharatnamberry", "confianca": 1.0, "epistemico": "fato", "origem": "fase_geometrica_berry.md", "rel": "parte_de"}, {"alvo": "colapso_koch_operador_ψ_da_tiffany", "confianca": 1.0, "epistemico": "fato", "origem": "colapso_koch_tiffany.md", "rel": "parte_de"}, {"alvo": "tiffany_cabeca_conversacional_broca-so", "confianca": 1.0, "epistemico": "fato", "origem": "tiffany_cabeca_broca.md", "rel": "parte_de"}, {"alvo": "geometria", "confianca": 0.95, "epistemico": "fato", "origem": "wiki", "rel": "referencia"}, {"alvo": "semantica", "confianca": 0.95, "epistemico": "fato", "origem": "wiki", "rel": "confirma"}, {"alvo": "observador_e_observado", "confianca": 0.95, "epistemico": "fato", "origem": "wiki", "rel": "confirma"}, {"alvo": "a_medida_o_calor_de_volta", "confianca": 0.95, "epistemico": "fato", "origem": "wiki", "rel": "confirma"}, {"alvo": "a_dinamica_a_precessao_de_bloch", "confianca": 0.95, "epistemico": "fato", "origem": "wiki", "rel": "confirma"}, {"alvo": "gato", "confianca": 0.95, "epistemico": "fato", "origem": "wiki", "rel": "usa"}, {"alvo": "olho_de_koch", "confianca": 0.95, "epistemico": "fato", "origem": "wiki", "rel": "usa"}, {"alvo": "codigo_conversa_colapso", "confianca": 0.95, "epistemico": "fato", "origem": "wiki", "rel": "define"}, {"alvo": "painel_estelar", "confianca": 0.95, "epistemico": "fato", "origem": "wiki", "rel": "explica"}, {"alvo": "lobo_frontal", "confianca": 0.95, "epistemico": "fato", "origem": "wiki", "rel": "explica"}, {"alvo": "assistente", "confianca": 0.95, "epistemico": "fato", "origem": "wiki", "rel": "explica"}, {"alvo": "tiffany", "confianca": 0.95, "epistemico": "fato", "origem": "wiki", "rel": "explica"}, {"alvo": "colapso", "confianca": 0.95, "epistemico": "fato", "origem": "wiki", "rel": "usa"}, {"alvo": "hopfield", "confianca": 0.95, "epistemico": "fato", "origem": "wiki", "rel": "usa"}, {"alvo": "ipc", "confianca": 0.95, "epistemico": "fato", "origem": "wiki", "rel": "usa"}, {"alvo": "a_decomposicao_de_peirce_celulas_e_canais", "confianca": 0.95, "epistemico": "fato", "origem": "wiki", "rel": "usa"}, {"alvo": "broca-so_camada_de_trabalho_isomorfica", "confianca": 0.95, "epistemico": "fato", "origem": "wiki", "rel": "especializa"}, {"alvo": "goldchain", "confianca": 0.95, "epistemico": "fato", "origem": "wiki", "rel": "usa"}, {"alvo": "computacao", "confianca": 0.95, "epistemico": "fato", "origem": "wiki", "rel": "parte_de"}], "confianca": 1.0, "contraexemplos": [], "descricao": "| Broca OS | Nó | relação | |----------|-----|---------| | núcleo | broca_os | confirma | | álgebra gato | gato | depende | | IPC Peirce | ipc | usa | | decomposição | a_decomposicao_de_peirce_celulas_e_canais | usa | | processos SO | processos | depende | | computação | computacao | parte_de | | cabeça | tiffany | explica | | colapso | colapso | usa | | Hopfield | hopfield | usa | | camada C | broca-so_camada_de_trabalho_isomorfica | especializa |", "epistemico": "fato", "exemplos": [], "id": "ponte_broca-so", "memoria": "semantica", "meta": {"indexar": false, "verificacao": "slug idêntico — enriquecer, não duplicar"}, "origem": "manual", "palavras_chave": [], "sinonimos": [], "tipo": "conceito", "titulo": "Ponte Broca-SO"}
\section{Ponte Broca-SO}
| Broca OS | Nó | relação | |----------|-----|---------| | núcleo | broca\_os | confirma | | álgebra gato | gato | depende | | IPC Peirce | ipc | usa | | decomposição | a\_decomposicao\_de\_peirce\_celulas\_e\_canais | usa | | processos SO | processos | depende | | computação | computacao | parte\_de | | cabeça | tiffany | explica | | colapso | colapso | usa | | Hopfield | hopfield | usa | | camada C | broca-so\_camada\_de\_trabalho\_isomorfica | especializa |
\prov{relações: parte\_de microscopia\_de\_fotoionizacao\_orbitais\_de\_hidrogenio\_stodolna\_2013; parte\_de correntes\_de\_birkeland\_filamentos\_de\_plasma; parte\_de semiotica\_triadica\_de\_charles\_sanders\_peirce; parte\_de ginzburg\_toro\_helice\_e\_3d-sst; analogia a\_escala\_logaritmica\_do\_ouro; analogia birkeland\_correntes\_plasma; parte\_de fase\_geometrica\_de\_berry\_pancharatnamberry; parte\_de colapso\_koch\_operador\_\ensuremath{\psi}\_da\_tiffany; parte\_de tiffany\_cabeca\_conversacional\_broca-so; referencia geometria; confirma semantica; confirma observador\_e\_observado; confirma a\_medida\_o\_calor\_de\_volta; confirma a\_dinamica\_a\_precessao\_de\_bloch; usa gato; usa olho\_de\_koch; define codigo\_conversa\_colapso; explica painel\_estelar; explica lobo\_frontal; explica assistente; explica tiffany; usa colapso; usa hopfield; usa ipc; usa a\_decomposicao\_de\_peirce\_celulas\_e\_canais; especializa broca-so\_camada\_de\_trabalho\_isomorfica; usa goldchain; parte\_de computacao}
\prov{proveniência: manual \textperiodcentered{} fato \textperiodcentered{} confiança 1.0 \textperiodcentered{} domínio (sem) \textperiodcentered{} história 98 versões no jornal}

%CRISTAL {"arestas": [{"alvo": "bai", "confianca": 0.85, "epistemico": "fato", "origem": "manual", "rel": "referencia"}, {"alvo": "bai_biblioteca_de_autorizacao_isomorfica", "confianca": 0.95, "epistemico": "fato", "origem": "manual", "rel": "define"}], "confianca": 1.0, "contraexemplos": [], "descricao": "Por construção (Definição~ ), uma capacidade emitida com profundidade admite cadeias de re-delegação de comprimento no máximo: [ 1 salto 1 0 ( não re-delegável ). ] [nosep] =0: o destinatário a capacidade, mas não a passa a ninguém. É o caso do “dar acesso sem permitir compartilhar”. =k: re-delegável por mais k saltos; cada salto decrementa. =: irrestrito (o comportamento legado, sem limite).", "epistemico": "fato", "exemplos": [], "id": "por_construcao_definicao_uma_capacidade_emitida_com_pro", "memoria": "semantica", "meta": {}, "origem": "papers/casl-propagation.tex", "palavras_chave": [], "sinonimos": [], "tipo": "conceito", "titulo": "Por construção (Definição~ ), uma capacidade emitida com pro…"}
\section{Por construção (Definição\~{} ), uma capacidade emitida com pro…}
Por construção (Definição\~{} ), uma capacidade emitida com profundidade admite cadeias de re-delegação de comprimento no máximo: [ 1 salto 1 0 ( não re-delegável ). ] [nosep] =0: o destinatário a capacidade, mas não a passa a ninguém. É o caso do “dar acesso sem permitir compartilhar”. =k: re-delegável por mais k saltos; cada salto decrementa. =: irrestrito (o comportamento legado, sem limite).
\prov{relações: referencia bai; define bai\_biblioteca\_de\_autorizacao\_isomorfica}
\prov{proveniência: papers/casl-propagation.tex \textperiodcentered{} fato \textperiodcentered{} confiança 1.0 \textperiodcentered{} domínio (sem) \textperiodcentered{} história 64 versões no jornal}

%CRISTAL {"arestas": [{"alvo": "kernel", "confianca": 0.68, "epistemico": "fato", "origem": "manual", "rel": "referencia"}, {"alvo": "vida-morte", "confianca": 0.5, "epistemico": "fato", "origem": "wiki", "rel": "relaciona"}, {"alvo": "pow_metal_pre_fecho", "confianca": 0.5, "epistemico": "fato", "origem": "wiki", "rel": "relaciona"}, {"alvo": "transcricao", "confianca": 0.5, "epistemico": "fato", "origem": "wiki", "rel": "relaciona"}, {"alvo": "resposta_a_pergunta_dificil", "confianca": 0.5, "epistemico": "fato", "origem": "wiki", "rel": "relaciona"}, {"alvo": "teste_ordem", "confianca": 0.5, "epistemico": "fato", "origem": "wiki", "rel": "relaciona"}, {"alvo": "pow_metal_ressonancia", "confianca": 0.5, "epistemico": "fato", "origem": "wiki", "rel": "relaciona"}, {"alvo": "pow_metal_setor", "confianca": 0.5, "epistemico": "fato", "origem": "wiki", "rel": "relaciona"}, {"alvo": "pow_metal_funil", "confianca": 0.5, "epistemico": "fato", "origem": "wiki", "rel": "relaciona"}, {"alvo": "testes", "confianca": 0.5, "epistemico": "fato", "origem": "wiki", "rel": "relaciona"}, {"alvo": "pow_geometria", "confianca": 0.5, "epistemico": "fato", "origem": "wiki", "rel": "relaciona"}, {"alvo": "topologia", "confianca": 0.5, "epistemico": "fato", "origem": "wiki", "rel": "relaciona"}, {"alvo": "tiffany_cabeca_broca", "confianca": 0.5, "epistemico": "fato", "origem": "wiki", "rel": "relaciona"}, {"alvo": "projeto_cristaljson", "confianca": 0.5, "epistemico": "fato", "origem": "wiki", "rel": "relaciona"}, {"alvo": "recursao", "confianca": 0.5, "epistemico": "fato", "origem": "wiki", "rel": "relaciona"}, {"alvo": "thread", "confianca": 0.5, "epistemico": "fato", "origem": "wiki", "rel": "relaciona"}], "confianca": 1.0, "contraexemplos": [], "descricao": "```", "epistemico": "fato", "exemplos": [], "id": "por_sessao_conversadadostelemetriasessaojsonl", "memoria": "semantica", "meta": {}, "origem": "CASL.md", "palavras_chave": [], "sinonimos": [], "tipo": "conceito", "titulo": "por sessão:  conversa/dados/telemetria/<sessao>.jsonl"}
\section{por sessão:  conversa/dados/telemetria/<sessao>.jsonl}
```
\prov{relações: referencia kernel; relaciona vida-morte; relaciona pow\_metal\_pre\_fecho; relaciona transcricao; relaciona resposta\_a\_pergunta\_dificil; relaciona teste\_ordem; relaciona pow\_metal\_ressonancia; relaciona pow\_metal\_setor; relaciona pow\_metal\_funil; relaciona testes; relaciona pow\_geometria; relaciona topologia; relaciona tiffany\_cabeca\_broca; relaciona projeto\_cristaljson; relaciona recursao; relaciona thread}
\prov{proveniência: CASL.md \textperiodcentered{} fato \textperiodcentered{} confiança 1.0 \textperiodcentered{} domínio (sem) \textperiodcentered{} história 19 versões no jornal}

%CRISTAL {"arestas": [{"alvo": "metais", "confianca": 0.95, "epistemico": "fato", "origem": "paper", "rel": "parte_de"}, {"alvo": "word", "confianca": 0.5, "epistemico": "fato", "origem": "wiki", "rel": "relaciona"}, {"alvo": "virtualizacao", "confianca": 0.5, "epistemico": "fato", "origem": "wiki", "rel": "relaciona"}], "confianca": 1.0, "contraexemplos": [], "descricao": "prata _ ourovivo 117 = ourovivo 117 () ourovivo 117 T id ourovivo 117 fail-closed Parte I --- Biblioteca matemática", "epistemico": "fato", "exemplos": [], "id": "prata_ourovivo_117_ourovivo_117_ourovivo_117_t", "memoria": "semantica", "meta": {"indexar": false, "verificacao": "slug idêntico — enriquecer, não duplicar"}, "origem": "manual", "palavras_chave": [], "sinonimos": [], "tipo": "conceito", "titulo": "prata_ourovivo_117_ourovivo_117_ourovivo_117_t"}
\section{prata\_ourovivo\_117\_ourovivo\_117\_ourovivo\_117\_t}
prata \_ ourovivo 117 = ourovivo 117 () ourovivo 117 T id ourovivo 117 fail-closed Parte I --- Biblioteca matemática
\prov{relações: parte\_de metais; relaciona word; relaciona virtualizacao}
\prov{proveniência: manual \textperiodcentered{} fato \textperiodcentered{} confiança 1.0 \textperiodcentered{} domínio (sem) \textperiodcentered{} história 62 versões no jornal}

%CRISTAL {"arestas": [{"alvo": "bai_biblioteca_de_autorizacao_isomorfica", "confianca": 0.078, "epistemico": "fato", "origem": "manual", "rel": "referencia"}, {"alvo": "bai", "confianca": 0.117, "epistemico": "fato", "origem": "manual", "rel": "referencia"}, {"alvo": "kappa", "confianca": 1.0, "epistemico": "fato", "origem": "paper", "rel": "define"}, {"alvo": "delta", "confianca": 1.0, "epistemico": "fato", "origem": "paper", "rel": "define"}, {"alvo": "word", "confianca": 0.5, "epistemico": "fato", "origem": "wiki", "rel": "relaciona"}], "confianca": 1.0, "contraexemplos": [], "descricao": "**Quinteto:** (κ, δ, Φ, T, Ψ) — teorema-mãe (Thm. I).", "epistemico": "fato", "exemplos": [], "id": "primitivas_parte_i", "memoria": "semantica", "meta": {}, "origem": "CASL.md", "palavras_chave": [], "sinonimos": [], "tipo": "conceito", "titulo": "Primitivas (Parte I)"}
\section{Primitivas (Parte I)}
**Quinteto:** (\ensuremath{\kappa}, \ensuremath{\delta}, \ensuremath{\Phi}, T, \ensuremath{\Psi}) — teorema-mãe (Thm. I).
\prov{relações: referencia bai\_biblioteca\_de\_autorizacao\_isomorfica; referencia bai; define kappa; define delta; relaciona word}
\prov{proveniência: CASL.md \textperiodcentered{} fato \textperiodcentered{} confiança 1.0 \textperiodcentered{} domínio (sem) \textperiodcentered{} história 118 versões no jornal}

%CRISTAL {"arestas": [{"alvo": "delta", "confianca": 0.95, "epistemico": "fato", "origem": "manual", "rel": "especializa"}], "confianca": 1.0, "contraexemplos": [], "descricao": "[Propagação limitada] Uma capacidade emitida com profundidade só pode aparecer (atenuada) ao fim de cadeias de re-delegação de comprimento . Em particular, =0 não re-delegável: ela permanece exclusivamente no destinatário direto. Cada aresta válida decrementa em pelo menos 1 e exige _ _0 1 (Definição~ ); uma cadeia de comprimento requer . Para =0 não há aresta de saída possível.", "epistemico": "fato", "exemplos": [], "id": "propagacao_limitada_uma_capacidade_emitida_com_profundidad", "memoria": "semantica", "meta": {}, "origem": "papers/casl-propagation.tex", "palavras_chave": [], "sinonimos": [], "tipo": "conceito", "titulo": "[Propagação limitada] Uma capacidade emitida com profundidad…"}
\section{[Propagação limitada] Uma capacidade emitida com profundidad…}
[Propagação limitada] Uma capacidade emitida com profundidade só pode aparecer (atenuada) ao fim de cadeias de re-delegação de comprimento . Em particular, =0 não re-delegável: ela permanece exclusivamente no destinatário direto. Cada aresta válida decrementa em pelo menos 1 e exige \_ \_0 1 (Definição\~{} ); uma cadeia de comprimento requer . Para =0 não há aresta de saída possível.
\prov{relações: especializa delta}
\prov{proveniência: papers/casl-propagation.tex \textperiodcentered{} fato \textperiodcentered{} confiança 1.0 \textperiodcentered{} domínio (sem) \textperiodcentered{} história 60 versões no jornal}

%CRISTAL {"arestas": [{"alvo": "complexos", "confianca": 1.0, "epistemico": "fato", "origem": "curriculo", "rel": "especializa"}, {"alvo": "word", "confianca": 0.5, "epistemico": "fato", "origem": "wiki", "rel": "relaciona"}, {"alvo": "virtualizacao", "confianca": 0.5, "epistemico": "fato", "origem": "wiki", "rel": "relaciona"}, {"alvo": "vida-morte", "confianca": 0.5, "epistemico": "fato", "origem": "wiki", "rel": "relaciona"}, {"alvo": "vontade_de_poder", "confianca": 0.5, "epistemico": "fato", "origem": "wiki", "rel": "relaciona"}], "confianca": 1.0, "contraexemplos": [], "descricao": "Completude — Cauchy/Cantor (Lopes); codificação Bergman e transformada 𝒢 (engenharia). Dois charts, um corpo ℝ.", "epistemico": "fato", "exemplos": [], "id": "reais", "memoria": "semantica", "meta": {"serve": "Hub da ponte simbólica", "verificacao": "slug idêntico — enriquecer, não duplicar"}, "origem": "manual", "palavras_chave": [], "sinonimos": ["ℝ"], "tipo": "conceito", "titulo": "reais"}
\section{reais}
Completude — Cauchy/Cantor (Lopes); codificação Bergman e transformada \ensuremath{\mathcal{G}} (engenharia). Dois charts, um corpo \ensuremath{\mathbb{R}}.
\prov{sinónimos: \ensuremath{\mathbb{R}}}
\prov{relações: especializa complexos; relaciona word; relaciona virtualizacao; relaciona vida-morte; relaciona vontade\_de\_poder}
\prov{proveniência: manual \textperiodcentered{} fato \textperiodcentered{} confiança 1.0 \textperiodcentered{} domínio (sem) \textperiodcentered{} história 49 versões no jornal}

%CRISTAL {"arestas": [{"alvo": "microscopia_de_fotoionizacao_orbitais_de_hidrogenio_stodolna_2013", "confianca": 1.0, "epistemico": "fato", "origem": "microscopia_fotoionizacao_hidrogenio.md", "rel": "parte_de"}, {"alvo": "fibracao_de_hopf_e_esfera_de_bloch", "confianca": 1.0, "epistemico": "fato", "origem": "fibracao_hopf_bloch.md", "rel": "parte_de"}, {"alvo": "correntes_de_birkeland_filamentos_de_plasma", "confianca": 1.0, "epistemico": "fato", "origem": "birkeland_correntes_plasma.md", "rel": "parte_de"}, {"alvo": "semiotica_triadica_de_charles_sanders_peirce", "confianca": 1.0, "epistemico": "fato", "origem": "semiotica_peirce_triade.md", "rel": "parte_de"}, {"alvo": "ginzburg_toro_helice_e_3d-sst", "confianca": 1.0, "epistemico": "fato", "origem": "ginzburg_toro_helice.md", "rel": "parte_de"}, {"alvo": "fase_geometrica_de_berry_pancharatnamberry", "confianca": 1.0, "epistemico": "fato", "origem": "fase_geometrica_berry.md", "rel": "parte_de"}, {"alvo": "colapso_koch_operador_ψ_da_tiffany", "confianca": 1.0, "epistemico": "fato", "origem": "colapso_koch_tiffany.md", "rel": "parte_de"}, {"alvo": "tiffany_cabeca_conversacional_broca-so", "confianca": 1.0, "epistemico": "fato", "origem": "tiffany_cabeca_broca.md", "rel": "parte_de"}, {"alvo": "geometria_hub_broca-so", "confianca": 1.0, "epistemico": "fato", "origem": "geometria_hub_broca.md", "rel": "parte_de"}], "confianca": 1.0, "contraexemplos": [], "descricao": "- **Paper:** `papers/broca_os.tex`, `papers/computacao_fundamentos_broca.tex` - **Manual:** `ula/BROCA_SO.md`, `ula/BROCA_CRISTAL.md` - **Código:** `neuronio/neuronio.c`, `neuronio/radio.py`, `neuronio/primitivas_broca.py` - **Currículo:** fase 6 programação/SO → Broca OS → Tiffany", "epistemico": "fato", "exemplos": [], "id": "referencia", "memoria": "semantica", "meta": {"indexar": false, "verificacao": "slug idêntico — enriquecer, não duplicar"}, "origem": "manual", "palavras_chave": [], "sinonimos": [], "tipo": "conceito", "titulo": "Referência"}
\section{Referência}
- **Paper:** `papers/broca\_os.tex`, `papers/computacao\_fundamentos\_broca.tex` - **Manual:** `ula/BROCA\_SO.md`, `ula/BROCA\_CRISTAL.md` - **Código:** `neuronio/neuronio.c`, `neuronio/radio.py`, `neuronio/primitivas\_broca.py` - **Currículo:** fase 6 programação/SO \ensuremath{\to} Broca OS \ensuremath{\to} Tiffany
\prov{relações: parte\_de microscopia\_de\_fotoionizacao\_orbitais\_de\_hidrogenio\_stodolna\_2013; parte\_de fibracao\_de\_hopf\_e\_esfera\_de\_bloch; parte\_de correntes\_de\_birkeland\_filamentos\_de\_plasma; parte\_de semiotica\_triadica\_de\_charles\_sanders\_peirce; parte\_de ginzburg\_toro\_helice\_e\_3d-sst; parte\_de fase\_geometrica\_de\_berry\_pancharatnamberry; parte\_de colapso\_koch\_operador\_\ensuremath{\psi}\_da\_tiffany; parte\_de tiffany\_cabeca\_conversacional\_broca-so; parte\_de geometria\_hub\_broca-so}
\prov{proveniência: manual \textperiodcentered{} fato \textperiodcentered{} confiança 1.0 \textperiodcentered{} domínio (sem) \textperiodcentered{} história 51 versões no jornal}

%CRISTAL {"fusao": [{"arestas": [{"alvo": "algebra", "confianca": 1.0, "epistemico": "fato", "origem": "semilla", "rel": "relaciona"}, {"alvo": "monotono", "confianca": 1.0, "epistemico": "fato", "origem": "semilla", "rel": "relaciona"}, {"alvo": "reticulados", "confianca": 1.0, "epistemico": "fato", "origem": "manual", "rel": "generaliza"}, {"alvo": "ordem", "confianca": 1.0, "epistemico": "fato", "origem": "manual", "rel": "depende"}], "confianca": 1.0, "contraexemplos": ["ordem total finita sem meet de elementos incomparáveis"], "descricao": "Estrutura parcialmente ordenada em que join e meet existem para pares comparáveis.", "epistemico": "fato", "exemplos": ["(P(X), ⊆) — reticulado completo", " divisores de 60"], "id": "reticulado", "memoria": "semantica", "meta": {"verificacao": "slug idêntico — enriquecer, não duplicar"}, "origem": "manual", "palavras_chave": ["ordem", "álgebra", "monótono", "join", "meet"], "sinonimos": ["lattice", "retículo"], "tipo": "conceito", "titulo": "reticulado"}, {"arestas": [{"alvo": "topologia", "confianca": 1.0, "epistemico": "fato", "origem": "curriculo", "rel": "depende"}, {"alvo": "bai_biblioteca_de_autorizacao_isomorfica", "confianca": 1.0, "epistemico": "fato", "origem": "codigo", "rel": "generaliza"}, {"alvo": "ordem", "confianca": 1.0, "epistemico": "fato", "origem": "manual", "rel": "depende"}], "confianca": 1.0, "contraexemplos": ["ℕ com divisibilidade (falta meet de 2 e 3)", "grafo cíclico sem join global"], "descricao": "Ordem parcial com join e meet; estrutura de autorização BAI.", "epistemico": "fato", "exemplos": ["subconjuntos de ℤ com divisibilidade", "partições de um conjunto finito"], "id": "reticulados", "memoria": "semantica", "meta": {"serve": "Modelar autorização monótona (BAI); join=∪, meet=∩", "verificacao": "slug idêntico — enriquecer, não duplicar"}, "origem": "manual", "palavras_chave": ["ordem", "join", "meet", "monótono"], "sinonimos": ["lattice", "reticulado"], "tipo": "conceito", "titulo": "reticulados"}], "id": "reticulado", "tipo": "conceito"}
\section{reticulado}
Estrutura parcialmente ordenada em que join e meet existem para pares comparáveis.
\prov{exemplos: (P(X), \ensuremath{\subseteq}) — reticulado completo; divisores de 60}
\prov{contraexemplos: ordem total finita sem meet de elementos incomparáveis}
\prov{sinónimos: lattice; retículo}
\prov{relações: relaciona algebra; relaciona monotono; generaliza reticulados; depende ordem}
\prov{proveniência: manual \textperiodcentered{} fato \textperiodcentered{} confiança 1.0 \textperiodcentered{} domínio (sem) \textperiodcentered{} história 160 versões no jornal}
\prov{a segunda face --- reticulados:}
Ordem parcial com join e meet; estrutura de autorização BAI.
\prov{fusão de curadoria (julgada): absorve reticulados --- as duas faces acima, intactas no registo}

%CRISTAL {"arestas": [{"alvo": "bai", "confianca": -0.023, "epistemico": "fato", "origem": "manual", "rel": "referencia"}, {"alvo": "bai_biblioteca_de_autorizacao_isomorfica", "confianca": 0.95, "epistemico": "fato", "origem": "manual", "rel": "parte_de"}], "confianca": 1.0, "contraexemplos": [], "descricao": "Seja P o conjunto de (operador, donos de, perfis, usuários). Cada principal um conjunto de capacidades Hold (p) (Seção~ ). Para um par fixo (s,a), a autorização efetiva é a dos grants menos as proibições: [ Ep(s,a);=; ( _ Hold (p) _ =0, (s,a) c_ ) ( _ Hold (p) _ =1, (s,a) c_ ), ] onde (s,a) significa que casa o subject e a ação (com / como curingas).", "epistemico": "fato", "exemplos": [], "id": "seja_p_o_conjunto_de_operador_donos_de_perfis_usuarios", "memoria": "semantica", "meta": {}, "origem": "papers/casl-propagation.tex", "palavras_chave": [], "sinonimos": [], "tipo": "conceito", "titulo": "Seja P o conjunto de (operador, donos de , perfis, usuários)…"}
\section{Seja P o conjunto de (operador, donos de , perfis, usuários)…}
Seja P o conjunto de (operador, donos de, perfis, usuários). Cada principal um conjunto de capacidades Hold (p) (Seção\~{} ). Para um par fixo (s,a), a autorização efetiva é a dos grants menos as proibições: [ Ep(s,a);=; ( \_ Hold (p) \_ =0, (s,a) c\_ ) ( \_ Hold (p) \_ =1, (s,a) c\_ ), ] onde (s,a) significa que casa o subject e a ação (com / como curingas).
\prov{relações: referencia bai; parte\_de bai\_biblioteca\_de\_autorizacao\_isomorfica}
\prov{proveniência: papers/casl-propagation.tex \textperiodcentered{} fato \textperiodcentered{} confiança 1.0 \textperiodcentered{} domínio (sem) \textperiodcentered{} história 64 versões no jornal}

%CRISTAL {"arestas": [{"alvo": "logica", "confianca": 1.0, "epistemico": "fato", "origem": "curriculo", "rel": "depende"}, {"alvo": "linguistica", "confianca": 1.0, "epistemico": "fato", "origem": "wiki", "rel": "parte_de"}], "confianca": 1.0, "contraexemplos": [], "descricao": "Significado de signos, termos e proposições; referência e sentido.", "epistemico": "fato", "exemplos": ["Peirce: signo só significa ao ser interpretado (triádico, não diádico)"], "id": "semantica", "memoria": "semantica", "meta": {"fundamento_web": "semiotica_peirce_triade", "verificacao": "slug idêntico — enriquecer, não duplicar"}, "origem": "wiki", "palavras_chave": ["denotação", "conotação", "verdade"], "sinonimos": ["semântica"], "tipo": "conceito", "titulo": "semantica"}
\section{semantica}
Significado de signos, termos e proposições; referência e sentido.
\prov{exemplos: Peirce: signo só significa ao ser interpretado (triádico, não diádico)}
\prov{sinónimos: semântica}
\prov{relações: depende logica; parte\_de linguistica}
\prov{proveniência: wiki \textperiodcentered{} fato \textperiodcentered{} confiança 1.0 \textperiodcentered{} domínio (sem) \textperiodcentered{} história 10 versões no jornal}

%CRISTAL {"arestas": [{"alvo": "bai_biblioteca_de_autorizacao_isomorfica", "confianca": 0.008, "epistemico": "fato", "origem": "manual", "rel": "referencia"}, {"alvo": "bai", "confianca": 0.85, "epistemico": "fato", "origem": "manual", "rel": "referencia"}, {"alvo": "vida-morte", "confianca": 0.5, "epistemico": "fato", "origem": "wiki", "rel": "relaciona"}, {"alvo": "word", "confianca": 0.5, "epistemico": "fato", "origem": "wiki", "rel": "relaciona"}, {"alvo": "virtualizacao", "confianca": 0.5, "epistemico": "fato", "origem": "wiki", "rel": "relaciona"}, {"alvo": "vontade_de_poder", "confianca": 0.5, "epistemico": "fato", "origem": "wiki", "rel": "relaciona"}], "confianca": 1.0, "contraexemplos": [], "descricao": "Sem alterar a resposta — observa `(Φ, δ)`:", "epistemico": "fato", "exemplos": [], "id": "telemetria_shadow_producao", "memoria": "semantica", "meta": {}, "origem": "CASL.md", "palavras_chave": [], "sinonimos": [], "tipo": "conceito", "titulo": "Telemetria shadow (produção)"}
\section{Telemetria shadow (produção)}
Sem alterar a resposta — observa `(\ensuremath{\Phi}, \ensuremath{\delta})`:
\prov{relações: referencia bai\_biblioteca\_de\_autorizacao\_isomorfica; referencia bai; relaciona vida-morte; relaciona word; relaciona virtualizacao; relaciona vontade\_de\_poder}
\prov{proveniência: CASL.md \textperiodcentered{} fato \textperiodcentered{} confiança 1.0 \textperiodcentered{} domínio (sem) \textperiodcentered{} história 12 versões no jornal}

%CRISTAL {"arestas": [{"alvo": "geometria_riemanniana", "confianca": 1.0, "epistemico": "fato", "origem": "curriculo", "rel": "continua"}, {"alvo": "corpo_dual", "confianca": 1.0, "epistemico": "fato", "origem": "curriculo", "rel": "referencia"}, {"alvo": "analise", "confianca": 1.0, "epistemico": "fato", "origem": "paper", "rel": "depende"}, {"alvo": "topologia", "confianca": 1.0, "epistemico": "fato", "origem": "paper", "rel": "referencia"}, {"alvo": "mecanica_quantica", "confianca": 1.0, "epistemico": "fato", "origem": "curriculo", "rel": "continua"}], "confianca": 1.0, "contraexemplos": [], "descricao": "Entropia, medida de Parry, temperatura = densidade do atrator; 1ª/2ª lei na balança.", "epistemico": "fato", "exemplos": ["máquina térmica", "segunda lei"], "id": "termodinamica", "memoria": "semantica", "meta": {"serve": "Energia, entropia, irreversibilidade", "verificacao": "slug idêntico — enriquecer, não duplicar"}, "origem": "manual", "palavras_chave": ["entropia", "calor", "Parry"], "sinonimos": ["termodinâmica"], "tipo": "conceito", "titulo": "termodinamica"}
\section{termodinamica}
Entropia, medida de Parry, temperatura = densidade do atrator; 1ª/2ª lei na balança.
\prov{exemplos: máquina térmica; segunda lei}
\prov{sinónimos: termodinâmica}
\prov{relações: continua geometria\_riemanniana; referencia corpo\_dual; depende analise; referencia topologia; continua mecanica\_quantica}
\prov{proveniência: manual \textperiodcentered{} fato \textperiodcentered{} confiança 1.0 \textperiodcentered{} domínio (sem) \textperiodcentered{} história 192 versões no jornal}

%CRISTAL {"arestas": [{"alvo": "bai", "confianca": 0.148, "epistemico": "fato", "origem": "manual", "rel": "referencia"}, {"alvo": "colapso", "confianca": -0.078, "epistemico": "fato", "origem": "manual", "rel": "referencia"}, {"alvo": "bai_biblioteca_de_autorizacao_isomorfica", "confianca": 0.95, "epistemico": "fato", "origem": "manual", "rel": "prova"}], "confianca": 1.0, "contraexemplos": [], "descricao": "[Tese-mãe da ] Toda instância da é determinada pelo [ (T, ), ] onde é capacidade abstrata (Definição~ ), orçamento de propagação (Definição~ ), medida de organização estrutural (Definição~ ), T operador de atenuação deflacionário e operador de colapso (Definições~ e~ ). CASL, e são distintas destas primitivas --- não motivações do formalismo (Seção~ ).", "epistemico": "fato", "exemplos": [], "id": "tese-mae_da_toda_instancia_da_e_determinada_pelo", "memoria": "semantica", "meta": {}, "origem": "papers/casl-propagation.tex", "palavras_chave": [], "sinonimos": [], "tipo": "conceito", "titulo": "[Tese-mãe da ] Toda instância da é determinada pelo [ (…"}
\section{[Tese-mãe da ] Toda instância da é determinada pelo [ (…}
[Tese-mãe da ] Toda instância da é determinada pelo [ (T, ), ] onde é capacidade abstrata (Definição\~{} ), orçamento de propagação (Definição\~{} ), medida de organização estrutural (Definição\~{} ), T operador de atenuação deflacionário e operador de colapso (Definições\~{} e\~{} ). CASL, e são distintas destas primitivas --- não motivações do formalismo (Seção\~{} ).
\prov{relações: referencia bai; referencia colapso; prova bai\_biblioteca\_de\_autorizacao\_isomorfica}
\prov{proveniência: papers/casl-propagation.tex \textperiodcentered{} fato \textperiodcentered{} confiança 1.0 \textperiodcentered{} domínio (sem) \textperiodcentered{} história 68 versões no jornal}

%CRISTAL {"arestas": [{"alvo": "kernel", "confianca": 0.789, "epistemico": "fato", "origem": "manual", "rel": "referencia"}, {"alvo": "transcricao", "confianca": 0.5, "epistemico": "fato", "origem": "wiki", "rel": "relaciona"}, {"alvo": "vida-morte", "confianca": 0.5, "epistemico": "fato", "origem": "wiki", "rel": "relaciona"}, {"alvo": "testes", "confianca": 0.5, "epistemico": "fato", "origem": "wiki", "rel": "relaciona"}, {"alvo": "tiffany_cabeca_broca", "confianca": 0.5, "epistemico": "fato", "origem": "wiki", "rel": "relaciona"}, {"alvo": "topologia", "confianca": 0.5, "epistemico": "fato", "origem": "wiki", "rel": "relaciona"}, {"alvo": "thread", "confianca": 0.5, "epistemico": "fato", "origem": "wiki", "rel": "relaciona"}, {"alvo": "virtualizacao", "confianca": 0.5, "epistemico": "fato", "origem": "wiki", "rel": "relaciona"}], "confianca": 1.0, "contraexemplos": [], "descricao": "ordem parcial", "epistemico": "fato", "exemplos": [], "id": "teste_ordem", "memoria": "semantica", "meta": {"verificacao": "slug idêntico — enriquecer, não duplicar"}, "origem": "manual", "palavras_chave": [], "sinonimos": [], "tipo": "conceito", "titulo": "teste_ordem"}
\section{teste\_ordem}
ordem parcial
\prov{relações: referencia kernel; relaciona transcricao; relaciona vida-morte; relaciona testes; relaciona tiffany\_cabeca\_broca; relaciona topologia; relaciona thread; relaciona virtualizacao}
\prov{proveniência: manual \textperiodcentered{} fato \textperiodcentered{} confiança 1.0 \textperiodcentered{} domínio (sem) \textperiodcentered{} história 84 versões no jornal}

%CRISTAL {"arestas": [{"alvo": "broca_os", "confianca": 1.0, "epistemico": "fato", "origem": "curriculo", "rel": "parte_de"}, {"alvo": "inteligencia_artificial", "confianca": 1.0, "epistemico": "fato", "origem": "curriculo", "rel": "especializa"}, {"alvo": "painel_estelar", "confianca": 1.0, "epistemico": "fato", "origem": "manual", "rel": "explica"}, {"alvo": "lobo_frontal", "confianca": 1.0, "epistemico": "fato", "origem": "paper", "rel": "explica"}, {"alvo": "assistente", "confianca": 1.0, "epistemico": "fato", "origem": "paper", "rel": "explica"}, {"alvo": "codigo_neuronio_experimentos_tiffany", "confianca": 1.0, "epistemico": "fato", "origem": "codigo", "rel": "usa"}, {"alvo": "colapso", "confianca": 1.0, "epistemico": "fato", "origem": "paper", "rel": "usa"}, {"alvo": "hopfield", "confianca": 1.0, "epistemico": "fato", "origem": "codigo", "rel": "especializa"}, {"alvo": "joalheria_hub", "confianca": 1.0, "epistemico": "fato", "origem": "wiki", "rel": "parte_de"}], "confianca": 1.0, "contraexemplos": ["Tiffany ≠ LLM autoregressivo (GPT)", "Tiffany ≠ Hopfield puro — acrescenta colapso Koch + BAI", "Tiffany ≠ Cristal — Cristal administra tecido"], "descricao": "Cabeça conversacional Broca-SO: M=U·P, colapso Ψ sobre tecido, confirmação via experimentos reprodutíveis. Encarna-se no cristal_cli (a identidade do disco); sua continuidade entre sessões vem da ponte de recall (síntese compactada → contexto do colapso) e seu floco (memória não-associativa) roda no metal (libtiffany.so). A régua associativa (overlap/Hopfield) é a outra metade — barata, em Python.", "epistemico": "fato", "exemplos": ["lab.py reproduzir — 62 testes (experimentos_tiffany.py)", "lab.py porque <fala> — explica face vencedora", "colapso.py — camadas dialogos > conhecimento > corpus", "painel_estelar nonce 1572863 — três braços do gato"], "id": "tiffany", "memoria": "semantica", "meta": {"confirmado": "tiffany_cabeca_broca", "fechado_fase7": true, "serve": "Cabeça conversacional — percebe, colapsa, confirma"}, "origem": "wiki", "palavras_chave": ["colapso", "função de onda", "Broca OS"], "sinonimos": ["Tiffany"], "tipo": "conceito", "titulo": "tiffany"}
\section{tiffany}
Cabeça conversacional Broca-SO: M=U\textperiodcentered P, colapso \ensuremath{\Psi} sobre tecido, confirmação via experimentos reprodutíveis. Encarna-se no cristal\_cli (a identidade do disco); sua continuidade entre sessões vem da ponte de recall (síntese compactada \ensuremath{\to} contexto do colapso) e seu floco (memória não-associativa) roda no metal (libtiffany.so). A régua associativa (overlap/Hopfield) é a outra metade — barata, em Python.
\prov{exemplos: lab.py reproduzir — 62 testes (experimentos\_tiffany.py); lab.py porque <fala> — explica face vencedora; colapso.py — camadas dialogos > conhecimento > corpus; painel\_estelar nonce 1572863 — três braços do gato}
\prov{contraexemplos: Tiffany \ensuremath{\ne} LLM autoregressivo (GPT); Tiffany \ensuremath{\ne} Hopfield puro — acrescenta colapso Koch + BAI; Tiffany \ensuremath{\ne} Cristal — Cristal administra tecido}
\prov{sinónimos: Tiffany}
\prov{relações: parte\_de broca\_os; especializa inteligencia\_artificial; explica painel\_estelar; explica lobo\_frontal; explica assistente; usa codigo\_neuronio\_experimentos\_tiffany; usa colapso; especializa hopfield; parte\_de joalheria\_hub}
\prov{proveniência: wiki \textperiodcentered{} fato \textperiodcentered{} confiança 1.0 \textperiodcentered{} domínio (sem) \textperiodcentered{} história 201 versões no jornal}

%CRISTAL {"arestas": [{"alvo": "bai", "confianca": 0.195, "epistemico": "fato", "origem": "manual", "rel": "referencia"}, {"alvo": "casl-propagation", "confianca": 0.95, "epistemico": "fato", "origem": "manual", "rel": "parte_de"}, {"alvo": "bai_biblioteca_de_autorizacao_isomorfica", "confianca": 0.95, "epistemico": "fato", "origem": "manual", "rel": "explica"}], "confianca": 1.0, "contraexemplos": [], "descricao": "tinta!85 Uma biblioteca matemática de propagação e atenuação. CASL, e são instâncias --- não motivações.", "epistemico": "fato", "exemplos": [], "id": "tinta85_uma_biblioteca_matematica_de_propagacao_e_atenuacao", "memoria": "semantica", "meta": {}, "origem": "papers/casl-propagation.tex", "palavras_chave": [], "sinonimos": [], "tipo": "conceito", "titulo": "tinta!85 Uma biblioteca matemática de propagação e atenuação…"}
\section{tinta!85 Uma biblioteca matemática de propagação e atenuação…}
tinta!85 Uma biblioteca matemática de propagação e atenuação. CASL, e são instâncias --- não motivações.
\prov{relações: referencia bai; parte\_de casl-propagation; explica bai\_biblioteca\_de\_autorizacao\_isomorfica}
\prov{proveniência: papers/casl-propagation.tex \textperiodcentered{} fato \textperiodcentered{} confiança 1.0 \textperiodcentered{} domínio (sem) \textperiodcentered{} história 123 versões no jornal}

%CRISTAL {"arestas": [{"alvo": "matrizes", "confianca": 1.0, "epistemico": "fato", "origem": "curriculo", "rel": "usa"}, {"alvo": "espaco_vetorial", "confianca": 1.0, "epistemico": "fato", "origem": "curriculo", "rel": "depende"}], "confianca": 1.0, "contraexemplos": [], "descricao": "Mapa entre espaços vetoriais que preserva soma e produto escalar.", "epistemico": "fato", "exemplos": ["matriz representa mapa linear", "campo espectral PoW"], "id": "transformacao_linear", "memoria": "semantica", "meta": {"verificacao": "slug idêntico — enriquecer, não duplicar"}, "origem": "manual", "palavras_chave": ["kernel", "imagem", "autovalor"], "sinonimos": ["transformação linear"], "tipo": "conceito", "titulo": "transformacao_linear"}
\section{transformacao\_linear}
Mapa entre espaços vetoriais que preserva soma e produto escalar.
\prov{exemplos: matriz representa mapa linear; campo espectral PoW}
\prov{sinónimos: transformação linear}
\prov{relações: usa matrizes; depende espaco\_vetorial}
\prov{proveniência: manual \textperiodcentered{} fato \textperiodcentered{} confiança 1.0 \textperiodcentered{} domínio (sem) \textperiodcentered{} história 164 versões no jornal}

%CRISTAL {"arestas": [{"alvo": "fase_2_tecido_de_conhecimento", "confianca": 1.0, "epistemico": "fato", "origem": "CONHECIMENTO.md", "rel": "parte_de"}, {"alvo": "reticulado", "confianca": 0.078, "epistemico": "fato", "origem": "manual", "rel": "referencia"}], "confianca": 1.0, "contraexemplos": [], "descricao": "| Memória | Tag / tipo | Exemplo | |---------|------------|---------| | Episódica | `fluxo:…` em dialogos | turnos de cadeia | | Semântica | `skn:conceito:…` | reticulado, gato | | Procedural | `skn:procedimento:…` | como colapsar |", "epistemico": "fato", "exemplos": [], "id": "tres_memorias_mesmo_tecido_tags_distintas", "memoria": "semantica", "meta": {}, "origem": "CONHECIMENTO.md", "palavras_chave": [], "sinonimos": [], "tipo": "conceito", "titulo": "Três memórias (mesmo tecido, tags distintas)"}
\section{Três memórias (mesmo tecido, tags distintas)}
| Memória | Tag / tipo | Exemplo | |---------|------------|---------| | Episódica | `fluxo:…` em dialogos | turnos de cadeia | | Semântica | `skn:conceito:…` | reticulado, gato | | Procedural | `skn:procedimento:…` | como colapsar |
\prov{relações: parte\_de fase\_2\_tecido\_de\_conhecimento; referencia reticulado}
\prov{proveniência: CONHECIMENTO.md \textperiodcentered{} fato \textperiodcentered{} confiança 1.0 \textperiodcentered{} domínio (sem) \textperiodcentered{} história 5 versões no jornal}

%CRISTAL {"arestas": [{"alvo": "reticulado", "confianca": 0.039, "epistemico": "fato", "origem": "manual", "rel": "referencia"}, {"alvo": "reticulados", "confianca": 0.95, "epistemico": "fato", "origem": "manual", "rel": "confirma"}], "confianca": 1.0, "contraexemplos": [], "descricao": "[União, não interseção] Compor grants com interseção seria um erro: ``ler onde =A '' e ``ler onde =B '' devem autorizar A B . Só as proibições subtraem.", "epistemico": "fato", "exemplos": [], "id": "uniao_nao_intersecao_compor_grants_com_intersecao_seria_u", "memoria": "semantica", "meta": {}, "origem": "papers/casl-propagation.tex", "palavras_chave": [], "sinonimos": [], "tipo": "conceito", "titulo": "[União, não interseção] Compor grants com interseção seria u…"}
\section{[União, não interseção] Compor grants com interseção seria u…}
[União, não interseção] Compor grants com interseção seria um erro: ``ler onde =A '' e ``ler onde =B '' devem autorizar A B . Só as proibições subtraem.
\prov{relações: referencia reticulado; confirma reticulados}
\prov{proveniência: papers/casl-propagation.tex \textperiodcentered{} fato \textperiodcentered{} confiança 1.0 \textperiodcentered{} domínio (sem) \textperiodcentered{} história 63 versões no jornal}

%CRISTAL {"arestas": [{"alvo": "ponte_quantica_conjunto_dougherty_e_orbitais_de_hidrogenio", "confianca": 1.0, "epistemico": "fato", "origem": "ponte_quantica_dougherty.md", "rel": "parte_de"}, {"alvo": "geometria", "confianca": 0.52, "epistemico": "hipotese", "origem": "ingest", "rel": "referencia"}], "confianca": 0.52, "contraexemplos": [], "descricao": "Vladimir B. Ginzburg (3D-SST): toryces (toro) + helyces (hélice) como elementos primários. Hipótese externa — paralelo Dougherty, não validada.", "epistemico": "hipotese", "exemplos": ["2004 Relativistic Torus and Helix", "Natural Philosophy Wiki"], "id": "unificacao_torohelice_ginzburg", "memoria": "semantica", "meta": {"confirmado_web": "ginzburg_toro_helice", "falsificavel": true, "verificacao": "slug idêntico — enriquecer, não duplicar"}, "origem": "ingest", "palavras_chave": [], "sinonimos": [], "tipo": "conceito", "titulo": "Unificação toro–hélice (Ginzburg)"}
\section{Unificação toro–hélice (Ginzburg)}
Vladimir B. Ginzburg (3D-SST): toryces (toro) + helyces (hélice) como elementos primários. Hipótese externa — paralelo Dougherty, não validada.
\prov{exemplos: 2004 Relativistic Torus and Helix; Natural Philosophy Wiki}
\prov{relações: parte\_de ponte\_quantica\_conjunto\_dougherty\_e\_orbitais\_de\_hidrogenio; referencia geometria}
\prov{proveniência: ingest \textperiodcentered{} hipotese \textperiodcentered{} confiança 0.52 \textperiodcentered{} domínio (sem) \textperiodcentered{} história 10 versões no jornal}

%CRISTAL {"arestas": [{"alvo": "reticulado", "confianca": 0.047, "epistemico": "fato", "origem": "manual", "rel": "referencia"}, {"alvo": "bai", "confianca": -0.109, "epistemico": "fato", "origem": "manual", "rel": "referencia"}, {"alvo": "colapso", "confianca": 0.039, "epistemico": "fato", "origem": "manual", "rel": "referencia"}, {"alvo": "reticulados", "confianca": 0.95, "epistemico": "fato", "origem": "manual", "rel": "referencia"}, {"alvo": "bai_biblioteca_de_autorizacao_isomorfica", "confianca": 0.95, "epistemico": "fato", "origem": "manual", "rel": "prova"}], "confianca": 1.0, "contraexemplos": [], "descricao": "[Universalidade da ] Qualquer sistema cuja autoridade satisfaça: [nosep,=( )] estados e condições formam um reticulado booleano (2X, ); existe operador de atenuação T id sobre extents; existe operador de propagação P id com orçamento decrementado a cada delegação válida; existe colapso = ( _, ) sobre candidatos vigentes, com fail-closed em; admite uma na --- isto é, capacidades, quinteto (T, ) e teoremas de não-escalada, propagação limitada e deny-overrides como corolário", "epistemico": "fato", "exemplos": [], "id": "universalidade_da_qualquer_sistema_cuja_autoridade_satisf", "memoria": "semantica", "meta": {}, "origem": "papers/casl-propagation.tex", "palavras_chave": [], "sinonimos": [], "tipo": "conceito", "titulo": "[Universalidade da ] Qualquer sistema cuja autoridade satisf…"}
\section{[Universalidade da ] Qualquer sistema cuja autoridade satisf…}
[Universalidade da ] Qualquer sistema cuja autoridade satisfaça: [nosep,=( )] estados e condições formam um reticulado booleano (2X, ); existe operador de atenuação T id sobre extents; existe operador de propagação P id com orçamento decrementado a cada delegação válida; existe colapso = ( \_, ) sobre candidatos vigentes, com fail-closed em; admite uma na --- isto é, capacidades, quinteto (T, ) e teoremas de não-escalada, propagação limitada e deny-overrides como corolário
\prov{relações: referencia reticulado; referencia bai; referencia colapso; referencia reticulados; prova bai\_biblioteca\_de\_autorizacao\_isomorfica}
\prov{proveniência: papers/casl-propagation.tex \textperiodcentered{} fato \textperiodcentered{} confiança 1.0 \textperiodcentered{} domínio (sem) \textperiodcentered{} história 129 versões no jornal}

%CRISTAL {"arestas": [{"alvo": "grava_fato_resposta_normal_em_seguida", "confianca": 1.0, "epistemico": "fato", "origem": "CONHECIMENTO.md", "rel": "parte_de"}, {"alvo": "virtualizacao", "confianca": 0.5, "epistemico": "fato", "origem": "wiki", "rel": "relaciona"}, {"alvo": "word", "confianca": 0.5, "epistemico": "fato", "origem": "wiki", "rel": "relaciona"}, {"alvo": "vida-morte", "confianca": 0.5, "epistemico": "fato", "origem": "wiki", "rel": "relaciona"}], "confianca": 1.0, "contraexemplos": [], "descricao": "Ver [`conhecimento/UNIVERSIDADE.md`](conhecimento/UNIVERSIDADE.md).", "epistemico": "fato", "exemplos": [], "id": "universidade_curriculo_em_camadas", "memoria": "semantica", "meta": {}, "origem": "CONHECIMENTO.md", "palavras_chave": [], "sinonimos": [], "tipo": "conceito", "titulo": "Universidade (currículo em camadas)"}
\section{Universidade (currículo em camadas)}
Ver [`conhecimento/UNIVERSIDADE.md`](conhecimento/UNIVERSIDADE.md).
\prov{relações: parte\_de grava\_fato\_resposta\_normal\_em\_seguida; relaciona virtualizacao; relaciona word; relaciona vida-morte}
\prov{proveniência: CONHECIMENTO.md \textperiodcentered{} fato \textperiodcentered{} confiança 1.0 \textperiodcentered{} domínio (sem) \textperiodcentered{} história 5 versões no jornal}

%CRISTAL {"arestas": [{"alvo": "reticulado", "confianca": 0.07, "epistemico": "fato", "origem": "manual", "rel": "referencia"}, {"alvo": "reticulados", "confianca": 0.95, "epistemico": "fato", "origem": "manual", "rel": "especializa"}], "confianca": 1.0, "contraexemplos": [], "descricao": "[Universo e extent] Fixe um universo X. Uma c é um predicado c:X, identificado ao c X. As condições formam o reticulado booleano (2X, X).", "epistemico": "fato", "exemplos": [], "id": "universo_e_extent_fixe_um_universo_x_uma_c_e_um_predicad", "memoria": "semantica", "meta": {}, "origem": "papers/casl-propagation.tex", "palavras_chave": [], "sinonimos": [], "tipo": "conceito", "titulo": "[Universo e extent] Fixe um universo X . Uma c é um predicad…"}
\section{[Universo e extent] Fixe um universo X . Uma c é um predicad…}
[Universo e extent] Fixe um universo X. Uma c é um predicado c:X, identificado ao c X. As condições formam o reticulado booleano (2X, X).
\prov{relações: referencia reticulado; especializa reticulados}
\prov{proveniência: papers/casl-propagation.tex \textperiodcentered{} fato \textperiodcentered{} confiança 1.0 \textperiodcentered{} domínio (sem) \textperiodcentered{} história 64 versões no jornal}

\end{document}
